\nonstopmode{}
\documentclass[a4paper]{book}
\usepackage[times,hyper]{Rd}
\usepackage{makeidx}
\makeatletter\@ifl@t@r\fmtversion{2018/04/01}{}{\usepackage[utf8]{inputenc}}\makeatother
% \usepackage{graphicx} % @USE GRAPHICX@
\makeindex{}
\begin{document}
\chapter*{}
\begin{center}
{\textbf{\huge Package `EventStudy'}}
\par\bigskip{\large \today}
\end{center}
\ifthenelse{\boolean{Rd@use@hyper}}{\hypersetup{pdftitle = {EventStudy: Event Study Analysis}}}{}
\ifthenelse{\boolean{Rd@use@hyper}}{\hypersetup{pdfauthor = {Simon Mueller}}}{}
\begin{description}
\raggedright{}
\item[Type]\AsIs{Package}
\item[Title]\AsIs{Event Study Analysis}
\item[Version]\AsIs{0.62.0}
\item[Date]\AsIs{2026-09-05}
\item[Maintainer]\AsIs{Simon Mueller }\email{sm@data-zoo.de}\AsIs{}
\item[Description]\AsIs{Perform financial event study analysis in R. Implements the
classical event study methodology ('MacKinlay' 1997) with multiple return
models (Market Model, Market Adjusted, Comparison Period Mean Adjusted,
'Fama'-French 3- and 5-factor, 'Carhart' 4-factor, GARCH, rolling window,
DCC-GARCH), parametric and non-parametric test statistics (including
'Kolari'-'Pynnönen' adjusted BMP), buy-and-hold abnormal returns ('BHAR'),
volume and volatility event studies, cross-sectional regression,
'intraday' event studies, panel ('DiD') event studies with modern
estimators ('Callaway' & 'Sant Anna', 'de Chaisemartin' &
'D Haultfoeuille', 'Borusyak' 'Jaravel' & 'Spiess'), synthetic
control methods, wild bootstrap inference, multiple testing corrections,
power analysis via simulation, data download helpers, automated
reporting, and visualization tools. Results can be exported to CSV,
Excel, or 'LaTeX'.}
\item[URL]\AsIs{}\url{https://github.com/sipemu/eventstudy}\AsIs{,
}\url{https://sipemu.github.io/eventstudy/}\AsIs{}
\item[BugReports]\AsIs{}\url{https://github.com/sipemu/eventstudy/issues}\AsIs{}
\item[License]\AsIs{AGPL-3}
\item[Depends]\AsIs{R (>= 4.1.0)}
\item[Imports]\AsIs{R6, distributional, plotly, tibble, ggplot2, dplyr, purrr,
tidyr, stringr, rlang, stats, tools, utils}
\item[Suggests]\AsIs{testthat (>= 3.0.0), withr, gridExtra, knitr, rmarkdown,
openxlsx, rugarch, sandwich, covr, rmgarch, quadprog, callr,
did, DIDmultiplegt, didimputation, tidyquant, quantmod, DT,
zoo, httr2 (>= 1.0.0), jsonlite}
\item[Encoding]\AsIs{UTF-8}
\item[Config/testthat/edition]\AsIs{3}
\item[VignetteBuilder]\AsIs{knitr}
\item[Config/roxygen2/version]\AsIs{8.0.0}
\item[LazyData]\AsIs{true}
\item[NeedsCompilation]\AsIs{no}
\item[Author]\AsIs{Simon Mueller [cre, aut]}
\end{description}
\Rdcontents{Contents}
\HeaderA{EventStudy-package}{EventStudy: Event Study Analysis in R}{EventStudy.Rdash.package}
\aliasA{EventStudy}{EventStudy-package}{EventStudy}
\keyword{internal}{EventStudy-package}
%
\begin{Description}
Perform financial event study analysis in R. Implements the classical
event study methodology (MacKinlay 1997) with multiple return models,
parametric and non-parametric test statistics, and visualization tools.
\end{Description}
%
\begin{Section}{Return Models}

\begin{itemize}

\item{} \code{\LinkA{MarketModel}{MarketModel}} -- OLS market model (single factor, optional HAC SEs)
\item{} \code{\LinkA{MarketAdjustedModel}{MarketAdjustedModel}} -- Market-adjusted returns
\item{} \code{\LinkA{ComparisonPeriodMeanAdjustedModel}{ComparisonPeriodMeanAdjustedModel}} -- Mean-adjusted returns
\item{} \code{\LinkA{FamaFrench3FactorModel}{FamaFrench3FactorModel}} -- Fama-French 3-factor model
\item{} \code{\LinkA{FamaFrench5FactorModel}{FamaFrench5FactorModel}} -- Fama-French 5-factor model
\item{} \code{\LinkA{Carhart4FactorModel}{Carhart4FactorModel}} -- Carhart 4-factor model
\item{} \code{\LinkA{GARCHModel}{GARCHModel}} -- GARCH(1,1) model
\item{} \code{\LinkA{RollingWindowModel}{RollingWindowModel}} -- Rolling-window OLS (time-varying beta)
\item{} \code{\LinkA{DCCGARCHModel}{DCCGARCHModel}} -- DCC-GARCH (time-varying beta)
\item{} \code{\LinkA{BHARModel}{BHARModel}} -- Buy-and-hold abnormal returns
\item{} \code{\LinkA{VolumeModel}{VolumeModel}} -- Volume event study model
\item{} \code{\LinkA{VolatilityModel}{VolatilityModel}} -- Volatility event study model

\end{itemize}

\end{Section}
%
\begin{Section}{Test Statistics}

\strong{Single-event:}
\begin{itemize}

\item{} \code{\LinkA{ARTTest}{ARTTest}} -- Abnormal return t-test
\item{} \code{\LinkA{CARTTest}{CARTTest}} -- Cumulative abnormal return t-test
\item{} \code{\LinkA{BHARTTest}{BHARTTest}} -- BHAR t-test

\end{itemize}

\strong{Multi-event:}
\begin{itemize}

\item{} \code{\LinkA{CSectTTest}{CSectTTest}} -- Cross-sectional t-test
\item{} \code{\LinkA{PatellZTest}{PatellZTest}} -- Patell standardized residual test
\item{} \code{\LinkA{SignTest}{SignTest}} -- Sign test
\item{} \code{\LinkA{GeneralizedSignTest}{GeneralizedSignTest}} -- Generalized sign test (Cowan 1992)
\item{} \code{\LinkA{RankTest}{RankTest}} -- Rank test (Corrado 1989)
\item{} \code{\LinkA{BMPTest}{BMPTest}} -- Boehmer, Musumeci \& Poulsen (1991) test
\item{} \code{\LinkA{KolariPynnonenTest}{KolariPynnonenTest}} -- Kolari-Pynnönen adjusted BMP test
\item{} \code{\LinkA{CalendarTimePortfolioTest}{CalendarTimePortfolioTest}} -- Calendar-time portfolio test

\end{itemize}

\end{Section}
%
\begin{Section}{Inference \& Robustness}

\begin{itemize}

\item{} \code{\LinkA{adjust\_p\_values}{adjust.Rul.p.Rul.values}} -- Multiple testing corrections (BH, Bonferroni, etc.)
\item{} \code{\LinkA{bootstrap\_test}{bootstrap.Rul.test}} -- Wild bootstrap inference
\item{} HAC standard errors via \code{use\_hac = TRUE} in model constructors

\end{itemize}

\end{Section}
%
\begin{Section}{Simulation}

\begin{itemize}

\item{} \code{\LinkA{simulate\_event\_study}{simulate.Rul.event.Rul.study}} -- Monte Carlo power analysis

\end{itemize}

\end{Section}
%
\begin{Section}{Pipeline}

The standard workflow is:
\begin{enumerate}

\item{} Create an \code{\LinkA{EventStudyTask}{EventStudyTask}} with data
\item{} Define a \code{\LinkA{ParameterSet}{ParameterSet}} with model and test choices
\item{} Run \code{\LinkA{run\_event\_study}{run.Rul.event.Rul.study}(task, parameter\_set)} (or
the individual steps: \code{\LinkA{prepare\_event\_study}{prepare.Rul.event.Rul.study}},
\code{\LinkA{fit\_model}{fit.Rul.model}}, \code{\LinkA{calculate\_statistics}{calculate.Rul.statistics}})
\item{} Extract results with \code{get\_ar()}, \code{get\_car()},
\code{get\_aar()}, or \code{\LinkA{tidy.EventStudyTask}{tidy.EventStudyTask}}
\item{} Export with \code{\LinkA{export\_results}{export.Rul.results}}
\item{} Visualize with \code{\LinkA{plot\_event\_study}{plot.Rul.event.Rul.study}}
\item{} Generate reports with \code{\LinkA{generate\_report}{generate.Rul.report}}

\end{enumerate}

\end{Section}
%
\begin{Section}{Extensions}

\begin{itemize}

\item{} \code{\LinkA{cross\_sectional\_regression}{cross.Rul.sectional.Rul.regression}} -- Explain CARs with firm characteristics
\item{} \code{\LinkA{IntradayEventStudyTask}{IntradayEventStudyTask}} -- Intraday event studies
\item{} \code{\LinkA{PanelEventStudyTask}{PanelEventStudyTask}} -- Panel DiD event studies (TWFE,
Sun-Abraham, Callaway-Sant'Anna, de Chaisemartin-D'Haultfoeuille,
Borusyak-Jaravel-Spiess)
\item{} \code{\LinkA{SyntheticControlTask}{SyntheticControlTask}} -- Synthetic control methods
\item{} \code{\LinkA{download\_stock\_data}{download.Rul.stock.Rul.data}}, \code{\LinkA{download\_factor\_data}{download.Rul.factor.Rul.data}} -- Data helpers

\end{itemize}

\end{Section}
%
\begin{Author}
\strong{Maintainer}: Simon Mueller \email{sm@data-zoo.de}

Authors:
\begin{itemize}

\item{} Simon Mueller \email{sm@data-zoo.de}

\end{itemize}


\end{Author}
%
\begin{References}
MacKinlay, A. C. (1997). Event Studies in Economics and Finance.
\emph{Journal of Economic Literature}, 35(1), 13--39.

Fama, E. F. and French, K. R. (1993). Common risk factors in the
returns on stocks and bonds. \emph{Journal of Financial Economics},
33(1), 3--56.

Carhart, M. M. (1997). On persistence in mutual fund performance.
\emph{The Journal of Finance}, 52(1), 57--82.
\end{References}
%
\begin{SeeAlso}
Useful links:
\begin{itemize}

\item{} \url{https://github.com/sipemu/eventstudy}
\item{} \url{https://sipemu.github.io/eventstudy/}
\item{} Report bugs at \url{https://github.com/sipemu/eventstudy/issues}

\end{itemize}


\end{SeeAlso}
\HeaderA{adjust\_p\_values}{Adjust P-Values for Multiple Testing}{adjust.Rul.p.Rul.values}
%
\begin{Description}
Computes adjusted p-values for AAR and CAAR test statistics across the
event window, correcting for the multiple comparisons problem. Supports
all methods available in \code{\LinkA{p.adjust}{p.adjust}}.
\end{Description}
%
\begin{Usage}
\begin{verbatim}
adjust_p_values(task, method = "BH", stat_name = "CSectT", group = NULL)
\end{verbatim}
\end{Usage}
%
\begin{Arguments}
\begin{ldescription}
\item[\code{task}] A fitted EventStudyTask with \code{aar\_caar\_tbl} populated.

\item[\code{method}] Adjustment method passed to \code{\LinkA{p.adjust}{p.adjust}}.
Common choices: \code{"BH"} (Benjamini-Hochberg, default), \code{"bonferroni"},
\code{"holm"}, \code{"hochberg"}, \code{"BY"}, \code{"none"}.

\item[\code{stat\_name}] Name of the multi-event test statistic to adjust.
Must match a column name in \code{task\$aar\_caar\_tbl}. Default \code{"CSectT"}.

\item[\code{group}] Optional group name to filter. If NULL, adjusts all groups.
\end{ldescription}
\end{Arguments}
%
\begin{Value}
A tibble with columns from the original test statistic result plus
\code{p\_raw\_aar}, \code{p\_adj\_aar}, \code{p\_raw\_caar}, \code{p\_adj\_caar}.
\end{Value}
\HeaderA{advisor\_pro}{Advisor Pro — Future Retrieval-Grounded Paid Tier}{advisor.Rul.pro}
\keyword{internal}{advisor\_pro}
%
\begin{Description}
\strong{Advisor Pro} is a planned premium add-on for the EventStudy package
that will provide retrieval-augmented, evidence-grounded AI advice backed by
a curated academic knowledge base.  It is not yet available; this help topic
documents the waitlist mechanism.
\end{Description}
%
\begin{Section}{What Advisor Pro will include}

\begin{itemize}

\item{} Retrieval-grounded recommendations sourced from peer-reviewed
methodology (MacKinlay 1997, Fama 1991, Boehmer et al. 1991, and more)
\item{} Automatic context assembly from \code{\LinkA{es\_diagnostics}{es.Rul.diagnostics}} output
\item{} Structured \code{Advice} objects with traceable evidence chains
\item{} Priority access to new KB rules and model integrations

\end{itemize}

\end{Section}
%
\begin{Section}{Current offline tier}

The current package already ships an offline advice layer based on a
deterministic knowledge base:
\begin{itemize}

\item{} \code{\LinkA{recommend\_stat}{recommend.Rul.stat}()} — KB-based test-statistic recommendation
\item{} \code{\LinkA{flag\_robustness}{flag.Rul.robustness}()} — KB-based robustness flags
\item{} \code{\LinkA{es\_advise}{es.Rul.advise}()} — LLM-backed advice (requires provider)

\end{itemize}

\end{Section}
%
\begin{Section}{Waitlist}

To join the Advisor Pro waitlist, visit:
\url{https://github.com/sipemu/eventstudy\#advisor-pro-waitlist}

To enable an optional footer reminder after printing advice objects:
\begin{alltt}options(eventstudy.advisor_pro_footer = TRUE)\end{alltt}

The footer is silent by default; enabling it appends a static URL — no
network connection is made.
\end{Section}
\HeaderA{AnthropicProvider}{AnthropicProvider}{AnthropicProvider}
%
\begin{Description}
A provider that issues \code{POST \{base\_url\}/v1/messages} against
the native Anthropic Messages API. The request carries the mandatory
\code{max\_tokens} field and the \code{anthropic-version} header; when a
\code{schema} is supplied it adds a tool-use \code{input\_schema} block
(\code{tools} + \code{tool\_choice}) to obtain structured output, and the
response extractor prefers that tool call's structured \code{input} while
falling back to the first text block — so the provider never crashes when the
model replies with plain text instead of a tool call.

The API key is resolved at CALL time from \code{ANTHROPIC\_API\_KEY} (never at
construction, never stored on the object) and attached with
\code{req\_headers\_redacted("x-api-key" = key)}, which redacts the key in
every print/error path. Note \code{req\_auth\_bearer\_token} does NOT redact an
\code{x-api-key} header — the redacted-headers helper is mandatory here. A
missing key, transport/timeout failure, non-2xx status, malformed body, or
missing completion text each degrades to exactly one \code{warning()} plus an
`es\_provider\_response` with \code{text = NA\_character\_} — it never crashes the
session and never returns a fabricated completion.

\code{httr2} is required only for this provider and is guarded by
\code{requireNamespace()} at the top of \code{complete()}, so the package
installs and \code{R CMD check}s cleanly with \code{httr2} absent.
\end{Description}
%
\begin{Section}{Super class}
\code{\LinkA{ProviderBase}{ProviderBase}} -> \code{AnthropicProvider}
\end{Section}
%
\begin{Section}{Public fields}

\begin{description}

\item[\code{max\_tokens}] Integer max completion tokens (Anthropic REQUIRES this).

\end{description}


\end{Section}
%
\begin{Section}{Methods}
%
\begin{SubSection}{Public methods}
\begin{itemize}

\item{} \Rhref{\#method-AnthropicProvider-initialize}{\code{AnthropicProvider\$new()}}
\item{} \Rhref{\#method-AnthropicProvider-complete}{\code{AnthropicProvider\$complete()}}
\item{} \Rhref{\#method-AnthropicProvider-clone}{\code{AnthropicProvider\$clone()}}

\end{itemize}

\end{SubSection}



\hypertarget{method-AnthropicProvider-initialize}{}
%
\begin{SubSection}{\code{AnthropicProvider\$new()}}
Construct an AnthropicProvider. Stores non-secret config only;
the API key is resolved at call time inside \code{complete()}.
%
\begin{SubSubSection}{Usage}

\begin{alltt}AnthropicProvider$new(
  model,
  base_url = "https://api.anthropic.com",
  max_tokens = 1024L,
  ...
)\end{alltt}



\end{SubSubSection}

%
\begin{SubSubSection}{Arguments}

\begin{description}

\item[\code{model}] Character model identifier, e.g. "claude-opus-4-8".
\item[\code{base\_url}] Character base URL. Defaults to the Anthropic cloud
endpoint.
\item[\code{max\_tokens}] Integer maximum number of tokens to generate. The
Anthropic Messages API REQUIRES this field; defaults to 1024.
\item[\code{...}] Ignored; accepted for forward-compatibility.

\end{description}



\end{SubSubSection}

\end{SubSection}




\hypertarget{method-AnthropicProvider-complete}{}
%
\begin{SubSection}{\code{AnthropicProvider\$complete()}}
Complete a prompt via \code{POST \{base\_url\}/v1/messages}.
Resolves the key at call time; on a missing key or any HTTP/parse failure
returns one warning + \code{NA}. Structured output is OPTIONAL: when
\code{schema} is supplied a tool-use \code{input\_schema} block is added and
the extractor prefers the tool call's \code{input}, falling back to the
first text block.
%
\begin{SubSubSection}{Usage}

\begin{alltt}AnthropicProvider$complete(prompt, schema = NULL, ...)\end{alltt}



\end{SubSubSection}

%
\begin{SubSubSection}{Arguments}

\begin{description}

\item[\code{prompt}] Character prompt.
\item[\code{schema}] Optional structured-output schema (list) or \code{NULL}.
\item[\code{...}] Ignored; accepted for forward-compatibility.

\end{description}



\end{SubSubSection}

%
\begin{SubSubSection}{Returns}
An `es\_provider\_response` (see \code{ProviderBase\$complete}).

\end{SubSubSection}

\end{SubSection}




\hypertarget{method-AnthropicProvider-clone}{}
%
\begin{SubSection}{\code{AnthropicProvider\$clone()}}
The objects of this class are cloneable with this method.
%
\begin{SubSubSection}{Usage}

\begin{alltt}AnthropicProvider$clone(deep = FALSE)\end{alltt}



\end{SubSubSection}

%
\begin{SubSubSection}{Arguments}

\begin{description}

\item[\code{deep}] Whether to make a deep clone.

\end{description}



\end{SubSubSection}

\end{SubSection}


\end{Section}
%
\begin{Examples}
\begin{ExampleCode}
## Not run: 
# Reads ANTHROPIC_API_KEY from the environment at call time:
p <- AnthropicProvider$new(model = "claude-opus-4-8")
p$complete("Summarise these event-study diagnostics")

# Structured output via a tool-use input_schema:
schema <- list(type = "object",
               properties = list(advice = list(type = "string")))
p$complete("Recommend a test statistic", schema = schema)

## End(Not run)
\end{ExampleCode}
\end{Examples}
\HeaderA{ARTTest}{Abnormal Return T Statistic (ART)}{ARTTest}
%
\begin{Description}
The AR t-test is a statistical method used to determine whether the abnormal
return of a security on a specific day is significantly different from zero.
This test helps researchers identify whether the event of interest has a
significant impact on the security’s return at a particular point in time.

See also \url{https://eventstudy.de/statistics/ar_car_statistics.html}
\end{Description}
%
\begin{Section}{Super class}
\code{\LinkA{TestStatisticBase}{TestStatisticBase}} -> \code{ARTTest}
\end{Section}
%
\begin{Section}{Public fields}

\begin{description}

\item[\code{name}] Short code of the test statistic.

\end{description}


\end{Section}
%
\begin{Section}{Methods}
%
\begin{SubSection}{Public methods}
\begin{itemize}

\item{} \Rhref{\#method-ARTTest-compute}{\code{ARTTest\$compute()}}
\item{} \Rhref{\#method-ARTTest-clone}{\code{ARTTest\$clone()}}

\end{itemize}

\end{SubSection}




\hypertarget{method-ARTTest-compute}{}
%
\begin{SubSection}{\code{ARTTest\$compute()}}
Computes the test AR test statistics for a single event.
%
\begin{SubSubSection}{Usage}

\begin{alltt}ARTTest$compute(data_tbl, model)\end{alltt}



\end{SubSubSection}

%
\begin{SubSubSection}{Arguments}

\begin{description}

\item[\code{data\_tbl}] The data for a single event with
calculated abnormal returns.
\item[\code{model}] The fitted model that includes the
necessary information for calculating the test
statistic.

\end{description}



\end{SubSubSection}

\end{SubSection}




\hypertarget{method-ARTTest-clone}{}
%
\begin{SubSection}{\code{ARTTest\$clone()}}
The objects of this class are cloneable with this method.
%
\begin{SubSubSection}{Usage}

\begin{alltt}ARTTest$clone(deep = FALSE)\end{alltt}



\end{SubSubSection}

%
\begin{SubSubSection}{Arguments}

\begin{description}

\item[\code{deep}] Whether to make a deep clone.

\end{description}



\end{SubSubSection}

\end{SubSection}


\end{Section}
\HeaderA{BHARModel}{Buy-and-Hold Abnormal Returns (BHAR) Model}{BHARModel}
%
\begin{Description}
Implements Buy-and-Hold Abnormal Returns for long-horizon event studies.
BHAR compounds returns over the event window instead of summing:
\deqn{BHAR_i = \prod(1 + R_{i,t}) - \prod(1 + R_{benchmark,t})}{}

The benchmark is the market/index return by default. This model is
appropriate for long-horizon studies (months/years) where compounding
effects matter.
\end{Description}
%
\begin{Section}{Super class}
\code{\LinkA{ModelBase}{ModelBase}} -> \code{BHARModel}
\end{Section}
%
\begin{Section}{Public fields}

\begin{description}

\item[\code{model\_name}] Name of the model.

\end{description}


\end{Section}
%
\begin{Section}{Methods}
%
\begin{SubSection}{Public methods}
\begin{itemize}

\item{} \Rhref{\#method-BHARModel-fit}{\code{BHARModel\$fit()}}
\item{} \Rhref{\#method-BHARModel-abnormal_returns}{\code{BHARModel\$abnormal\_returns()}}
\item{} \Rhref{\#method-BHARModel-clone}{\code{BHARModel\$clone()}}

\end{itemize}

\end{SubSection}



\hypertarget{method-BHARModel-fit}{}
%
\begin{SubSection}{\code{BHARModel\$fit()}}
Fit the BHAR model. Computes estimation window statistics.
%
\begin{SubSubSection}{Usage}

\begin{alltt}BHARModel$fit(data_tbl)\end{alltt}



\end{SubSubSection}

%
\begin{SubSubSection}{Arguments}

\begin{description}

\item[\code{data\_tbl}] Data frame or tibble.

\end{description}



\end{SubSubSection}

\end{SubSection}




\hypertarget{method-BHARModel-abnormal_returns}{}
%
\begin{SubSection}{\code{BHARModel\$abnormal\_returns()}}
Calculate abnormal returns using buy-and-hold compounding.
%
\begin{SubSubSection}{Usage}

\begin{alltt}BHARModel$abnormal_returns(data_tbl)\end{alltt}



\end{SubSubSection}

%
\begin{SubSubSection}{Arguments}

\begin{description}

\item[\code{data\_tbl}] Data frame or tibble.

\end{description}



\end{SubSubSection}

\end{SubSection}




\hypertarget{method-BHARModel-clone}{}
%
\begin{SubSection}{\code{BHARModel\$clone()}}
The objects of this class are cloneable with this method.
%
\begin{SubSubSection}{Usage}

\begin{alltt}BHARModel$clone(deep = FALSE)\end{alltt}



\end{SubSubSection}

%
\begin{SubSubSection}{Arguments}

\begin{description}

\item[\code{deep}] Whether to make a deep clone.

\end{description}



\end{SubSubSection}

\end{SubSection}


\end{Section}
\HeaderA{BHARTTest}{Buy-and-Hold Abnormal Return T Test (BHARTTest)}{BHARTTest}
%
\begin{Description}
Tests whether the BHAR for a single event is significantly different from
zero. The BHAR is the difference between compounded firm returns and
compounded benchmark returns over the event window.
\end{Description}
%
\begin{Section}{Super class}
\code{\LinkA{TestStatisticBase}{TestStatisticBase}} -> \code{BHARTTest}
\end{Section}
%
\begin{Section}{Public fields}

\begin{description}

\item[\code{name}] Short code of the test statistic.

\end{description}


\end{Section}
%
\begin{Section}{Methods}
%
\begin{SubSection}{Public methods}
\begin{itemize}

\item{} \Rhref{\#method-BHARTTest-compute}{\code{BHARTTest\$compute()}}
\item{} \Rhref{\#method-BHARTTest-clone}{\code{BHARTTest\$clone()}}

\end{itemize}

\end{SubSection}




\hypertarget{method-BHARTTest-compute}{}
%
\begin{SubSection}{\code{BHARTTest\$compute()}}
Computes the BHAR t test for a single event.
%
\begin{SubSubSection}{Usage}

\begin{alltt}BHARTTest$compute(data_tbl, model)\end{alltt}



\end{SubSubSection}

%
\begin{SubSubSection}{Arguments}

\begin{description}

\item[\code{data\_tbl}] The data for a single event with
calculated abnormal returns.
\item[\code{model}] The fitted model.

\end{description}



\end{SubSubSection}

\end{SubSection}




\hypertarget{method-BHARTTest-clone}{}
%
\begin{SubSection}{\code{BHARTTest\$clone()}}
The objects of this class are cloneable with this method.
%
\begin{SubSubSection}{Usage}

\begin{alltt}BHARTTest$clone(deep = FALSE)\end{alltt}



\end{SubSubSection}

%
\begin{SubSubSection}{Arguments}

\begin{description}

\item[\code{deep}] Whether to make a deep clone.

\end{description}



\end{SubSubSection}

\end{SubSection}


\end{Section}
\HeaderA{BMPTest}{BMP Test (Boehmer, Musumeci, Poulsen 1991)}{BMPTest}
%
\begin{Description}
Standardized cross-sectional test that is robust to event-induced variance
increases. The BMP test standardizes abnormal returns by their forecast
error corrected standard deviation, then applies a cross-sectional
t-test to these standardized residuals.
\end{Description}
%
\begin{Section}{Super class}
\code{\LinkA{TestStatisticBase}{TestStatisticBase}} -> \code{BMPTest}
\end{Section}
%
\begin{Section}{Public fields}

\begin{description}

\item[\code{name}] Short code of the test statistic.

\end{description}


\end{Section}
%
\begin{Section}{Methods}
%
\begin{SubSection}{Public methods}
\begin{itemize}

\item{} \Rhref{\#method-BMPTest-compute}{\code{BMPTest\$compute()}}
\item{} \Rhref{\#method-BMPTest-clone}{\code{BMPTest\$clone()}}

\end{itemize}

\end{SubSection}




\hypertarget{method-BMPTest-compute}{}
%
\begin{SubSection}{\code{BMPTest\$compute()}}
Computes the BMP standardized cross-sectional test.
%
\begin{SubSubSection}{Usage}

\begin{alltt}BMPTest$compute(data_tbl, model)\end{alltt}



\end{SubSubSection}

%
\begin{SubSubSection}{Arguments}

\begin{description}

\item[\code{data\_tbl}] The data for a multiple event with
calculated abnormal returns.
\item[\code{model}] The fitted model containing sigma estimates.

\end{description}



\end{SubSubSection}

\end{SubSection}




\hypertarget{method-BMPTest-clone}{}
%
\begin{SubSection}{\code{BMPTest\$clone()}}
The objects of this class are cloneable with this method.
%
\begin{SubSubSection}{Usage}

\begin{alltt}BMPTest$clone(deep = FALSE)\end{alltt}



\end{SubSubSection}

%
\begin{SubSubSection}{Arguments}

\begin{description}

\item[\code{deep}] Whether to make a deep clone.

\end{description}



\end{SubSubSection}

\end{SubSection}


\end{Section}
\HeaderA{bootstrap\_test}{Wild Bootstrap Inference for Event Studies}{bootstrap.Rul.test}
%
\begin{Description}
Computes bootstrap p-values for AAR and CAAR test statistics using
the wild bootstrap approach. Per-firm random weights preserve the
cross-sectional dependence structure while randomizing the sign of
abnormal returns under the null hypothesis.
\end{Description}
%
\begin{Usage}
\begin{verbatim}
bootstrap_test(
  task,
  n_boot = 999L,
  weight_type = "rademacher",
  statistic = "both",
  group = NULL,
  seed = NULL
)
\end{verbatim}
\end{Usage}
%
\begin{Arguments}
\begin{ldescription}
\item[\code{task}] A fitted EventStudyTask with abnormal returns computed.

\item[\code{n\_boot}] Number of bootstrap replications. Default 999.

\item[\code{weight\_type}] Type of bootstrap weights: \code{"rademacher"} (default,
+1/-1 with equal probability) or \code{"mammen"} (Mammen two-point
distribution).

\item[\code{statistic}] Which statistic to bootstrap: \code{"aar"}, \code{"caar"},
or \code{"both"} (default).

\item[\code{group}] Optional group name to filter.

\item[\code{seed}] Optional seed for reproducibility.
\end{ldescription}
\end{Arguments}
%
\begin{Value}
A tibble with columns: \code{relative\_index}, \code{observed\_aar},
\code{observed\_caar}, \code{boot\_p\_aar}, \code{boot\_p\_caar}.
\end{Value}
\HeaderA{calculate\_statistics}{Calculate test statistics}{calculate.Rul.statistics}
%
\begin{Description}
Calculate test statistics
\end{Description}
%
\begin{Usage}
\begin{verbatim}
calculate_statistics(task, parameter_set)
\end{verbatim}
\end{Usage}
%
\begin{Arguments}
\begin{ldescription}
\item[\code{task}] The event study task

\item[\code{parameter\_set}] The parameter set that defines the event study.
\end{ldescription}
\end{Arguments}
%
\begin{Value}
task
\end{Value}
\HeaderA{CalendarTimePortfolioTest}{Calendar-Time Portfolio Test}{CalendarTimePortfolioTest}
%
\begin{Description}
Aggregates event-firm returns into calendar-time portfolios and tests
whether the portfolio intercept (alpha) is significantly different from
zero. This approach naturally handles cross-sectional dependence that
arises when events cluster in calendar time.

For each relative event day, the test forms an equal-weighted portfolio
of all event firms' abnormal returns and computes a t-statistic of the
mean portfolio return.
\end{Description}
%
\begin{Section}{Super class}
\code{\LinkA{TestStatisticBase}{TestStatisticBase}} -> \code{CalendarTimePortfolioTest}
\end{Section}
%
\begin{Section}{Public fields}

\begin{description}

\item[\code{name}] Short code of the test statistic.

\end{description}


\end{Section}
%
\begin{Section}{Methods}
%
\begin{SubSection}{Public methods}
\begin{itemize}

\item{} \Rhref{\#method-CalendarTimePortfolioTest-compute}{\code{CalendarTimePortfolioTest\$compute()}}
\item{} \Rhref{\#method-CalendarTimePortfolioTest-clone}{\code{CalendarTimePortfolioTest\$clone()}}

\end{itemize}

\end{SubSection}




\hypertarget{method-CalendarTimePortfolioTest-compute}{}
%
\begin{SubSection}{\code{CalendarTimePortfolioTest\$compute()}}
Computes the calendar-time portfolio test.
%
\begin{SubSubSection}{Usage}

\begin{alltt}CalendarTimePortfolioTest$compute(data_tbl, model)\end{alltt}



\end{SubSubSection}

%
\begin{SubSubSection}{Arguments}

\begin{description}

\item[\code{data\_tbl}] The data for multiple events with
calculated abnormal returns.
\item[\code{model}] The fitted models (unused directly).

\end{description}



\end{SubSubSection}

\end{SubSection}




\hypertarget{method-CalendarTimePortfolioTest-clone}{}
%
\begin{SubSection}{\code{CalendarTimePortfolioTest\$clone()}}
The objects of this class are cloneable with this method.
%
\begin{SubSubSection}{Usage}

\begin{alltt}CalendarTimePortfolioTest$clone(deep = FALSE)\end{alltt}



\end{SubSubSection}

%
\begin{SubSubSection}{Arguments}

\begin{description}

\item[\code{deep}] Whether to make a deep clone.

\end{description}



\end{SubSubSection}

\end{SubSection}


\end{Section}
\HeaderA{Carhart4FactorModel}{Carhart Four-Factor Model}{Carhart4FactorModel}
%
\begin{Description}
Implements the Carhart (1997) four-factor model, which extends the
Fama-French three-factor model with a momentum factor:
\deqn{R_i - R_f = \alpha + \beta_m (R_m - R_f) + \beta_s SMB + \beta_h HML + \beta_{mom} MOM + \epsilon}{}

Requires columns: \code{excess\_return}, \code{market\_excess}, \code{smb},
\code{hml}, \code{mom}.
\end{Description}
%
\begin{Section}{Super classes}
\code{\LinkA{ModelBase}{ModelBase}} -> \code{\LinkA{LinearFactorModel}{LinearFactorModel}} -> \code{Carhart4FactorModel}
\end{Section}
%
\begin{Section}{Public fields}

\begin{description}

\item[\code{model\_name}] Name of the model.

\item[\code{formula}] The four-factor regression formula.

\item[\code{required\_columns}] Required data columns.

\end{description}


\end{Section}
%
\begin{Section}{Methods}
%
\begin{SubSection}{Public methods}
\begin{itemize}

\item{} \Rhref{\#method-Carhart4FactorModel-initialize}{\code{Carhart4FactorModel\$new()}}
\item{} \Rhref{\#method-Carhart4FactorModel-abnormal_returns}{\code{Carhart4FactorModel\$abnormal\_returns()}}
\item{} \Rhref{\#method-Carhart4FactorModel-clone}{\code{Carhart4FactorModel\$clone()}}

\end{itemize}

\end{SubSection}




\hypertarget{method-Carhart4FactorModel-initialize}{}
%
\begin{SubSection}{\code{Carhart4FactorModel\$new()}}
Create a new Carhart4FactorModel.
%
\begin{SubSubSection}{Usage}

\begin{alltt}Carhart4FactorModel$new(use_hac = FALSE, hac_lag = NULL)\end{alltt}



\end{SubSubSection}

%
\begin{SubSubSection}{Arguments}

\begin{description}

\item[\code{use\_hac}] Logical. Use HAC (Newey-West) standard errors.
\item[\code{hac\_lag}] Integer or NULL. Lag truncation for Newey-West.

\end{description}



\end{SubSubSection}

\end{SubSection}




\hypertarget{method-Carhart4FactorModel-abnormal_returns}{}
%
\begin{SubSection}{\code{Carhart4FactorModel\$abnormal\_returns()}}
Calculate abnormal returns using the four-factor model.
%
\begin{SubSubSection}{Usage}

\begin{alltt}Carhart4FactorModel$abnormal_returns(data_tbl)\end{alltt}



\end{SubSubSection}

%
\begin{SubSubSection}{Arguments}

\begin{description}

\item[\code{data\_tbl}] Data frame or tibble.

\end{description}



\end{SubSubSection}

\end{SubSection}




\hypertarget{method-Carhart4FactorModel-clone}{}
%
\begin{SubSection}{\code{Carhart4FactorModel\$clone()}}
The objects of this class are cloneable with this method.
%
\begin{SubSubSection}{Usage}

\begin{alltt}Carhart4FactorModel$clone(deep = FALSE)\end{alltt}



\end{SubSubSection}

%
\begin{SubSubSection}{Arguments}

\begin{description}

\item[\code{deep}] Whether to make a deep clone.

\end{description}



\end{SubSubSection}

\end{SubSection}


\end{Section}
\HeaderA{CARTTest}{Cumulative Abnormal Return T Statistic (CART)}{CARTTest}
%
\begin{Description}
The CAR t-test is a statistical method used to determine whether the
cumulative abnormal return of a security over an event window is
significantly different from zero. This test helps researchers identify
whether the event of interest has a significant impact on the security’s
return over the entire event window, considering the cumulative effects of
the event.

See also \url{https://eventstudy.de/statistics/ar_car_statistics.html}
\end{Description}
%
\begin{Section}{Super class}
\code{\LinkA{TestStatisticBase}{TestStatisticBase}} -> \code{CARTTest}
\end{Section}
%
\begin{Section}{Public fields}

\begin{description}

\item[\code{name}] Short code of the test statistic.

\end{description}


\end{Section}
%
\begin{Section}{Methods}
%
\begin{SubSection}{Public methods}
\begin{itemize}

\item{} \Rhref{\#method-CARTTest-compute}{\code{CARTTest\$compute()}}
\item{} \Rhref{\#method-CARTTest-clone}{\code{CARTTest\$clone()}}

\end{itemize}

\end{SubSection}




\hypertarget{method-CARTTest-compute}{}
%
\begin{SubSection}{\code{CARTTest\$compute()}}
Computes the test CAR test statistics for a single event.
%
\begin{SubSubSection}{Usage}

\begin{alltt}CARTTest$compute(data_tbl, model)\end{alltt}



\end{SubSubSection}

%
\begin{SubSubSection}{Arguments}

\begin{description}

\item[\code{data\_tbl}] The data for a single event with
calculated abnormal returns.
\item[\code{model}] The fitted model that includes the
necessary information for calculating the test
statistic.

\end{description}



\end{SubSubSection}

\end{SubSection}




\hypertarget{method-CARTTest-clone}{}
%
\begin{SubSection}{\code{CARTTest\$clone()}}
The objects of this class are cloneable with this method.
%
\begin{SubSubSection}{Usage}

\begin{alltt}CARTTest$clone(deep = FALSE)\end{alltt}



\end{SubSubSection}

%
\begin{SubSubSection}{Arguments}

\begin{description}

\item[\code{deep}] Whether to make a deep clone.

\end{description}



\end{SubSubSection}

\end{SubSection}


\end{Section}
\HeaderA{car\_by\_group}{Compare CARs Across Groups}{car.Rul.by.Rul.group}
%
\begin{Description}
Performs a t-test (two groups) or ANOVA (3+ groups) to test whether
CARs differ significantly across groups.
\end{Description}
%
\begin{Usage}
\begin{verbatim}
car_by_group(task, group_var = "group", car_window = NULL)
\end{verbatim}
\end{Usage}
%
\begin{Arguments}
\begin{ldescription}
\item[\code{task}] A fitted EventStudyTask.

\item[\code{group\_var}] Name of the grouping variable. Defaults to "group"
(the standard group column in EventStudyTask).

\item[\code{car\_window}] Optional two-element vector for CAR window.
\end{ldescription}
\end{Arguments}
%
\begin{Value}
A list with test results and group-level summary statistics.
\end{Value}
\HeaderA{car\_quantiles}{CAR Quantiles}{car.Rul.quantiles}
%
\begin{Description}
Compute quantiles of the CAR distribution across events.
\end{Description}
%
\begin{Usage}
\begin{verbatim}
car_quantiles(task, probs = c(0.05, 0.25, 0.5, 0.75, 0.95), car_window = NULL)
\end{verbatim}
\end{Usage}
%
\begin{Arguments}
\begin{ldescription}
\item[\code{task}] A fitted EventStudyTask.

\item[\code{probs}] Probability vector for quantiles. Default
\code{c(0.05, 0.25, 0.5, 0.75, 0.95)}.

\item[\code{car\_window}] Optional two-element vector for CAR window.
\end{ldescription}
\end{Arguments}
%
\begin{Value}
A named numeric vector of quantiles.
\end{Value}
\HeaderA{ComparisonPeriodMeanAdjustedModel}{Comparison Period Mean Adjusted Model}{ComparisonPeriodMeanAdjustedModel}
%
\begin{Description}
The Comparison Period Mean Adjusted Model is another relatively simple
approach used in event studies to estimate the expected returns of a stock
and calculate its abnormal returns during an event window. This model is
based on the assumption that a stock’s expected return during the event
window is equal to its average return during a comparison period (typically
a pre-event period). This model is particularly useful when researchers want
to control for a stock’s historical performance and do not wish to rely on
market return data.
\end{Description}
%
\begin{Section}{Super class}
\code{\LinkA{ModelBase}{ModelBase}} -> \code{ComparisonPeriodMeanAdjustedModel}
\end{Section}
%
\begin{Section}{Public fields}

\begin{description}

\item[\code{model\_name}] Name of the model.

\end{description}


\end{Section}
%
\begin{Section}{Methods}
%
\begin{SubSection}{Public methods}
\begin{itemize}

\item{} \Rhref{\#method-ComparisonPeriodMeanAdjustedModel-fit}{\code{ComparisonPeriodMeanAdjustedModel\$fit()}}
\item{} \Rhref{\#method-ComparisonPeriodMeanAdjustedModel-abnormal_returns}{\code{ComparisonPeriodMeanAdjustedModel\$abnormal\_returns()}}
\item{} \Rhref{\#method-ComparisonPeriodMeanAdjustedModel-clone}{\code{ComparisonPeriodMeanAdjustedModel\$clone()}}

\end{itemize}

\end{SubSection}



\hypertarget{method-ComparisonPeriodMeanAdjustedModel-fit}{}
%
\begin{SubSection}{\code{ComparisonPeriodMeanAdjustedModel\$fit()}}
Fit the model with given data.
%
\begin{SubSubSection}{Usage}

\begin{alltt}ComparisonPeriodMeanAdjustedModel$fit(data_tbl)\end{alltt}



\end{SubSubSection}

%
\begin{SubSubSection}{Arguments}

\begin{description}

\item[\code{data\_tbl}] Data frame or tibble containing the data to fit.

\end{description}



\end{SubSubSection}

\end{SubSection}




\hypertarget{method-ComparisonPeriodMeanAdjustedModel-abnormal_returns}{}
%
\begin{SubSection}{\code{ComparisonPeriodMeanAdjustedModel\$abnormal\_returns()}}
Calculate the abnormal returns with given data.
%
\begin{SubSubSection}{Usage}

\begin{alltt}ComparisonPeriodMeanAdjustedModel$abnormal_returns(data_tbl)\end{alltt}



\end{SubSubSection}

%
\begin{SubSubSection}{Arguments}

\begin{description}

\item[\code{data\_tbl}] Data frame or tibble containing the data to calculate abnormal returns.

\end{description}



\end{SubSubSection}

\end{SubSection}




\hypertarget{method-ComparisonPeriodMeanAdjustedModel-clone}{}
%
\begin{SubSection}{\code{ComparisonPeriodMeanAdjustedModel\$clone()}}
The objects of this class are cloneable with this method.
%
\begin{SubSubSection}{Usage}

\begin{alltt}ComparisonPeriodMeanAdjustedModel$clone(deep = FALSE)\end{alltt}



\end{SubSubSection}

%
\begin{SubSubSection}{Arguments}

\begin{description}

\item[\code{deep}] Whether to make a deep clone.

\end{description}



\end{SubSubSection}

\end{SubSection}


\end{Section}
\HeaderA{cross\_sectional\_regression}{Cross-Sectional Regression of CARs}{cross.Rul.sectional.Rul.regression}
%
\begin{Description}
Regress cumulative abnormal returns (CARs) on firm characteristics to
explain cross-sectional variation in event effects. Supports OLS with
heteroskedasticity-consistent (HC) standard errors.
\end{Description}
%
\begin{Usage}
\begin{verbatim}
cross_sectional_regression(
  task,
  formula,
  data,
  car_window = NULL,
  robust = TRUE
)
\end{verbatim}
\end{Usage}
%
\begin{Arguments}
\begin{ldescription}
\item[\code{task}] A fitted EventStudyTask with abnormal returns computed.

\item[\code{formula}] A formula with the response on the left (ignored; CAR is
always the dependent variable) and explanatory variables on the right,
e.g., \code{\textasciitilde{} log\_market\_cap + leverage}.

\item[\code{data}] A data frame of firm characteristics. Must contain an
\code{event\_id} column to merge with CARs.

\item[\code{car\_window}] Two-element integer vector specifying the CAR window
as \code{c(start, end)} relative indices. Default is the full event window.

\item[\code{robust}] Logical. If TRUE and the \pkg{sandwich} package is available,
compute HC1 robust standard errors. Default TRUE.
\end{ldescription}
\end{Arguments}
%
\begin{Value}
A list with class \code{"es\_cross\_sectional"} containing:
\begin{description}

\item[model] The fitted \code{lm} object
\item[coefficients] Coefficient table with (robust) standard errors
\item[r\_squared] R-squared of the regression
\item[n\_obs] Number of observations
\item[car\_data] The merged CAR + characteristics data

\end{description}

\end{Value}
\HeaderA{CSectTTest}{Cross-Sectional T Test (CSectTTest)}{CSectTTest}
%
\begin{Description}
The Cross-Sectional Test (CSect T) is a statistical tool employed in event
studies to evaluate the null hypothesis that the average abnormal return at
the event date is zero. The test statistic is distributed as \eqn{t_{N-1}}{}, where
N is the number of events in that group.

See also \url{https://eventstudy.de/statistics/aar_caar_statistics.html}
\end{Description}
%
\begin{Section}{Super class}
\code{\LinkA{TestStatisticBase}{TestStatisticBase}} -> \code{CSectTTest}
\end{Section}
%
\begin{Section}{Public fields}

\begin{description}

\item[\code{name}] Short code of the test statistic.

\end{description}


\end{Section}
%
\begin{Section}{Methods}
%
\begin{SubSection}{Public methods}
\begin{itemize}

\item{} \Rhref{\#method-CSectTTest-compute}{\code{CSectTTest\$compute()}}
\item{} \Rhref{\#method-CSectTTest-clone}{\code{CSectTTest\$clone()}}

\end{itemize}

\end{SubSection}




\hypertarget{method-CSectTTest-compute}{}
%
\begin{SubSection}{\code{CSectTTest\$compute()}}
Computes the AAR/CAAR cross-sectional t test statistics.
%
\begin{SubSubSection}{Usage}

\begin{alltt}CSectTTest$compute(data_tbl, model)\end{alltt}



\end{SubSubSection}

%
\begin{SubSubSection}{Arguments}

\begin{description}

\item[\code{data\_tbl}] The data for a multiple event with
calculated abnormal returns.
\item[\code{model}] The fitted model: not necessary for this
test statistics.

\end{description}



\end{SubSubSection}

\end{SubSection}




\hypertarget{method-CSectTTest-clone}{}
%
\begin{SubSection}{\code{CSectTTest\$clone()}}
The objects of this class are cloneable with this method.
%
\begin{SubSubSection}{Usage}

\begin{alltt}CSectTTest$clone(deep = FALSE)\end{alltt}



\end{SubSubSection}

%
\begin{SubSubSection}{Arguments}

\begin{description}

\item[\code{deep}] Whether to make a deep clone.

\end{description}



\end{SubSubSection}

\end{SubSection}


\end{Section}
\HeaderA{CustomProvider}{CustomProvider}{CustomProvider}
%
\begin{Description}
A provider that wraps a user-supplied
\code{function(prompt, schema, ...) -> list | character}. This is the
offline end-to-end seam and the escape hatch for backends the built-in HTTP
providers do not cover. It uses NO network and needs neither \code{httr2}
nor \code{jsonlite}.

The user function is run inside \code{tryCatch}: if it errors, the provider
degrades to exactly one \code{warning()} plus an `es\_provider\_response` with
\code{text = NA\_character\_} — it never crashes the session.
\end{Description}
%
\begin{Section}{Super class}
\code{\LinkA{ProviderBase}{ProviderBase}} -> \code{CustomProvider}
\end{Section}
%
\begin{Section}{Methods}
%
\begin{SubSection}{Public methods}
\begin{itemize}

\item{} \Rhref{\#method-CustomProvider-initialize}{\code{CustomProvider\$new()}}
\item{} \Rhref{\#method-CustomProvider-complete}{\code{CustomProvider\$complete()}}
\item{} \Rhref{\#method-CustomProvider-clone}{\code{CustomProvider\$clone()}}

\end{itemize}

\end{SubSection}



\hypertarget{method-CustomProvider-initialize}{}
%
\begin{SubSection}{\code{CustomProvider\$new()}}
Construct a CustomProvider.
%
\begin{SubSubSection}{Usage}

\begin{alltt}CustomProvider$new(fn, model = NULL, ...)\end{alltt}



\end{SubSubSection}

%
\begin{SubSubSection}{Arguments}

\begin{description}

\item[\code{fn}] A function of \code{(prompt, schema, ...)} returning a list or a
character completion. Required.
\item[\code{model}] Optional character model identifier (informational only).
\item[\code{...}] Ignored; forwarded for forward-compatibility.

\end{description}



\end{SubSubSection}

\end{SubSection}




\hypertarget{method-CustomProvider-complete}{}
%
\begin{SubSection}{\code{CustomProvider\$complete()}}
Run the wrapped user function and wrap its result in an
`es\_provider\_response`. A throwing function degrades to one warning + NA.
%
\begin{SubSubSection}{Usage}

\begin{alltt}CustomProvider$complete(prompt, schema = NULL, ...)\end{alltt}



\end{SubSubSection}

%
\begin{SubSubSection}{Arguments}

\begin{description}

\item[\code{prompt}] Character prompt passed to the user function.
\item[\code{schema}] Optional structured-output schema passed through to the user
function.
\item[\code{...}] Additional arguments forwarded to the user function.

\end{description}



\end{SubSubSection}

%
\begin{SubSubSection}{Returns}
An `es\_provider\_response` (see \code{ProviderBase\$complete}).

\end{SubSubSection}

\end{SubSection}




\hypertarget{method-CustomProvider-clone}{}
%
\begin{SubSection}{\code{CustomProvider\$clone()}}
The objects of this class are cloneable with this method.
%
\begin{SubSubSection}{Usage}

\begin{alltt}CustomProvider$clone(deep = FALSE)\end{alltt}



\end{SubSubSection}

%
\begin{SubSubSection}{Arguments}

\begin{description}

\item[\code{deep}] Whether to make a deep clone.

\end{description}



\end{SubSubSection}

\end{SubSection}


\end{Section}
%
\begin{Examples}
\begin{ExampleCode}
# In-process, no network: the user function supplies the completion text.
p <- CustomProvider$new(function(prompt, schema) "canned advice")
res <- p$complete("Summarise these diagnostics")
res$text
res$source
\end{ExampleCode}
\end{Examples}
\HeaderA{DCCGARCHModel}{DCC-GARCH Model}{DCCGARCHModel}
%
\begin{Description}
Event study model using Dynamic Conditional Correlation GARCH for
time-varying beta estimation. Requires the \pkg{rmgarch} package.
The bivariate DCC-GARCH model captures both time-varying volatility
and time-varying correlation between firm and market returns, yielding
a time-varying beta: \eqn{\beta_t = Cov(R_{firm}, R_{market})_t / Var(R_{market})_t}{}.
\end{Description}
%
\begin{Section}{Super class}
\code{\LinkA{ModelBase}{ModelBase}} -> \code{DCCGARCHModel}
\end{Section}
%
\begin{Section}{Public fields}

\begin{description}

\item[\code{model\_name}] Name of the model.

\item[\code{garch\_order}] GARCH order for each univariate model. Default c(1,1).

\item[\code{dcc\_order}] DCC order. Default c(1,1).

\end{description}


\end{Section}
%
\begin{Section}{Methods}
%
\begin{SubSection}{Public methods}
\begin{itemize}

\item{} \Rhref{\#method-DCCGARCHModel-initialize}{\code{DCCGARCHModel\$new()}}
\item{} \Rhref{\#method-DCCGARCHModel-fit}{\code{DCCGARCHModel\$fit()}}
\item{} \Rhref{\#method-DCCGARCHModel-abnormal_returns}{\code{DCCGARCHModel\$abnormal\_returns()}}
\item{} \Rhref{\#method-DCCGARCHModel-clone}{\code{DCCGARCHModel\$clone()}}

\end{itemize}

\end{SubSection}



\hypertarget{method-DCCGARCHModel-initialize}{}
%
\begin{SubSection}{\code{DCCGARCHModel\$new()}}
Create a new DCCGARCHModel.
%
\begin{SubSubSection}{Usage}

\begin{alltt}DCCGARCHModel$new(garch_order = c(1, 1), dcc_order = c(1, 1))\end{alltt}



\end{SubSubSection}

%
\begin{SubSubSection}{Arguments}

\begin{description}

\item[\code{garch\_order}] GARCH(p,q) order for univariate models.
\item[\code{dcc\_order}] DCC(a,b) order.

\end{description}



\end{SubSubSection}

\end{SubSection}




\hypertarget{method-DCCGARCHModel-fit}{}
%
\begin{SubSection}{\code{DCCGARCHModel\$fit()}}
Fit the DCC-GARCH model on the estimation window.
%
\begin{SubSubSection}{Usage}

\begin{alltt}DCCGARCHModel$fit(data_tbl)\end{alltt}



\end{SubSubSection}

%
\begin{SubSubSection}{Arguments}

\begin{description}

\item[\code{data\_tbl}] Data frame or tibble with firm\_returns,
index\_returns, estimation\_window, event\_window columns.

\end{description}



\end{SubSubSection}

\end{SubSection}




\hypertarget{method-DCCGARCHModel-abnormal_returns}{}
%
\begin{SubSection}{\code{DCCGARCHModel\$abnormal\_returns()}}
Calculate abnormal returns using the last conditional beta.
%
\begin{SubSubSection}{Usage}

\begin{alltt}DCCGARCHModel$abnormal_returns(data_tbl)\end{alltt}



\end{SubSubSection}

%
\begin{SubSubSection}{Arguments}

\begin{description}

\item[\code{data\_tbl}] Data frame or tibble.

\end{description}



\end{SubSubSection}

\end{SubSection}




\hypertarget{method-DCCGARCHModel-clone}{}
%
\begin{SubSection}{\code{DCCGARCHModel\$clone()}}
The objects of this class are cloneable with this method.
%
\begin{SubSubSection}{Usage}

\begin{alltt}DCCGARCHModel$clone(deep = FALSE)\end{alltt}



\end{SubSubSection}

%
\begin{SubSubSection}{Arguments}

\begin{description}

\item[\code{deep}] Whether to make a deep clone.

\end{description}



\end{SubSubSection}

\end{SubSection}


\end{Section}
\HeaderA{degenerate-input-contract}{Degenerate-Input Contract for EventStudy Models}{degenerate.Rdash.input.Rdash.contract}
%
\begin{Description}
The degenerate-input contract defines how all EventStudy return models
behave when the estimation data is degenerate. Two modes are supported:

\strong{Lenient (default):} The model sets \code{is\_fitted = FALSE},
emits exactly one \code{warning()} per \code{(event\_id, firm\_symbol)}
per fit call, and propagates \code{NA} through all abnormal returns
and downstream statistics. No event is silently dropped; zeros are
never substituted for \code{NA}.

\strong{Strict:} The model raises a descriptive \code{stop()} error
naming the component, \code{event\_id}, and \code{firm\_symbol}, plus
the specific reason for the degeneracy.

\strong{Configuration:}
\begin{itemize}

\item{} Via \code{ParameterSet}: \code{ParameterSet\$new(degenerate\_handling = "strict")}
\item{} Via package option: \code{options(EventStudy.degenerate\_handling = "strict")}

\end{itemize}

The ParameterSet field takes precedence over the package option; if
neither is set, the default \code{"lenient"} mode is used.

\strong{Degenerate conditions covered:}
\begin{itemize}

\item{} Fewer than 2 finite observations in the estimation window
(insufficient observations for OLS).
\item{} Zero or near-zero variance in index returns
(\code{sd < .Machine\$double.eps}), making OLS estimation undefined.
\item{} Single-event group (relevant for multi-event statistics;
applied in Phase 2).
\item{} \code{NA} propagation from upstream pipeline steps.

\end{itemize}


\strong{NA propagation semantics:}
When \code{is\_fitted = FALSE}, \code{model\$abnormal\_returns()} returns
a tibble with \code{abnormal\_returns = NA\_real\_} for all rows. All
downstream test statistics that depend on fitted models then receive
\code{NA} inputs and propagate \code{NA} to their outputs. The event
is retained in the output tibble — it is never silently dropped.
\end{Description}
%
\begin{Details}
Degenerate-Input Contract for EventStudy Models
\end{Details}
%
\begin{SeeAlso}
\code{\LinkA{ParameterSet}{ParameterSet}}, \code{\LinkA{MarketModel}{MarketModel}}
\end{SeeAlso}
\HeaderA{dieselgate}{Volkswagen "Dieselgate" Multi-Automaker Event Study Dataset}{dieselgate}
\keyword{datasets}{dieselgate}
%
\begin{Description}
A small, frozen dataset bundling daily prices for four German automakers and
the DAX benchmark index around the 2015 "dieselgate" emissions scandal, ready
to drive a complete multi-group event study (\code{prepare\_event\_study()} ->
\code{fit\_model()} -> \code{calculate\_statistics()}).
\end{Description}
%
\begin{Usage}
\begin{verbatim}
data(dieselgate)
\end{verbatim}
\end{Usage}
%
\begin{Format}
A named \code{list} with four elements:
\begin{description}

\item[firm] A tibble of daily prices for all four automakers combined,
with columns \code{symbol} (one of \code{"VOW.DE"}, \code{"PAH3.DE"},
\code{"BMW.DE"}, \code{"MBG.DE"}), \code{date} (character,
\code{"\%d.\%m.\%Y"} format), and \code{adjusted} (numeric adjusted
close). Rows for all firms are stacked (1 440 rows total).
\item[index] A tibble of DAX (\code{"\textasciicircum{}GDAXI"}) daily prices with the same
\code{symbol} / \code{date} / \code{adjusted} columns, used as the
benchmark / reference market for all four events.
\item[request] A four-row tibble giving the event-study request
specifications, one row per firm, with the nine columns expected by
\code{\LinkA{EventStudyTask}{EventStudyTask}}: \code{event\_id} (1L to 4L),
\code{firm\_symbol}, \code{index\_symbol}, \code{event\_date}
(\code{"18.09.2015"}), \code{group} (\code{"VW Group"} for
VOW.DE/PAH3.DE, \code{"Other"} for BMW.DE/MBG.DE),
\code{event\_window\_start} (-10), \code{event\_window\_end} (10),
\code{shift\_estimation\_window} (-11), and
\code{estimation\_window\_length} (250). \code{event\_id = 1} is
VOW.DE for backward compatibility.
\item[meta] A list of provenance metadata: \code{firm\_tickers} (character
vector of all four tickers), \code{groups} (named list mapping group
labels to tickers), \code{index\_ticker}, \code{event\_date}, \code{from},
\code{to}, \code{source}, \code{access\_date}, and \code{note}.

\end{description}

\end{Format}
%
\begin{Details}
The event is the U.S. Environmental Protection Agency's Notice of Violation
issued to Volkswagen on \strong{2015-09-18} (a Friday); the share-price crash
lands on the following trading days. The dataset covers two groups: the
"VW Group" (VOW.DE, PAH3.DE — directly implicated firms) and "Other"
(BMW.DE, MBG.DE — peer automakers). The bundled window layout uses a
250-trading-day estimation window ending 11 days before the event and an
event window of \code{[-10, +10]} trading days.

Fitting a market model on the "VW Group" events produces a strongly negative
cumulative average abnormal return (CAAR approximately -39
roughly -17\% and -13\% average abnormal returns on the first two trading days
after the disclosure, while the "Other" peer automakers show near-zero CAAR
(approximately +1\%), illustrating the idiosyncratic nature of the shock.


\strong{Firms:}
\begin{itemize}

\item{} \code{VOW.DE} — Volkswagen AG ordinary shares (Xetra), event\_id = 1
\item{} \code{PAH3.DE} — Porsche Automobil Holding SE (Xetra), event\_id = 2
\item{} \code{BMW.DE} — BMW AG (Xetra), event\_id = 3
\item{} \code{MBG.DE} — Mercedes-Benz Group AG (Xetra), event\_id = 4

\end{itemize}

\strong{Groups:} "VW Group" (VOW.DE, PAH3.DE) vs "Other" (BMW.DE, MBG.DE).
\strong{Benchmark:} DAX performance index (ticker \code{\textasciicircum{}GDAXI}).
\strong{Date range:} 2014-06-01 to 2015-11-01.
\end{Details}
%
\begin{Source}
Yahoo Finance daily adjusted prices, retrieved 2026-09-04 via the
package's own \code{\LinkA{download\_stock\_data}{download.Rul.stock.Rul.data}}. This is a small
illustrative sample bundled for academic / demonstration use only; see
\code{data-raw/dieselgate.R} for the reproducible fetch script.
\end{Source}
%
\begin{Examples}
\begin{ExampleCode}

data(dieselgate)

# Build a multi-group task and run the full pipeline
task <- EventStudyTask$new(dieselgate$firm, dieselgate$index,
                           dieselgate$request)
task <- run_event_study(task, ParameterSet$new())

# Single-firm: VW crash (event_id = 1)
task$get_car(1L)

# Multi-group: CAAR comparison
vw_caar    <- task$aar_caar_tbl[task$aar_caar_tbl$group == "VW Group", ]$CSectT[[1]]
other_caar <- task$aar_caar_tbl[task$aar_caar_tbl$group == "Other", ]$CSectT[[1]]
tail(vw_caar[, c("relative_index", "caar", "caar_t")], 1)
tail(other_caar[, c("relative_index", "caar", "caar_t")], 1)


\end{ExampleCode}
\end{Examples}
\HeaderA{download\_factor\_data}{Download Factor Data}{download.Rul.factor.Rul.data}
%
\begin{Description}
Downloads Fama-French factor data from the Kenneth French Data Library.
No external package dependencies required (uses base R download + CSV parsing).
\end{Description}
%
\begin{Usage}
\begin{verbatim}
download_factor_data(
  model = c("ff3", "ff5", "mom"),
  frequency = c("daily", "monthly"),
  format_for_task = TRUE
)
\end{verbatim}
\end{Usage}
%
\begin{Arguments}
\begin{ldescription}
\item[\code{model}] Factor model: \code{"ff3"} (default), \code{"ff5"}, \code{"mom"}
(momentum), or \code{"ff3\_monthly"}, \code{"ff5\_monthly"}.

\item[\code{frequency}] Data frequency: \code{"daily"} (default) or \code{"monthly"}.

\item[\code{format\_for\_task}] Logical. If TRUE, formats for EventStudyTask factor\_tbl.
\end{ldescription}
\end{Arguments}
%
\begin{Value}
A tibble with factor data.
\end{Value}
\HeaderA{download\_risk\_free\_rate}{Download Risk-Free Rate}{download.Rul.risk.Rul.free.Rul.rate}
%
\begin{Description}
Downloads risk-free rate data from the Kenneth French Data Library
(extracted from the Fama-French factor files).
\end{Description}
%
\begin{Usage}
\begin{verbatim}
download_risk_free_rate(
  frequency = c("daily", "monthly"),
  source = "french",
  format_for_task = TRUE
)
\end{verbatim}
\end{Usage}
%
\begin{Arguments}
\begin{ldescription}
\item[\code{frequency}] Data frequency: \code{"daily"} (default) or \code{"monthly"}.

\item[\code{source}] Data source: \code{"french"} (default).

\item[\code{format\_for\_task}] Logical. If TRUE, formats for EventStudyTask.
\end{ldescription}
\end{Arguments}
%
\begin{Value}
A tibble with \code{date} and \code{risk\_free\_rate} columns.
\end{Value}
\HeaderA{download\_stock\_data}{Download Stock Data}{download.Rul.stock.Rul.data}
%
\begin{Description}
Downloads historical stock price data for specified symbols. Uses
\pkg{tidyquant} if available, otherwise falls back to \pkg{quantmod}.
\end{Description}
%
\begin{Usage}
\begin{verbatim}
download_stock_data(
  symbols,
  from,
  to = Sys.Date(),
  source = "yahoo",
  format_for_task = TRUE
)
\end{verbatim}
\end{Usage}
%
\begin{Arguments}
\begin{ldescription}
\item[\code{symbols}] Character vector of stock ticker symbols.

\item[\code{from}] Start date (Date or character in YYYY-MM-DD format).

\item[\code{to}] End date. Default \code{Sys.Date()}.

\item[\code{source}] Data source: \code{"yahoo"} (default).

\item[\code{format\_for\_task}] Logical. If TRUE, formats output for
\code{EventStudyTask}: columns \code{symbol}, \code{date} (dd.mm.yyyy),
\code{adjusted}.
\end{ldescription}
\end{Arguments}
%
\begin{Value}
A tibble with stock data.
\end{Value}
\HeaderA{earnings\_surprises}{Earnings Surprise Multi-Firm Event Study Dataset}{earnings.Rul.surprises}
\keyword{datasets}{earnings\_surprises}
%
\begin{Description}
A small, frozen dataset bundling daily prices for three U.S. large-cap firms
and the S\&P 500 benchmark around the Q1 2023 earnings surprise quarter,
ready to drive a complete event study pipeline (\code{prepare\_event\_study()}
-> \code{fit\_model()} -> \code{calculate\_statistics()}).
\end{Description}
%
\begin{Usage}
\begin{verbatim}
data(earnings_surprises)
\end{verbatim}
\end{Usage}
%
\begin{Format}
A named \code{list} with four elements:
\begin{description}

\item[firm] A tibble of daily prices for all three firms combined,
with columns \code{symbol} (one of \code{"AAPL"}, \code{"MSFT"},
\code{"GOOGL"}), \code{date} (character,
\code{"\%d.\%m.\%Y"} format), and \code{adjusted} (numeric adjusted
close). Rows for all firms are stacked (813 rows total).
\item[index] A tibble of S\&P 500 (\code{"\textasciicircum{}GSPC"}) daily prices with the
same \code{symbol} / \code{date} / \code{adjusted} columns, used as the
benchmark for all three events.
\item[request] A three-row tibble giving the event-study request
specifications, one row per firm, with the nine columns expected by
\code{\LinkA{EventStudyTask}{EventStudyTask}}: \code{event\_id} (1L to 3L),
\code{firm\_symbol}, \code{index\_symbol}, \code{event\_date}
(\code{"04.05.2023"} for AAPL, \code{"25.04.2023"} for MSFT and GOOGL),
\code{group} (\code{"Earnings Beat"} for all firms),
\code{event\_window\_start} (-5), \code{event\_window\_end} (5),
\code{shift\_estimation\_window} (-6), and
\code{estimation\_window\_length} (200).
\item[meta] A list of provenance metadata: \code{firm\_tickers},
\code{index\_ticker}, \code{event\_dates}, \code{from}, \code{to},
\code{source}, \code{access\_date}, and \code{note}.

\end{description}

\end{Format}
%
\begin{Details}
All three firms beat consensus EPS estimates in Q1 2023: Apple Inc. reported
on \strong{2023-05-04} and beat by approximately 8\%, driving a strong
next-day return; Microsoft Corporation reported on \strong{2023-04-25},
beating cloud (Azure) estimates with a roughly +7\% next-day response;
Alphabet Inc. reported on \strong{2023-04-25}, with advertising revenue
exceeding expectations. The bundled window layout uses a 200-trading-day
estimation window ending 6 days before each event and an event window of
\code{[-5, +5]} trading days. Running a market model on this single-group
panel produces a positive cumulative average abnormal return (CAAR) over the
event window, reflecting the shared earnings-beat signal across all three firms.


\strong{Firms:}
\begin{itemize}

\item{} \code{AAPL} — Apple Inc. (NASDAQ), event\_id = 1
\item{} \code{MSFT} — Microsoft Corporation (NASDAQ), event\_id = 2
\item{} \code{GOOGL} — Alphabet Inc. Class A (NASDAQ), event\_id = 3

\end{itemize}

\strong{Group:} "Earnings Beat" (all firms).
\strong{Benchmark:} S\&P 500 index (ticker \code{\textasciicircum{}GSPC}).
\strong{Date range:} 2022-06-01 to 2023-06-30.
\end{Details}
%
\begin{Source}
Yahoo Finance daily adjusted prices, retrieved 2026-09-05 via the
package's own \code{\LinkA{download\_stock\_data}{download.Rul.stock.Rul.data}}. This is a small
illustrative sample bundled for academic / demonstration use only; see
\code{data-raw/earnings\_surprises.R} for the reproducible fetch script.
\end{Source}
%
\begin{Examples}
\begin{ExampleCode}

data(earnings_surprises)

# Build task and run the full pipeline
task <- EventStudyTask$new(earnings_surprises$firm,
                           earnings_surprises$index,
                           earnings_surprises$request)
task <- run_event_study(task, ParameterSet$new())

# Multi-event: AAR/CAAR across all firms
caar_tbl <- task$aar_caar_tbl$CSectT[[1]]
tail(caar_tbl[, c("relative_index", "caar", "caar_t")], 1)


\end{ExampleCode}
\end{Examples}
\HeaderA{estimate\_panel\_event\_study}{Estimate Panel Event Study}{estimate.Rul.panel.Rul.event.Rul.study}
%
\begin{Description}
Estimate a panel event study using two-way fixed effects (TWFE) or
dynamic TWFE. Returns event-time coefficient estimates suitable for
event study plots.
\end{Description}
%
\begin{Usage}
\begin{verbatim}
estimate_panel_event_study(
  task,
  method = c("static_twfe", "dynamic_twfe", "sun_abraham", "callaway_santanna",
    "dechaisemartin_dhaultfoeuille", "borusyak_jaravel_spiess"),
  leads = 5,
  lags = 5,
  base_period = -1,
  cluster = NULL,
  ...
)
\end{verbatim}
\end{Usage}
%
\begin{Arguments}
\begin{ldescription}
\item[\code{task}] A \code{PanelEventStudyTask}.

\item[\code{method}] Estimation method. One of \code{"static\_twfe"},
\code{"dynamic\_twfe"}, or \code{"sun\_abraham"}.

\item[\code{leads}] Number of pre-treatment periods to include. Default 5.

\item[\code{lags}] Number of post-treatment periods to include. Default 5.

\item[\code{base\_period}] The reference period (relative to treatment) to omit.
Default -1 (the period just before treatment).

\item[\code{cluster}] Name of the clustering variable for standard errors.
Defaults to the unit ID.

\item[\code{...}] Additional arguments passed to the underlying estimator
(used by \code{callaway\_santanna}, \code{dechaisemartin\_dhaultfoeuille},
and \code{borusyak\_jaravel\_spiess} methods).
\end{ldescription}
\end{Arguments}
%
\begin{Value}
The task with \code{results} populated, containing:
\begin{description}

\item[coefficients] Tibble of event-time coefficients with std errors
\item[model] The fitted model object
\item[method] The estimation method used

\end{description}

\end{Value}
\HeaderA{estimate\_synthetic\_control}{Estimate Synthetic Control}{estimate.Rul.synthetic.Rul.control}
%
\begin{Description}
Estimate synthetic control weights by minimizing the pre-treatment
mean squared prediction error. Weights are non-negative and sum to one.
\end{Description}
%
\begin{Usage}
\begin{verbatim}
estimate_synthetic_control(
  task,
  method = c("quadprog", "optim"),
  covariates = NULL
)
\end{verbatim}
\end{Usage}
%
\begin{Arguments}
\begin{ldescription}
\item[\code{task}] A \code{SyntheticControlTask}.

\item[\code{method}] Optimization method: \code{"quadprog"} (default, requires
\pkg{quadprog}) or \code{"optim"} (uses \code{stats::optim} L-BFGS-B).

\item[\code{covariates}] Optional character vector of covariate column names
in both treated and donor data to include in the matching.
\end{ldescription}
\end{Arguments}
%
\begin{Value}
The task with \code{results} populated.
\end{Value}
\HeaderA{es\_advise}{Grounded AI Advice for Event Study Results}{es.Rul.advise}
%
\begin{Description}
Produces a grounded \code{Advice} S3 object by routing through a task-type
dispatch, calling an optional LLM provider, parsing the JSON response, and
running the runtime grounding guard — which drops any recommendation whose
\code{evidence[]} cites a diagnostic key absent from the computed diagnostics
or a value mismatching beyond numeric tolerance.
\end{Description}
%
\begin{Usage}
\begin{verbatim}
es_advise(diagnostics, task_type, provider = NULL, model = NULL, ...)
\end{verbatim}
\end{Usage}
%
\begin{Arguments}
\begin{ldescription}
\item[\code{diagnostics}] An \code{es\_diagnostics} object returned by
\code{\LinkA{es\_diagnostics}{es.Rul.diagnostics}()}.

\item[\code{task\_type}] Character. One of \code{"interpret"}, \code{"recommend\_stat"},
\code{"recommend\_model"}, \code{"flag\_robustness"}, \code{"design\_discussion"},
\code{"report\_writing"}.

\item[\code{provider}] An optional provider R6 object (from \code{\LinkA{provider}{provider}()}).
Required for LLM-only task types. When \code{NULL} and task type is KB-based,
falls back to the Phase 5 offline path.

\item[\code{model}] Optional character model identifier. Reserved for forward
compatibility: the effective model is the one the provider was constructed
with (see \code{\LinkA{provider}{provider}()}), so set the model there. Accepted here
without error so calling code can pass it, but it does not override the
provider's configured model.

\item[\code{...}] Additional arguments (currently ignored; reserved for future use).
\end{ldescription}
\end{Arguments}
%
\begin{Details}
\strong{Grounding guarantee:} Every returned recommendation is provably tied
to a value the package actually computed (\code{es\_diagnostics()}). The guard
is enforced in R, independent of the prompt (ADV-04).

\strong{Task type routing:}
\begin{description}

\item[\code{recommend\_stat}, \code{flag\_robustness}] 
No provider: returns the Phase 5 \code{es\_advice} object (offline KB
path — deterministic, \code{is\_deterministic = TRUE}).
With provider: KB produces grounded evidence[], LLM adds prose;
returns an \code{Advice} object (\code{is\_deterministic = FALSE}).

\item[\code{interpret}, \code{recommend\_model}, \code{design\_discussion},
\code{report\_writing}] 
LLM-required: \code{stop()} when \code{provider = NULL} (ADV-06).
With provider: LLM produces full advice; guard runs.


\end{description}


\strong{Failure discipline:} Any provider failure, malformed JSON, or empty
response degrades to one \code{warning()} + an empty \code{Advice} object —
never a crash, never a fabricated result (mirrors \code{.handle\_degenerate()}).
\end{Details}
%
\begin{Value}
For KB task types without a provider: an \code{es\_advice} S3 object
(Phase 5 offline path, \code{is\_deterministic = TRUE}).
For all other paths: an \code{Advice} S3 object with fields
\code{source}, \code{is\_deterministic}, \code{task\_type},
\code{interpretation}, \code{recommendations}, \code{caveats},
\code{n\_dropped}.
\end{Value}
%
\begin{SeeAlso}
\code{\LinkA{es\_diagnostics}{es.Rul.diagnostics}}, \code{\LinkA{recommend\_stat}{recommend.Rul.stat}},
\code{\LinkA{flag\_robustness}{flag.Rul.robustness}}, \code{\LinkA{provider}{provider}}
\end{SeeAlso}
%
\begin{Examples}
\begin{ExampleCode}
## Not run: 
task    <- run_event_study(my_task, ParameterSet$new())
diag    <- es_diagnostics(task)

# Offline KB path (no LLM):
advice_kb <- es_advise(diag, task_type = "recommend_stat")
print(advice_kb)   # es_advice S3

# LLM-grounded path:
p         <- provider("openai")
advice    <- es_advise(diag, task_type = "recommend_stat", provider = p)
print(advice)      # Advice S3 with grounding guarantee

## End(Not run)

\end{ExampleCode}
\end{Examples}
\HeaderA{es\_diagnostics}{Harvest Diagnostics from a Fitted EventStudyTask}{es.Rul.diagnostics}
%
\begin{Description}
Extracts already-computed statistical signals from a fitted
\code{EventStudyTask} into a flat, JSON-ready S3 list. This is the
always-available grounding foundation for the knowledge base and offline
advice layer — it recomputes nothing; all signals already exist in the
fitted task.
\end{Description}
%
\begin{Usage}
\begin{verbatim}
es_diagnostics(task, max_events = 20L)
\end{verbatim}
\end{Usage}
%
\begin{Arguments}
\begin{ldescription}
\item[\code{task}] A fitted \code{EventStudyTask} (after \code{fit\_model()} and
optionally \code{calculate\_statistics()} have been called).

\item[\code{max\_events}] Integer. Maximum number of events to include in the
per-event sections (\code{estimation\_window}, \code{event\_window},
\code{contract\_state}). Events are ranked by anomaly score (degenerate
events first, then by absolute final CAR). The remainder is summarised
in \code{aggregate\_summary}. Default: 20.
\end{ldescription}
\end{Arguments}
%
\begin{Value}
A named list of class \code{"es\_diagnostics"} with six sections:
\begin{description}

\item[\code{meta}] List: \code{n\_events\_total} (integer), \code{n\_events\_shown}
(integer), \code{n\_events\_summarized} (integer), \code{event\_ids\_shown}
(integer vector).
\item[\code{estimation\_window}] List of plain numeric vectors (length
\code{n\_events\_shown}): \code{r2}, \code{sigma}, \code{degree\_of\_freedom},
\code{shapiro\_p}, \code{dw\_stat}, \code{ljung\_box\_p}, \code{acf1}.
Each entry is \code{NA\_real\_} when the model was not fitted or residuals
are insufficient for the test.
\item[\code{event\_window}] List of plain numeric vectors (length
\code{n\_events\_shown}): \code{ar\_t} (last event-window day AR t-stat),
\code{ar\_p} (two-sided p-value), \code{car\_t} (full-window CAR t-stat),
\code{car\_p} (two-sided p-value), \code{final\_car}. All \code{NA\_real\_}
when the \code{ART}/\code{CART} columns are absent.
\item[\code{cross\_sectional}] List of scalars aggregated across ALL events:
\code{n\_events} (total), \code{n\_valid\_events} (fitted count),
\code{car\_iqr}, \code{car\_sd}, \code{n\_overlap\_pairs}, \code{any\_overlap}.
Multi-event fields degrade to \code{NA} when \code{task\$aar\_caar\_tbl}
is \code{NULL}.
\item[\code{contract\_state}] List of logical/numeric vectors (length
\code{n\_events\_shown}): \code{is\_fitted}, \code{na\_ar\_count},
\code{na\_est\_count}, \code{insufficient\_obs}, \code{zero\_var\_index}.
\item[\code{aggregate\_summary}] Named list summarising remainder events
(those beyond \code{max\_events}): \code{n\_summarized},
\code{mean\_r2}, \code{median\_r2}, \code{mean\_final\_car},
\code{n\_fitted}, \code{n\_degenerate}. \code{NULL} when no events
are summarised (all shown).

\end{description}

\end{Value}
%
\begin{SeeAlso}
\code{\LinkA{model\_diagnostics}{model.Rul.diagnostics}}, \code{\LinkA{recommend\_stat}{recommend.Rul.stat}},
\code{\LinkA{flag\_robustness}{flag.Rul.robustness}}
\end{SeeAlso}
\HeaderA{es\_kb}{Access the EventStudy Grounding Knowledge Base}{es.Rul.kb}
%
\begin{Description}
Returns \code{EVENTSTUDY\_KB}, the package-level pure-R list of rule records
that maps diagnostic conditions to grounded methodological recommendations.
Each rule carries a structured academic citation and can be evaluated against
an \code{\LinkA{es\_diagnostics}{es.Rul.diagnostics}} object.
\end{Description}
%
\begin{Usage}
\begin{verbatim}
es_kb()
\end{verbatim}
\end{Usage}
%
\begin{Details}
The KB is the correctness-critical layer consumed by
\code{\LinkA{recommend\_stat}{recommend.Rul.stat}} and \code{\LinkA{flag\_robustness}{flag.Rul.robustness}}, and
exported here so that Phase 7 can inject its contents into an LLM
system prompt without accessing an internal package object.

\strong{KB-04 note:} This function delivers the \emph{structure} only —
exported and serializable, ready for Phase 7 system-prompt injection.
The actual prompt-injection behavior is a Phase 7 deliverable and is
deliberately out of scope here.
\end{Details}
%
\begin{Value}
A named list of rule records, each with fields:
\code{id}, \code{category}, \code{condition}, \code{recommendation},
\code{citation} (list of \code{author}/\code{year}/\code{key}/\code{venue}),
and \code{severity}.
\end{Value}
%
\begin{Note}
Each rule's \code{condition} field is an R function closure and is
therefore \strong{not} JSON-serializable. Any consumer that serializes
the KB (e.g. Phase 7 system-prompt injection) must drop the
\code{condition} field from each rule record before encoding to JSON.
\end{Note}
%
\begin{SeeAlso}
\code{\LinkA{recommend\_stat}{recommend.Rul.stat}}, \code{\LinkA{flag\_robustness}{flag.Rul.robustness}},
\code{\LinkA{es\_diagnostics}{es.Rul.diagnostics}}
\end{SeeAlso}
%
\begin{Examples}
\begin{ExampleCode}
kb <- es_kb()
length(kb)             # number of rules
kb[[1]]$id             # id of first rule
kb[[1]]$citation       # citation list

\end{ExampleCode}
\end{Examples}
\HeaderA{EventStudyTask}{Initialization of an Event Study task.}{EventStudyTask}
%
\begin{Description}
Event Study Task

The Event Study task contains all necessary data for performing an event
study. Furthermore, all calculations are saved in the internal dataframe
named data\_tbl.
\end{Description}
%
\begin{Section}{Public fields}

\begin{description}

\item[\code{data\_tbl}] All calculations are saved in
this dataframe. The dataframe is a nested
object with each row is one event, identified
by the event id, the group, and the firm
symbol.

\item[\code{aar\_caar\_tbl}] Placeholder for AAR and
CAAR test statistics.

\item[\code{.keys}] event identifier. Do not change.

\item[\code{.index}] The time column name

\item[\code{.target}] The price column name.

\item[\code{.request\_file\_columns}] The necessary
column names of the request dataframe.

\item[\code{factor\_tbl}] Optional factor data for multi-factor models
(Fama-French, Carhart). Must contain a \code{date} column.

\end{description}


\end{Section}
%
\begin{Section}{Active bindings}

\begin{description}

\item[\code{symbols}] Read-only field to get the firm symbols.

\item[\code{symbol\_data}] Read-only field to get the firm symbol data.

\item[\code{group\_level\_data}] Read-only field to get the group data.

\end{description}


\end{Section}
%
\begin{Section}{Methods}
%
\begin{SubSection}{Public methods}
\begin{itemize}

\item{} \Rhref{\#method-EventStudyTask-initialize}{\code{EventStudyTask\$new()}}
\item{} \Rhref{\#method-EventStudyTask-print}{\code{EventStudyTask\$print()}}
\item{} \Rhref{\#method-EventStudyTask-summary}{\code{EventStudyTask\$summary()}}
\item{} \Rhref{\#method-EventStudyTask-get_ar}{\code{EventStudyTask\$get\_ar()}}
\item{} \Rhref{\#method-EventStudyTask-get_car}{\code{EventStudyTask\$get\_car()}}
\item{} \Rhref{\#method-EventStudyTask-get_aar}{\code{EventStudyTask\$get\_aar()}}
\item{} \Rhref{\#method-EventStudyTask-get_model_stats}{\code{EventStudyTask\$get\_model\_stats()}}
\item{} \Rhref{\#method-EventStudyTask-clone}{\code{EventStudyTask\$clone()}}

\end{itemize}

\end{SubSection}



\hypertarget{method-EventStudyTask-initialize}{}
%
\begin{SubSection}{\code{EventStudyTask\$new()}}
Create a new EventStudyTask from stock data, reference
market data, and an event request specification.
%
\begin{SubSubSection}{Usage}

\begin{alltt}EventStudyTask$new(
  firm_stock_data_tbl,
  reference_tbl,
  request_tbl,
  factor_tbl = NULL
)\end{alltt}



\end{SubSubSection}

%
\begin{SubSubSection}{Arguments}

\begin{description}

\item[\code{firm\_stock\_data\_tbl}] Dataframe with firm
stock data. This dataframe must contain the
stock (col: symbol) the date (col: date)
and the price (col: adjusted) column.
\item[\code{reference\_tbl}] Dataframe with firm
reference data.
\item[\code{request\_tbl}] The request dataframe for
each event.
\item[\code{factor\_tbl}] Optional dataframe with factor data
(e.g., Fama-French factors). Must contain a \code{date}
column plus factor columns.

\end{description}



\end{SubSubSection}

\end{SubSection}




\hypertarget{method-EventStudyTask-print}{}
%
\begin{SubSection}{\code{EventStudyTask\$print()}}
Print a summary of the EventStudyTask.
%
\begin{SubSubSection}{Usage}

\begin{alltt}EventStudyTask$print(...)\end{alltt}



\end{SubSubSection}

\end{SubSection}




\hypertarget{method-EventStudyTask-summary}{}
%
\begin{SubSection}{\code{EventStudyTask\$summary()}}
Summarize the event study results.
%
\begin{SubSubSection}{Usage}

\begin{alltt}EventStudyTask$summary()\end{alltt}



\end{SubSubSection}

%
\begin{SubSubSection}{Returns}
A list with summary information about the event study results.

\end{SubSubSection}

\end{SubSection}




\hypertarget{method-EventStudyTask-get_ar}{}
%
\begin{SubSection}{\code{EventStudyTask\$get\_ar()}}
Extract abnormal returns for a single event.
%
\begin{SubSubSection}{Usage}

\begin{alltt}EventStudyTask$get_ar(event_id = NULL)\end{alltt}



\end{SubSubSection}

%
\begin{SubSubSection}{Arguments}

\begin{description}

\item[\code{event\_id}] The event identifier.

\end{description}



\end{SubSubSection}

%
\begin{SubSubSection}{Returns}
A tibble with abnormal returns for the event window.

\end{SubSubSection}

\end{SubSection}




\hypertarget{method-EventStudyTask-get_car}{}
%
\begin{SubSection}{\code{EventStudyTask\$get\_car()}}
Extract cumulative abnormal returns for a single event.
%
\begin{SubSubSection}{Usage}

\begin{alltt}EventStudyTask$get_car(event_id = NULL)\end{alltt}



\end{SubSubSection}

%
\begin{SubSubSection}{Arguments}

\begin{description}

\item[\code{event\_id}] The event identifier.

\end{description}



\end{SubSubSection}

%
\begin{SubSubSection}{Returns}
A tibble with cumulative abnormal returns for the event window.

\end{SubSubSection}

\end{SubSection}




\hypertarget{method-EventStudyTask-get_aar}{}
%
\begin{SubSection}{\code{EventStudyTask\$get\_aar()}}
Extract average abnormal returns across events.
%
\begin{SubSubSection}{Usage}

\begin{alltt}EventStudyTask$get_aar(group = NULL, stat_name = "CSectT")\end{alltt}



\end{SubSubSection}

%
\begin{SubSubSection}{Arguments}

\begin{description}

\item[\code{group}] Optional group name to filter by.
\item[\code{stat\_name}] Name of the multi-event test statistic
to extract. Defaults to "CSectT".

\end{description}



\end{SubSubSection}

%
\begin{SubSubSection}{Returns}
A tibble with AAR/CAAR results.

\end{SubSubSection}

\end{SubSection}




\hypertarget{method-EventStudyTask-get_model_stats}{}
%
\begin{SubSection}{\code{EventStudyTask\$get\_model\_stats()}}
Extract model statistics for a single event.
%
\begin{SubSubSection}{Usage}

\begin{alltt}EventStudyTask$get_model_stats(event_id = NULL)\end{alltt}



\end{SubSubSection}

%
\begin{SubSubSection}{Arguments}

\begin{description}

\item[\code{event\_id}] The event identifier.

\end{description}



\end{SubSubSection}

%
\begin{SubSubSection}{Returns}
A list with model statistics.

\end{SubSubSection}

\end{SubSection}




\hypertarget{method-EventStudyTask-clone}{}
%
\begin{SubSection}{\code{EventStudyTask\$clone()}}
The objects of this class are cloneable with this method.
%
\begin{SubSubSection}{Usage}

\begin{alltt}EventStudyTask$clone(deep = FALSE)\end{alltt}



\end{SubSubSection}

%
\begin{SubSubSection}{Arguments}

\begin{description}

\item[\code{deep}] Whether to make a deep clone.

\end{description}



\end{SubSubSection}

\end{SubSection}


\end{Section}
\HeaderA{EVENTSTUDY\_KB}{Grounding Knowledge Base for Event Study Diagnostics}{EVENTSTUDY.Rul.KB}
\keyword{internal}{EVENTSTUDY\_KB}
%
\begin{Description}
\code{EVENTSTUDY\_KB} is a pure-R list of rule records that maps diagnostic
conditions (read from an \code{\LinkA{es\_diagnostics}{es.Rul.diagnostics}} object) to grounded
methodological recommendations, each carrying a structured academic citation
and a severity level.

Package-level list; use \code{\LinkA{es\_kb}{es.Rul.kb}()} to access.
\end{Description}
%
\begin{Usage}
\begin{verbatim}
EVENTSTUDY_KB
\end{verbatim}
\end{Usage}
%
\begin{Details}
The KB is built at package-load time from calls to \code{.kb\_rule()} and
never depends on a network connection or LLM provider. It is the
correctness-critical layer that the offline advice functions
(\code{\LinkA{recommend\_stat}{recommend.Rul.stat}}, \code{\LinkA{flag\_robustness}{flag.Rul.robustness}}) and the
Phase 7 LLM layer both evaluate against.

\strong{KB-04 scope note:} Phase 5 delivers the KB data structure
only — exported and serializable, ready for Phase 7 system-prompt
injection. The actual prompt-injection behavior is a Phase 7
deliverable and is deliberately out of scope here.
\end{Details}
%
\begin{Section}{Rule record fields}

Each element of \code{EVENTSTUDY\_KB} is a list with:
\begin{description}

\item[\code{id}] Unique character identifier (e.g. \code{"KB-NORM-PATELL"}).
\item[\code{category}] Either \code{"stat\_choice"} (steers which test to
use) or \code{"robustness"} (data-quality / reliability warning).
\item[\code{condition}] A \code{function(diag)} that accepts an
\code{es\_diagnostics} list and returns a length-1 logical. Returns
\code{FALSE} when any required field is \code{NA} (NA-safe; never
errors on missing data).
\item[\code{recommendation}] Character. The methodological action to take.
\item[\code{citation}] Named list: \code{author} (character),
\code{year} (integer), \code{key} (character), \code{venue} (character).
\item[\code{severity}] One of \code{"info"}, \code{"warning"},
\code{"error"}.

\end{description}

\end{Section}
%
\begin{Section}{Threshold notes}

The following thresholds are literature-informed defaults and are marked
\strong{[ASSUMED]} — they can be tuned without breaking the KB contract:
\begin{itemize}

\item{} Shapiro-Wilk p-value threshold: 0.05
\item{} Proportion for normality-holds rule: \bsl{}>= 70\% of events
\item{} Proportion for non-normality rule: \bsl{}>= 50\% of events
\item{} Durbin-Watson bounds: [1.5, 2.5]
\item{} R-squared low-fit threshold: < 0.05
\item{} Small-N threshold: < 10 valid events
\item{} Degenerate-event threshold: < 80\% of events fitted
\item{} High-dispersion CAR IQR threshold: > 0.10 or SD > 0.15

\end{itemize}

\end{Section}
\HeaderA{export\_results}{Export Event Study Results}{export.Rul.results}
%
\begin{Description}
Export event study results to CSV, Excel, or LaTeX formats.
\end{Description}
%
\begin{Usage}
\begin{verbatim}
export_results(
  task,
  file,
  format = NULL,
  which = c("ar", "car", "aar", "model"),
  stat_name = "CSectT",
  ...
)
\end{verbatim}
\end{Usage}
%
\begin{Arguments}
\begin{ldescription}
\item[\code{task}] A fitted EventStudyTask with statistics computed.

\item[\code{file}] Path to the output file. The file extension determines the format
unless \code{format} is explicitly specified.

\item[\code{format}] Output format: "csv", "xlsx", or "latex". If NULL, inferred
from the file extension.

\item[\code{which}] Which results to export. One or more of "ar", "car", "aar",
"model". Defaults to all available.

\item[\code{stat\_name}] Name of the multi-event test statistic to use for AAR/CAAR
tables. Defaults to "CSectT".

\item[\code{...}] Additional arguments passed to the format-specific writer.
\end{ldescription}
\end{Arguments}
%
\begin{Value}
The file path (invisibly).
\end{Value}
\HeaderA{FamaFrench3FactorModel}{Fama-French Three-Factor Model}{FamaFrench3FactorModel}
%
\begin{Description}
Implements the Fama and French (1993) three-factor model:
\deqn{R_i - R_f = \alpha + \beta_m (R_m - R_f) + \beta_s SMB + \beta_h HML + \epsilon}{}

The data must contain columns: \code{excess\_return} (firm return minus risk-free),
\code{market\_excess} (market return minus risk-free), \code{smb}, and \code{hml}.
These can be joined via a factor table in \code{EventStudyTask}.
\end{Description}
%
\begin{Section}{Super classes}
\code{\LinkA{ModelBase}{ModelBase}} -> \code{\LinkA{LinearFactorModel}{LinearFactorModel}} -> \code{FamaFrench3FactorModel}
\end{Section}
%
\begin{Section}{Public fields}

\begin{description}

\item[\code{model\_name}] Name of the model.

\item[\code{formula}] The three-factor regression formula.

\item[\code{required\_columns}] Required data columns.

\end{description}


\end{Section}
%
\begin{Section}{Methods}
%
\begin{SubSection}{Public methods}
\begin{itemize}

\item{} \Rhref{\#method-FamaFrench3FactorModel-initialize}{\code{FamaFrench3FactorModel\$new()}}
\item{} \Rhref{\#method-FamaFrench3FactorModel-abnormal_returns}{\code{FamaFrench3FactorModel\$abnormal\_returns()}}
\item{} \Rhref{\#method-FamaFrench3FactorModel-clone}{\code{FamaFrench3FactorModel\$clone()}}

\end{itemize}

\end{SubSection}




\hypertarget{method-FamaFrench3FactorModel-initialize}{}
%
\begin{SubSection}{\code{FamaFrench3FactorModel\$new()}}
Create a new FamaFrench3FactorModel.
%
\begin{SubSubSection}{Usage}

\begin{alltt}FamaFrench3FactorModel$new(use_hac = FALSE, hac_lag = NULL)\end{alltt}



\end{SubSubSection}

%
\begin{SubSubSection}{Arguments}

\begin{description}

\item[\code{use\_hac}] Logical. Use HAC (Newey-West) standard errors.
\item[\code{hac\_lag}] Integer or NULL. Lag truncation for Newey-West.

\end{description}



\end{SubSubSection}

\end{SubSection}




\hypertarget{method-FamaFrench3FactorModel-abnormal_returns}{}
%
\begin{SubSection}{\code{FamaFrench3FactorModel\$abnormal\_returns()}}
Calculate abnormal returns using the three-factor model.
%
\begin{SubSubSection}{Usage}

\begin{alltt}FamaFrench3FactorModel$abnormal_returns(data_tbl)\end{alltt}



\end{SubSubSection}

%
\begin{SubSubSection}{Arguments}

\begin{description}

\item[\code{data\_tbl}] Data frame or tibble.

\end{description}



\end{SubSubSection}

\end{SubSection}




\hypertarget{method-FamaFrench3FactorModel-clone}{}
%
\begin{SubSection}{\code{FamaFrench3FactorModel\$clone()}}
The objects of this class are cloneable with this method.
%
\begin{SubSubSection}{Usage}

\begin{alltt}FamaFrench3FactorModel$clone(deep = FALSE)\end{alltt}



\end{SubSubSection}

%
\begin{SubSubSection}{Arguments}

\begin{description}

\item[\code{deep}] Whether to make a deep clone.

\end{description}



\end{SubSubSection}

\end{SubSection}


\end{Section}
\HeaderA{FamaFrench5FactorModel}{Fama-French Five-Factor Model}{FamaFrench5FactorModel}
%
\begin{Description}
Implements the Fama and French (2015) five-factor model:
\deqn{R_i - R_f = \alpha + \beta_m (R_m - R_f) + \beta_s SMB + \beta_h HML + \beta_r RMW + \beta_c CMA + \epsilon}{}

Requires columns: \code{excess\_return}, \code{market\_excess}, \code{smb},
\code{hml}, \code{rmw}, \code{cma}.
\end{Description}
%
\begin{Section}{Super classes}
\code{\LinkA{ModelBase}{ModelBase}} -> \code{\LinkA{LinearFactorModel}{LinearFactorModel}} -> \code{FamaFrench5FactorModel}
\end{Section}
%
\begin{Section}{Public fields}

\begin{description}

\item[\code{model\_name}] Name of the model.

\item[\code{formula}] The five-factor regression formula.

\item[\code{required\_columns}] Required data columns.

\end{description}


\end{Section}
%
\begin{Section}{Methods}
%
\begin{SubSection}{Public methods}
\begin{itemize}

\item{} \Rhref{\#method-FamaFrench5FactorModel-initialize}{\code{FamaFrench5FactorModel\$new()}}
\item{} \Rhref{\#method-FamaFrench5FactorModel-abnormal_returns}{\code{FamaFrench5FactorModel\$abnormal\_returns()}}
\item{} \Rhref{\#method-FamaFrench5FactorModel-clone}{\code{FamaFrench5FactorModel\$clone()}}

\end{itemize}

\end{SubSection}




\hypertarget{method-FamaFrench5FactorModel-initialize}{}
%
\begin{SubSection}{\code{FamaFrench5FactorModel\$new()}}
Create a new FamaFrench5FactorModel.
%
\begin{SubSubSection}{Usage}

\begin{alltt}FamaFrench5FactorModel$new(use_hac = FALSE, hac_lag = NULL)\end{alltt}



\end{SubSubSection}

%
\begin{SubSubSection}{Arguments}

\begin{description}

\item[\code{use\_hac}] Logical. Use HAC (Newey-West) standard errors.
\item[\code{hac\_lag}] Integer or NULL. Lag truncation for Newey-West.

\end{description}



\end{SubSubSection}

\end{SubSection}




\hypertarget{method-FamaFrench5FactorModel-abnormal_returns}{}
%
\begin{SubSection}{\code{FamaFrench5FactorModel\$abnormal\_returns()}}
Calculate abnormal returns using the five-factor model.
%
\begin{SubSubSection}{Usage}

\begin{alltt}FamaFrench5FactorModel$abnormal_returns(data_tbl)\end{alltt}



\end{SubSubSection}

%
\begin{SubSubSection}{Arguments}

\begin{description}

\item[\code{data\_tbl}] Data frame or tibble.

\end{description}



\end{SubSubSection}

\end{SubSection}




\hypertarget{method-FamaFrench5FactorModel-clone}{}
%
\begin{SubSection}{\code{FamaFrench5FactorModel\$clone()}}
The objects of this class are cloneable with this method.
%
\begin{SubSubSection}{Usage}

\begin{alltt}FamaFrench5FactorModel$clone(deep = FALSE)\end{alltt}



\end{SubSubSection}

%
\begin{SubSubSection}{Arguments}

\begin{description}

\item[\code{deep}] Whether to make a deep clone.

\end{description}



\end{SubSubSection}

\end{SubSection}


\end{Section}
\HeaderA{fit\_model}{Train the defined model on each event and calculate the abnormal return.}{fit.Rul.model}
%
\begin{Description}
Train the defined model on each event and calculate the abnormal return.
\end{Description}
%
\begin{Usage}
\begin{verbatim}
fit_model(task, parameter_set)
\end{verbatim}
\end{Usage}
%
\begin{Arguments}
\begin{ldescription}
\item[\code{task}] The event study task

\item[\code{parameter\_set}] The parameter set that defines the event study.
\end{ldescription}
\end{Arguments}
%
\begin{Value}
task
\end{Value}
\HeaderA{flag\_robustness}{Flag Robustness Issues via Offline KB Matching}{flag.Rul.robustness}
\methaliasA{flag\_robustness.default}{flag\_robustness}{flag.Rul.robustness.default}
\methaliasA{flag\_robustness.es\_diagnostics}{flag\_robustness}{flag.Rul.robustness.es.Rul.diagnostics}
\methaliasA{flag\_robustness.EventStudyTask}{flag\_robustness}{flag.Rul.robustness.EventStudyTask}
%
\begin{Description}
Evaluates the KB decision table (category \code{"robustness"} rules) against
the diagnostic signals extracted from a fitted \code{EventStudyTask} or a
precomputed \code{es\_diagnostics} object, and returns a severity-ranked
\code{es\_advice} S3 object.
\end{Description}
%
\begin{Usage}
\begin{verbatim}
flag_robustness(x, provider = NULL, ...)

## Default S3 method:
flag_robustness(x, provider = NULL, ...)

## S3 method for class 'EventStudyTask'
flag_robustness(x, provider = NULL, ...)

## S3 method for class 'es_diagnostics'
flag_robustness(x, provider = NULL, ...)
\end{verbatim}
\end{Usage}
%
\begin{Arguments}
\begin{ldescription}
\item[\code{x}] A fitted \code{EventStudyTask} (after \code{fit\_model()}) or a
precomputed \code{es\_diagnostics} object returned by \code{es\_diagnostics()}.

\item[\code{provider}] Accepted but ignored in the offline path — present only so
the Phase 7 call shape is forward-compatible. Default \code{NULL}.

\item[\code{...}] Additional arguments (currently ignored).
\end{ldescription}
\end{Arguments}
%
\begin{Details}
No LLM provider, network connection, or API key is required. The returned
\code{es\_advice} object has the same shape as the Phase 7 Advice contract,
flagged \code{is\_deterministic = TRUE} and \code{source = "offline\_kb"}.
\end{Details}
%
\begin{Value}
An S3 object of class \code{"es\_advice"} — see \code{\LinkA{recommend\_stat}{recommend.Rul.stat}}
for field descriptions. Rules are filtered to \code{category == "robustness"}.
\end{Value}
%
\begin{SeeAlso}
\code{\LinkA{recommend\_stat}{recommend.Rul.stat}}, \code{\LinkA{es\_diagnostics}{es.Rul.diagnostics}},
\code{\LinkA{es\_kb}{es.Rul.kb}}
\end{SeeAlso}
%
\begin{Examples}
\begin{ExampleCode}
## Not run: 
task <- run_event_study(my_task, ParameterSet$new())
advice <- flag_robustness(task)
print(advice)

## End(Not run)

\end{ExampleCode}
\end{Examples}
\HeaderA{GARCHModel}{GARCH Model}{GARCHModel}
%
\begin{Description}
Event study model using GARCH(1,1) for time-varying volatility estimation.
Uses the \pkg{rugarch} package to fit a GARCH(1,1) model with a market
return regressor in the mean equation during the estimation window.
Abnormal returns are computed as the difference between observed returns
and the GARCH conditional mean. The time-varying sigma from GARCH can
be used for standardized test statistics.
\end{Description}
%
\begin{Section}{Super class}
\code{\LinkA{ModelBase}{ModelBase}} -> \code{GARCHModel}
\end{Section}
%
\begin{Section}{Public fields}

\begin{description}

\item[\code{model\_name}] Name of the model.

\item[\code{garch\_order}] GARCH order as c(p, q). Default c(1,1).

\end{description}


\end{Section}
%
\begin{Section}{Methods}
%
\begin{SubSection}{Public methods}
\begin{itemize}

\item{} \Rhref{\#method-GARCHModel-fit}{\code{GARCHModel\$fit()}}
\item{} \Rhref{\#method-GARCHModel-abnormal_returns}{\code{GARCHModel\$abnormal\_returns()}}
\item{} \Rhref{\#method-GARCHModel-clone}{\code{GARCHModel\$clone()}}

\end{itemize}

\end{SubSection}



\hypertarget{method-GARCHModel-fit}{}
%
\begin{SubSection}{\code{GARCHModel\$fit()}}
Fit the GARCH model on the estimation window.
%
\begin{SubSubSection}{Usage}

\begin{alltt}GARCHModel$fit(data_tbl)\end{alltt}



\end{SubSubSection}

%
\begin{SubSubSection}{Arguments}

\begin{description}

\item[\code{data\_tbl}] Data frame or tibble with firm\_returns,
index\_returns, estimation\_window, event\_window columns.

\end{description}



\end{SubSubSection}

\end{SubSection}




\hypertarget{method-GARCHModel-abnormal_returns}{}
%
\begin{SubSection}{\code{GARCHModel\$abnormal\_returns()}}
Calculate abnormal returns from the GARCH model.
%
\begin{SubSubSection}{Usage}

\begin{alltt}GARCHModel$abnormal_returns(data_tbl)\end{alltt}



\end{SubSubSection}

%
\begin{SubSubSection}{Arguments}

\begin{description}

\item[\code{data\_tbl}] Data frame or tibble.

\end{description}



\end{SubSubSection}

\end{SubSection}




\hypertarget{method-GARCHModel-clone}{}
%
\begin{SubSection}{\code{GARCHModel\$clone()}}
The objects of this class are cloneable with this method.
%
\begin{SubSubSection}{Usage}

\begin{alltt}GARCHModel$clone(deep = FALSE)\end{alltt}



\end{SubSubSection}

%
\begin{SubSubSection}{Arguments}

\begin{description}

\item[\code{deep}] Whether to make a deep clone.

\end{description}



\end{SubSubSection}

\end{SubSection}


\end{Section}
\HeaderA{GeneralizedSignTest}{Generalized Sign Test (Cowan 1992)}{GeneralizedSignTest}
%
\begin{Description}
Adjusts the sign test for the expected proportion of positive abnormal
returns estimated from the estimation window, rather than assuming 0.5.
This accounts for asymmetry in the return distribution.
\end{Description}
%
\begin{Section}{Super class}
\code{\LinkA{TestStatisticBase}{TestStatisticBase}} -> \code{GeneralizedSignTest}
\end{Section}
%
\begin{Section}{Public fields}

\begin{description}

\item[\code{name}] Short code of the test statistic.

\end{description}


\end{Section}
%
\begin{Section}{Methods}
%
\begin{SubSection}{Public methods}
\begin{itemize}

\item{} \Rhref{\#method-GeneralizedSignTest-compute}{\code{GeneralizedSignTest\$compute()}}
\item{} \Rhref{\#method-GeneralizedSignTest-clone}{\code{GeneralizedSignTest\$clone()}}

\end{itemize}

\end{SubSection}




\hypertarget{method-GeneralizedSignTest-compute}{}
%
\begin{SubSection}{\code{GeneralizedSignTest\$compute()}}
Computes the generalized sign test for multiple events.
%
\begin{SubSubSection}{Usage}

\begin{alltt}GeneralizedSignTest$compute(data_tbl, model)\end{alltt}



\end{SubSubSection}

%
\begin{SubSubSection}{Arguments}

\begin{description}

\item[\code{data\_tbl}] The data for a multiple event with
calculated abnormal returns.
\item[\code{model}] The fitted model (unused).

\end{description}



\end{SubSubSection}

\end{SubSection}




\hypertarget{method-GeneralizedSignTest-clone}{}
%
\begin{SubSection}{\code{GeneralizedSignTest\$clone()}}
The objects of this class are cloneable with this method.
%
\begin{SubSubSection}{Usage}

\begin{alltt}GeneralizedSignTest$clone(deep = FALSE)\end{alltt}



\end{SubSubSection}

%
\begin{SubSubSection}{Arguments}

\begin{description}

\item[\code{deep}] Whether to make a deep clone.

\end{description}



\end{SubSubSection}

\end{SubSection}


\end{Section}
\HeaderA{generate\_report}{Generate Event Study Report}{generate.Rul.report}
%
\begin{Description}
Renders an automated HTML or PDF report from a completed event study task.
Uses a bundled RMarkdown template with configurable sections.
\end{Description}
%
\begin{Usage}
\begin{verbatim}
generate_report(
  task,
  output_file = "event_study_report.html",
  format = c("html", "pdf"),
  title = "Event Study Report",
  author = NULL,
  sections = c("summary", "data", "diagnostics", "single_event", "multi_event",
    "cross_sectional", "appendix"),
  cross_sectional = NULL,
  confidence_level = 0.95,
  interactive = TRUE,
  advice = NULL,
  ...
)
\end{verbatim}
\end{Usage}
%
\begin{Arguments}
\begin{ldescription}
\item[\code{task}] A fitted \code{EventStudyTask} or \code{PanelEventStudyTask}.

\item[\code{output\_file}] Output file path. Default \code{"event\_study\_report.html"}.

\item[\code{format}] Output format: \code{"html"} (default) or \code{"pdf"}.

\item[\code{title}] Report title.

\item[\code{author}] Author name (optional).

\item[\code{sections}] Character vector of sections to include. Any subset of:
\code{"summary"}, \code{"data"}, \code{"diagnostics"}, \code{"single\_event"},
\code{"multi\_event"}, \code{"cross\_sectional"}, \code{"appendix"}.

\item[\code{cross\_sectional}] Optional cross-sectional regression results to include.

\item[\code{confidence\_level}] Confidence level for plots. Default 0.95.

\item[\code{interactive}] Logical. Use interactive plotly plots in HTML output.
Default TRUE.

\item[\code{advice}] An optional grounded \code{Advice} object returned by
\code{\LinkA{es\_advise}{es.Rul.advise}(task\_type = "report\_writing")}. When supplied and
a valid \code{Advice}, renders a new \strong{AI Advisor Interpretation}
section in the report. When \code{NULL} (the default), the existing render
path is completely unchanged (byte-identical output). A supplied but
invalid \code{advice} (not an \code{Advice} object) is silently coerced to
\code{NULL} with exactly one \code{warning()} — the report is never broken.

\item[\code{...}] Additional arguments passed to \code{rmarkdown::render}.
\end{ldescription}
\end{Arguments}
%
\begin{Value}
The path to the generated report (invisibly).
\end{Value}
\HeaderA{IntradayEventStudyTask}{Intraday Event Study Task}{IntradayEventStudyTask}
%
\begin{Description}
Extension of \code{\LinkA{EventStudyTask}{EventStudyTask}} for intraday (high-frequency)
event studies. Supports POSIXct timestamps and time-based windows
specified in minutes or hours.
\end{Description}
%
\begin{Section}{Differences from EventStudyTask}

\begin{itemize}

\item{} Date column contains POSIXct timestamps instead of date strings
\item{} Event windows specified as time offsets (e.g., \code{"-30min"}, \code{"+60min"})
\item{} Handles trading session boundaries

\end{itemize}

\end{Section}
%
\begin{Section}{Super class}
\code{\LinkA{EventStudyTask}{EventStudyTask}} -> \code{IntradayEventStudyTask}
\end{Section}
%
\begin{Section}{Public fields}

\begin{description}

\item[\code{.index}] The time column name (POSIXct).

\item[\code{.target}] The price column name.

\item[\code{.request\_file\_columns}] Required columns for intraday requests.

\end{description}


\end{Section}
%
\begin{Section}{Methods}
%
\begin{SubSection}{Public methods}
\begin{itemize}

\item{} \Rhref{\#method-IntradayEventStudyTask-initialize}{\code{IntradayEventStudyTask\$new()}}
\item{} \Rhref{\#method-IntradayEventStudyTask-print}{\code{IntradayEventStudyTask\$print()}}
\item{} \Rhref{\#method-IntradayEventStudyTask-clone}{\code{IntradayEventStudyTask\$clone()}}

\end{itemize}

\end{SubSection}




\hypertarget{method-IntradayEventStudyTask-initialize}{}
%
\begin{SubSection}{\code{IntradayEventStudyTask\$new()}}
Create a new IntradayEventStudyTask.
%
\begin{SubSubSection}{Usage}

\begin{alltt}IntradayEventStudyTask$new(
  firm_stock_data_tbl,
  reference_tbl,
  request_tbl,
  factor_tbl = NULL
)\end{alltt}



\end{SubSubSection}

%
\begin{SubSubSection}{Arguments}

\begin{description}

\item[\code{firm\_stock\_data\_tbl}] Dataframe with firm intraday data.
Must contain columns: \code{symbol}, \code{timestamp} (POSIXct),
and \code{price}.
\item[\code{reference\_tbl}] Dataframe with reference market intraday data.
\item[\code{request\_tbl}] Request dataframe with \code{event\_timestamp}
(POSIXct), and window specifications in minutes.
\item[\code{factor\_tbl}] Optional factor data.

\end{description}



\end{SubSubSection}

\end{SubSection}




\hypertarget{method-IntradayEventStudyTask-print}{}
%
\begin{SubSection}{\code{IntradayEventStudyTask\$print()}}
Print summary.
%
\begin{SubSubSection}{Usage}

\begin{alltt}IntradayEventStudyTask$print(...)\end{alltt}



\end{SubSubSection}

\end{SubSection}




\hypertarget{method-IntradayEventStudyTask-clone}{}
%
\begin{SubSection}{\code{IntradayEventStudyTask\$clone()}}
The objects of this class are cloneable with this method.
%
\begin{SubSubSection}{Usage}

\begin{alltt}IntradayEventStudyTask$clone(deep = FALSE)\end{alltt}



\end{SubSubSection}

%
\begin{SubSubSection}{Arguments}

\begin{description}

\item[\code{deep}] Whether to make a deep clone.

\end{description}



\end{SubSubSection}

\end{SubSection}


\end{Section}
\HeaderA{KolariPynnonenTest}{Kolari-Pynnönen Adjusted BMP Test}{KolariPynnonenTest}
%
\begin{Description}
Adjusts the BMP (Boehmer, Musumeci, Poulsen 1991) test for cross-sectional
correlation of abnormal returns using the Kolari and Pynnönen (2010)
correction. The adjustment scales the BMP statistic by a factor that
accounts for the average pairwise correlation of standardized abnormal
residuals in the estimation window.
\end{Description}
%
\begin{Section}{Super class}
\code{\LinkA{TestStatisticBase}{TestStatisticBase}} -> \code{KolariPynnonenTest}
\end{Section}
%
\begin{Section}{Public fields}

\begin{description}

\item[\code{name}] Short code of the test statistic.

\end{description}


\end{Section}
%
\begin{Section}{Methods}
%
\begin{SubSection}{Public methods}
\begin{itemize}

\item{} \Rhref{\#method-KolariPynnonenTest-compute}{\code{KolariPynnonenTest\$compute()}}
\item{} \Rhref{\#method-KolariPynnonenTest-clone}{\code{KolariPynnonenTest\$clone()}}

\end{itemize}

\end{SubSection}




\hypertarget{method-KolariPynnonenTest-compute}{}
%
\begin{SubSection}{\code{KolariPynnonenTest\$compute()}}
Computes the Kolari-Pynnönen adjusted BMP test.
%
\begin{SubSubSection}{Usage}

\begin{alltt}KolariPynnonenTest$compute(data_tbl, model)\end{alltt}



\end{SubSubSection}

%
\begin{SubSubSection}{Arguments}

\begin{description}

\item[\code{data\_tbl}] The data for a multiple event with
calculated abnormal returns.
\item[\code{model}] The fitted model containing sigma estimates.

\end{description}



\end{SubSubSection}

\end{SubSection}




\hypertarget{method-KolariPynnonenTest-clone}{}
%
\begin{SubSection}{\code{KolariPynnonenTest\$clone()}}
The objects of this class are cloneable with this method.
%
\begin{SubSubSection}{Usage}

\begin{alltt}KolariPynnonenTest$clone(deep = FALSE)\end{alltt}



\end{SubSubSection}

%
\begin{SubSubSection}{Arguments}

\begin{description}

\item[\code{deep}] Whether to make a deep clone.

\end{description}



\end{SubSubSection}

\end{SubSection}


\end{Section}
%
\begin{References}
Kolari, J. W. and Pynnönen, S. (2010). Event Study Testing with
Cross-sectional Correlation of Abnormal Returns.
\emph{The Review of Financial Studies}, 23(11), 3996--4025.
\end{References}
\HeaderA{LinearFactorModel}{Linear Factor Model Base}{LinearFactorModel}
%
\begin{Description}
Base class for multi-factor OLS models used in event studies. All factor
models (Market Model, Fama-French, Carhart) share the same estimation
approach: OLS regression of excess returns on factor returns during the
estimation window.
\end{Description}
%
\begin{Section}{Super class}
\code{\LinkA{ModelBase}{ModelBase}} -> \code{LinearFactorModel}
\end{Section}
%
\begin{Section}{Public fields}

\begin{description}

\item[\code{model\_name}] Name of the model.

\item[\code{formula}] The regression formula.

\item[\code{required\_columns}] Columns required in the data.

\item[\code{use\_hac}] Logical. Use HAC (Newey-West) standard errors.

\item[\code{hac\_lag}] Integer or NULL. Lag truncation for Newey-West.

\end{description}


\end{Section}
%
\begin{Section}{Methods}
%
\begin{SubSection}{Public methods}
\begin{itemize}

\item{} \Rhref{\#method-LinearFactorModel-initialize}{\code{LinearFactorModel\$new()}}
\item{} \Rhref{\#method-LinearFactorModel-fit}{\code{LinearFactorModel\$fit()}}
\item{} \Rhref{\#method-LinearFactorModel-abnormal_returns}{\code{LinearFactorModel\$abnormal\_returns()}}
\item{} \Rhref{\#method-LinearFactorModel-clone}{\code{LinearFactorModel\$clone()}}

\end{itemize}

\end{SubSection}



\hypertarget{method-LinearFactorModel-initialize}{}
%
\begin{SubSection}{\code{LinearFactorModel\$new()}}
Create a new LinearFactorModel.
%
\begin{SubSubSection}{Usage}

\begin{alltt}LinearFactorModel$new(use_hac = FALSE, hac_lag = NULL)\end{alltt}



\end{SubSubSection}

%
\begin{SubSubSection}{Arguments}

\begin{description}

\item[\code{use\_hac}] Logical. Use HAC (Newey-West) standard errors.
\item[\code{hac\_lag}] Integer or NULL. Lag truncation for Newey-West.

\end{description}



\end{SubSubSection}

\end{SubSection}




\hypertarget{method-LinearFactorModel-fit}{}
%
\begin{SubSection}{\code{LinearFactorModel\$fit()}}
Fit the linear factor model via OLS on the estimation window.
%
\begin{SubSubSection}{Usage}

\begin{alltt}LinearFactorModel$fit(data_tbl)\end{alltt}



\end{SubSubSection}

%
\begin{SubSubSection}{Arguments}

\begin{description}

\item[\code{data\_tbl}] Data frame or tibble containing the data to fit.

\end{description}



\end{SubSubSection}

\end{SubSection}




\hypertarget{method-LinearFactorModel-abnormal_returns}{}
%
\begin{SubSection}{\code{LinearFactorModel\$abnormal\_returns()}}
Calculate abnormal returns as observed minus predicted.
%
\begin{SubSubSection}{Usage}

\begin{alltt}LinearFactorModel$abnormal_returns(data_tbl)\end{alltt}



\end{SubSubSection}

%
\begin{SubSubSection}{Arguments}

\begin{description}

\item[\code{data\_tbl}] Data frame or tibble.

\end{description}



\end{SubSubSection}

\end{SubSection}




\hypertarget{method-LinearFactorModel-clone}{}
%
\begin{SubSection}{\code{LinearFactorModel\$clone()}}
The objects of this class are cloneable with this method.
%
\begin{SubSubSection}{Usage}

\begin{alltt}LinearFactorModel$clone(deep = FALSE)\end{alltt}



\end{SubSubSection}

%
\begin{SubSubSection}{Arguments}

\begin{description}

\item[\code{deep}] Whether to make a deep clone.

\end{description}



\end{SubSubSection}

\end{SubSection}


\end{Section}
\HeaderA{LogReturn}{R6 class for log return calculation}{LogReturn}
%
\begin{Description}
R6 class for log return calculation
\end{Description}
%
\begin{Section}{Super class}
\code{\LinkA{ReturnCalculation}{ReturnCalculation}} -> \code{LogReturn}
\end{Section}
%
\begin{Section}{Public fields}

\begin{description}

\item[\code{name}] Name of the log return calculation.

\end{description}


\end{Section}
%
\begin{Section}{Methods}
%
\begin{SubSection}{Public methods}
\begin{itemize}

\item{} \Rhref{\#method-LogReturn-calculate_return}{\code{LogReturn\$calculate\_return()}}
\item{} \Rhref{\#method-LogReturn-clone}{\code{LogReturn\$clone()}}

\end{itemize}

\end{SubSection}



\hypertarget{method-LogReturn-calculate_return}{}
%
\begin{SubSection}{\code{LogReturn\$calculate\_return()}}
Calculates the return for a
single stock.
%
\begin{SubSubSection}{Usage}

\begin{alltt}LogReturn$calculate_return(
  tbl,
  in_column = "adjusted",
  out_column = "adjusted_return"
)\end{alltt}



\end{SubSubSection}

%
\begin{SubSubSection}{Arguments}

\begin{description}

\item[\code{tbl}] The dataframe with the stock price.
\item[\code{in\_column}] The column name of the price
infromation.
\item[\code{out\_column}] The column name were the
return will be saved.

\end{description}



\end{SubSubSection}

\end{SubSection}




\hypertarget{method-LogReturn-clone}{}
%
\begin{SubSection}{\code{LogReturn\$clone()}}
The objects of this class are cloneable with this method.
%
\begin{SubSubSection}{Usage}

\begin{alltt}LogReturn$clone(deep = FALSE)\end{alltt}



\end{SubSubSection}

%
\begin{SubSubSection}{Arguments}

\begin{description}

\item[\code{deep}] Whether to make a deep clone.

\end{description}



\end{SubSubSection}

\end{SubSection}


\end{Section}
\HeaderA{MarketAdjustedModel}{Market Adjusted Model}{MarketAdjustedModel}
%
\begin{Description}
The Market Adjusted Model is another simple approach used in event studies
to estimate the expected returns of a stock and calculate its abnormal
returns during an event window. This model is less complex than the Market
Model, as it assumes that a stock’s expected return is equal to the market
return, without considering any stock-specific factors. The Market
Adjusted Model is particularly useful in situations where the estimation of
individual stock parameters (such as alpha and beta) is not feasible or
desired, and a basic benchmark for comparison is needed.
\end{Description}
%
\begin{Section}{Super class}
\code{\LinkA{ModelBase}{ModelBase}} -> \code{MarketAdjustedModel}
\end{Section}
%
\begin{Section}{Public fields}

\begin{description}

\item[\code{model\_name}] Name of the model.

\end{description}


\end{Section}
%
\begin{Section}{Methods}
%
\begin{SubSection}{Public methods}
\begin{itemize}

\item{} \Rhref{\#method-MarketAdjustedModel-fit}{\code{MarketAdjustedModel\$fit()}}
\item{} \Rhref{\#method-MarketAdjustedModel-abnormal_returns}{\code{MarketAdjustedModel\$abnormal\_returns()}}
\item{} \Rhref{\#method-MarketAdjustedModel-clone}{\code{MarketAdjustedModel\$clone()}}

\end{itemize}

\end{SubSection}



\hypertarget{method-MarketAdjustedModel-fit}{}
%
\begin{SubSection}{\code{MarketAdjustedModel\$fit()}}
fit Fit the model with given data.
%
\begin{SubSubSection}{Usage}

\begin{alltt}MarketAdjustedModel$fit(data_tbl)\end{alltt}



\end{SubSubSection}

%
\begin{SubSubSection}{Arguments}

\begin{description}

\item[\code{data\_tbl}] Data frame or tibble containing the data to fit.

\end{description}



\end{SubSubSection}

\end{SubSection}




\hypertarget{method-MarketAdjustedModel-abnormal_returns}{}
%
\begin{SubSection}{\code{MarketAdjustedModel\$abnormal\_returns()}}
abnormal\_returns Calculate the abnormal returns with given data.
%
\begin{SubSubSection}{Usage}

\begin{alltt}MarketAdjustedModel$abnormal_returns(data_tbl)\end{alltt}



\end{SubSubSection}

%
\begin{SubSubSection}{Arguments}

\begin{description}

\item[\code{data\_tbl}] Data frame or tibble containing the data to calculate abnormal returns.

\end{description}



\end{SubSubSection}

\end{SubSection}




\hypertarget{method-MarketAdjustedModel-clone}{}
%
\begin{SubSection}{\code{MarketAdjustedModel\$clone()}}
The objects of this class are cloneable with this method.
%
\begin{SubSubSection}{Usage}

\begin{alltt}MarketAdjustedModel$clone(deep = FALSE)\end{alltt}



\end{SubSubSection}

%
\begin{SubSubSection}{Arguments}

\begin{description}

\item[\code{deep}] Whether to make a deep clone.

\end{description}



\end{SubSubSection}

\end{SubSection}


\end{Section}
\HeaderA{MarketModel}{Market Model}{MarketModel}
%
\begin{Description}
The Market Model is a widely used method in event studies to estimate the
expected returns of a stock and calculate its abnormal returns during an
event window. The model is based on a simple linear regression framework and
captures the relationship between a stock’s return and the return of a market
index, such as the S\&P 500 or the Dow Jones Industrial Average. The
underlying assumption of the Market Model is that a stock’s return is
primarily influenced by market movements, along with a stock-specific
idiosyncratic component.
\end{Description}
%
\begin{Section}{Super class}
\code{\LinkA{ModelBase}{ModelBase}} -> \code{MarketModel}
\end{Section}
%
\begin{Section}{Public fields}

\begin{description}

\item[\code{model\_name}] Name of the model.

\item[\code{formula}] The formula applied for calculating the market model

\item[\code{use\_hac}] Logical. Use HAC (Newey-West) standard errors.

\item[\code{hac\_lag}] Integer or NULL. Lag truncation for Newey-West.
NULL uses the automatic bandwidth selection.

\end{description}


\end{Section}
%
\begin{Section}{Methods}
%
\begin{SubSection}{Public methods}
\begin{itemize}

\item{} \Rhref{\#method-MarketModel-initialize}{\code{MarketModel\$new()}}
\item{} \Rhref{\#method-MarketModel-set_formula}{\code{MarketModel\$set\_formula()}}
\item{} \Rhref{\#method-MarketModel-fit}{\code{MarketModel\$fit()}}
\item{} \Rhref{\#method-MarketModel-abnormal_returns}{\code{MarketModel\$abnormal\_returns()}}
\item{} \Rhref{\#method-MarketModel-clone}{\code{MarketModel\$clone()}}

\end{itemize}

\end{SubSection}



\hypertarget{method-MarketModel-initialize}{}
%
\begin{SubSection}{\code{MarketModel\$new()}}
Create a new MarketModel.
%
\begin{SubSubSection}{Usage}

\begin{alltt}MarketModel$new(use_hac = FALSE, hac_lag = NULL)\end{alltt}



\end{SubSubSection}

%
\begin{SubSubSection}{Arguments}

\begin{description}

\item[\code{use\_hac}] Logical. Use HAC (Newey-West) standard errors.
Requires the \pkg{sandwich} package.
\item[\code{hac\_lag}] Integer or NULL. Lag truncation for Newey-West.

\end{description}



\end{SubSubSection}

\end{SubSection}




\hypertarget{method-MarketModel-set_formula}{}
%
\begin{SubSection}{\code{MarketModel\$set\_formula()}}
Set the formula
%
\begin{SubSubSection}{Usage}

\begin{alltt}MarketModel$set_formula(formula)\end{alltt}



\end{SubSubSection}

%
\begin{SubSubSection}{Arguments}

\begin{description}

\item[\code{formula}] A formula.

\end{description}



\end{SubSubSection}

\end{SubSection}




\hypertarget{method-MarketModel-fit}{}
%
\begin{SubSection}{\code{MarketModel\$fit()}}
Fit the model with given data.
%
\begin{SubSubSection}{Usage}

\begin{alltt}MarketModel$fit(data_tbl)\end{alltt}



\end{SubSubSection}

%
\begin{SubSubSection}{Arguments}

\begin{description}

\item[\code{data\_tbl}] Data frame or tibble containing the data to fit.

\end{description}



\end{SubSubSection}

\end{SubSection}




\hypertarget{method-MarketModel-abnormal_returns}{}
%
\begin{SubSection}{\code{MarketModel\$abnormal\_returns()}}
Calculate the abnormal returns with given data.
%
\begin{SubSubSection}{Usage}

\begin{alltt}MarketModel$abnormal_returns(data_tbl)\end{alltt}



\end{SubSubSection}

%
\begin{SubSubSection}{Arguments}

\begin{description}

\item[\code{data\_tbl}] Data frame or tibble containing the data to calculate abnormal returns.

\end{description}



\end{SubSubSection}

\end{SubSection}




\hypertarget{method-MarketModel-clone}{}
%
\begin{SubSection}{\code{MarketModel\$clone()}}
The objects of this class are cloneable with this method.
%
\begin{SubSubSection}{Usage}

\begin{alltt}MarketModel$clone(deep = FALSE)\end{alltt}



\end{SubSubSection}

%
\begin{SubSubSection}{Arguments}

\begin{description}

\item[\code{deep}] Whether to make a deep clone.

\end{description}



\end{SubSubSection}

\end{SubSection}


\end{Section}
\HeaderA{ModelBase}{ModelBase}{ModelBase}
%
\begin{Description}
Base model class for event study. Each single event study will
get its own model initialization and fitting. Therefore, the input DataFrame
contains the data for a single Event Study. For custom models, except the
child modles of the market model, several statistics must be included, namely
sigma, degree\_of\_freedom, first\_order\_auto\_correlation, residuals,
forecast\_error\_corrected\_sigma, and forecast\_error\_corrected\_sigma\_car. Part
of these statistics are necessary for calculating the Event Study statistics.
\end{Description}
%
\begin{Section}{Public fields}

\begin{description}

\item[\code{model\_name}] Name of the model.

\item[\code{degenerate\_mode}] Resolved degenerate-input mode injected by
fit\_model() before fit() is called. Subclasses that implement the
degenerate-input contract (MarketModel and Phase-2 models) read this
field inside fit() via .resolve\_degenerate\_mode(self\$degenerate\_mode).

\item[\code{event\_id}] Event identifier threaded from the outer data\_tbl row.
Used in degenerate-input error/warning messages.

\item[\code{firm\_symbol}] Firm identifier threaded from the outer data\_tbl row.
Used in degenerate-input error/warning messages.

\end{description}


\end{Section}
%
\begin{Section}{Active bindings}

\begin{description}

\item[\code{statistics}] Read-only field to get statistics.

\item[\code{model}] Read-only field to get the fitted model.

\item[\code{is\_fitted}] Read-only field to check if the model is fitted.
Statistics object contains different model specific KPIs
that describes the fitted model.

\end{description}


\end{Section}
%
\begin{Section}{Methods}
%
\begin{SubSection}{Public methods}
\begin{itemize}

\item{} \Rhref{\#method-ModelBase-fit}{\code{ModelBase\$fit()}}
\item{} \Rhref{\#method-ModelBase-abnormal_returns}{\code{ModelBase\$abnormal\_returns()}}
\item{} \Rhref{\#method-ModelBase-clone}{\code{ModelBase\$clone()}}

\end{itemize}

\end{SubSection}



\hypertarget{method-ModelBase-fit}{}
%
\begin{SubSection}{\code{ModelBase\$fit()}}
Fits the model with given data.
%
\begin{SubSubSection}{Usage}

\begin{alltt}ModelBase$fit(data_tbl)\end{alltt}



\end{SubSubSection}

%
\begin{SubSubSection}{Arguments}

\begin{description}

\item[\code{data\_tbl}] A data frame or tibble containing the data to fit.

\end{description}



\end{SubSubSection}

\end{SubSection}




\hypertarget{method-ModelBase-abnormal_returns}{}
%
\begin{SubSection}{\code{ModelBase\$abnormal\_returns()}}
Calculate the abnormal returns with given data and fitted model.
%
\begin{SubSubSection}{Usage}

\begin{alltt}ModelBase$abnormal_returns(data_tbl)\end{alltt}



\end{SubSubSection}

%
\begin{SubSubSection}{Arguments}

\begin{description}

\item[\code{data\_tbl}] Data frame or tibble containing the data to calculate abnormal returns.

\end{description}



\end{SubSubSection}

\end{SubSection}




\hypertarget{method-ModelBase-clone}{}
%
\begin{SubSection}{\code{ModelBase\$clone()}}
The objects of this class are cloneable with this method.
%
\begin{SubSubSection}{Usage}

\begin{alltt}ModelBase$clone(deep = FALSE)\end{alltt}



\end{SubSubSection}

%
\begin{SubSubSection}{Arguments}

\begin{description}

\item[\code{deep}] Whether to make a deep clone.

\end{description}



\end{SubSubSection}

\end{SubSection}


\end{Section}
\HeaderA{model\_diagnostics}{Model Diagnostics for Event Study}{model.Rul.diagnostics}
%
\begin{Description}
Run diagnostic tests on fitted event study models. Tests include
normality of residuals, autocorrelation, and basic model fit quality.
\end{Description}
%
\begin{Usage}
\begin{verbatim}
model_diagnostics(task, event_id = NULL)
\end{verbatim}
\end{Usage}
%
\begin{Arguments}
\begin{ldescription}
\item[\code{task}] A fitted EventStudyTask (after fit\_model has been called).

\item[\code{event\_id}] Optional event identifier. If NULL, runs diagnostics
for all events.
\end{ldescription}
\end{Arguments}
%
\begin{Value}
A tibble with diagnostic test results for each event.
\end{Value}
\HeaderA{MultiEventStatisticsSet}{Multi Event Statistic Set}{MultiEventStatisticsSet}
%
\begin{Description}
Contains a set of multi event test statistics as, e.g., AAR and CAAR
cross-sectional t test. These are the default test statistics.
\end{Description}
%
\begin{Section}{Super class}
\code{\LinkA{StatisticsSetBase}{StatisticsSetBase}} -> \code{MultiEventStatisticsSet}
\end{Section}
%
\begin{Section}{Public fields}

\begin{description}

\item[\code{tests}] Container with the statistical tests.

\end{description}


\end{Section}
%
\begin{Section}{Methods}
%
\begin{SubSection}{Public methods}
\begin{itemize}

\item{} \Rhref{\#method-MultiEventStatisticsSet-clone}{\code{MultiEventStatisticsSet\$clone()}}

\end{itemize}

\end{SubSection}




\hypertarget{method-MultiEventStatisticsSet-clone}{}
%
\begin{SubSection}{\code{MultiEventStatisticsSet\$clone()}}
The objects of this class are cloneable with this method.
%
\begin{SubSubSection}{Usage}

\begin{alltt}MultiEventStatisticsSet$clone(deep = FALSE)\end{alltt}



\end{SubSubSection}

%
\begin{SubSubSection}{Arguments}

\begin{description}

\item[\code{deep}] Whether to make a deep clone.

\end{description}



\end{SubSubSection}

\end{SubSection}


\end{Section}
\HeaderA{nonparametric\_intraday\_test}{Non-Parametric Intraday Event Study Test}{nonparametric.Rul.intraday.Rul.test}
%
\begin{Description}
Implements the non-parametric intraday event study methodology of
Rinaudo \& Saha (2014). For each event time, cumulative abnormal returns
(CARs) on the event day are compared to the empirical distribution of CARs
from estimation-period days. Significance is assessed without distributional
assumptions.
\end{Description}
%
\begin{Usage}
\begin{verbatim}
nonparametric_intraday_test(
  estimation_window,
  event_window,
  event_times,
  p = 0.05,
  init_window = 5L,
  upper = FALSE
)
\end{verbatim}
\end{Usage}
%
\begin{Arguments}
\begin{ldescription}
\item[\code{estimation\_window}] Data frame with columns \code{day}, \code{time},
and \code{abnormalReturn} containing intraday abnormal returns for
multiple estimation days.

\item[\code{event\_window}] Data frame with columns \code{time} and
\code{abnormalReturn} for the single event day.

\item[\code{event\_times}] Character vector of intraday times at which events
occurred (e.g., \code{c("10:00", "14:30")}).

\item[\code{p}] Numeric; significance level in (0, 1). Default 0.05.

\item[\code{init\_window}] Integer; number of initial observations to accumulate
before testing significance. Default 5.

\item[\code{upper}] Logical; if \code{TRUE} test for positive abnormal returns
(upper tail), otherwise test for negative abnormal returns (lower tail).
Default \code{FALSE}.
\end{ldescription}
\end{Arguments}
%
\begin{Value}
A named list of tibbles, one per event time. Each tibble contains:
\begin{description}

\item[id] Observation index within the significant window.
\item[CAR] Cumulative abnormal return on the event day.
\item[CI] Empirical confidence interval boundary from estimation days.
\item[fitCI] Polynomial-smoothed CI (degree-4 polynomial fit).

\end{description}

When an event time is not significant, returns a single-row tibble with
all values equal to zero.
\end{Value}
%
\begin{References}
Rinaudo, J.B. \& Saha, A. (2014). Non-parametric intraday event studies.
\end{References}
%
\begin{Examples}
\begin{ExampleCode}
## Not run: 
results <- nonparametric_intraday_test(
  estimation_window = est_data,
  event_window = event_data,
  event_times = c("10:00", "14:30"),
  p = 0.05,
  init_window = 5
)

## End(Not run)

\end{ExampleCode}
\end{Examples}
\HeaderA{OpenAICompatProvider}{OpenAICompatProvider}{OpenAICompatProvider}
%
\begin{Description}
A provider that issues \code{POST \{base\_url\}/chat/completions}
with an OpenAI-shaped request body. Because the endpoint is selected purely
by \code{base\_url} + \code{model}, this single class covers OpenAI itself
AND every OpenAI-compatible backend — Ollama, LM Studio, or any gateway —
via a \code{base\_url} override, with no code change.

The API key is resolved at CALL time from \code{OPENAI\_API\_KEY} (never at
construction, never stored on the object) and attached with
\code{req\_auth\_bearer\_token()}, which redacts the \code{Authorization}
header in every print/error path. A missing key, transport/timeout failure,
non-2xx status, malformed body, or missing completion text each degrades to
exactly one \code{warning()} plus an `es\_provider\_response` with
\code{text = NA\_character\_} — it never crashes the session and never returns
a fabricated completion.

\code{httr2} is required only for this provider and is guarded by
\code{requireNamespace()} at the top of \code{complete()}, so the package
installs and \code{R CMD check}s cleanly with \code{httr2} absent.
\end{Description}
%
\begin{Section}{Super class}
\code{\LinkA{ProviderBase}{ProviderBase}} -> \code{OpenAICompatProvider}
\end{Section}
%
\begin{Section}{Methods}
%
\begin{SubSection}{Public methods}
\begin{itemize}

\item{} \Rhref{\#method-OpenAICompatProvider-initialize}{\code{OpenAICompatProvider\$new()}}
\item{} \Rhref{\#method-OpenAICompatProvider-complete}{\code{OpenAICompatProvider\$complete()}}
\item{} \Rhref{\#method-OpenAICompatProvider-clone}{\code{OpenAICompatProvider\$clone()}}

\end{itemize}

\end{SubSection}



\hypertarget{method-OpenAICompatProvider-initialize}{}
%
\begin{SubSection}{\code{OpenAICompatProvider\$new()}}
Construct an OpenAICompatProvider. Stores non-secret config
only; the API key is resolved at call time inside \code{complete()}.
%
\begin{SubSubSection}{Usage}

\begin{alltt}OpenAICompatProvider$new(model, base_url = "https://api.openai.com/v1", ...)\end{alltt}



\end{SubSubSection}

%
\begin{SubSubSection}{Arguments}

\begin{description}

\item[\code{model}] Character model identifier, e.g. "gpt-4o" or a local model
name like "llama3".
\item[\code{base\_url}] Character base URL. Defaults to the OpenAI cloud endpoint;
override for OpenAI-compatible servers (Ollama, LM Studio, gateways).
\item[\code{...}] Ignored; accepted for forward-compatibility.

\end{description}



\end{SubSubSection}

\end{SubSection}




\hypertarget{method-OpenAICompatProvider-complete}{}
%
\begin{SubSection}{\code{OpenAICompatProvider\$complete()}}
Complete a prompt via \code{POST \{base\_url\}/chat/completions}.
Resolves the key at call time; on a missing key or any HTTP/parse failure
returns one warning + \code{NA}. Structured output is OPTIONAL: when
\code{schema} is supplied a \code{response\_format} json\_schema is added to
the body, but omitting it uses the plain-text path (which local models
that lack \code{response\_format} support still handle).
%
\begin{SubSubSection}{Usage}

\begin{alltt}OpenAICompatProvider$complete(prompt, schema = NULL, ...)\end{alltt}



\end{SubSubSection}

%
\begin{SubSubSection}{Arguments}

\begin{description}

\item[\code{prompt}] Character prompt.
\item[\code{schema}] Optional structured-output schema (list) or \code{NULL}.
\item[\code{...}] Ignored; accepted for forward-compatibility.

\end{description}



\end{SubSubSection}

%
\begin{SubSubSection}{Returns}
An `es\_provider\_response` (see \code{ProviderBase\$complete}).

\end{SubSubSection}

\end{SubSection}




\hypertarget{method-OpenAICompatProvider-clone}{}
%
\begin{SubSection}{\code{OpenAICompatProvider\$clone()}}
The objects of this class are cloneable with this method.
%
\begin{SubSubSection}{Usage}

\begin{alltt}OpenAICompatProvider$clone(deep = FALSE)\end{alltt}



\end{SubSubSection}

%
\begin{SubSubSection}{Arguments}

\begin{description}

\item[\code{deep}] Whether to make a deep clone.

\end{description}



\end{SubSubSection}

\end{SubSection}


\end{Section}
%
\begin{Examples}
\begin{ExampleCode}
## Not run: 
# OpenAI (reads OPENAI_API_KEY from the environment at call time):
p <- OpenAICompatProvider$new(model = "gpt-4o")
p$complete("Summarise these event-study diagnostics")

# Any OpenAI-compatible endpoint works via a base_url override — e.g. a local
# Ollama server (no cloud key needed by the server, but OPENAI_API_KEY is still
# read as the bearer token; set it to any non-empty value for local servers):
p_local <- OpenAICompatProvider$new(
  model = "llama3",
  base_url = "http://localhost:11434/v1"   # Ollama; LM Studio: :1234/v1
)
p_local$complete("Summarise these diagnostics")

## End(Not run)
\end{ExampleCode}
\end{Examples}
\HeaderA{PanelEventStudyTask}{Panel Event Study Task}{PanelEventStudyTask}
%
\begin{Description}
Task container for panel (difference-in-differences style) event studies.
This is designed for settings where units receive treatment at potentially
different times, and we observe outcomes over multiple periods.
\end{Description}
%
\begin{Section}{Methods}

\begin{description}

\item[\code{new(panel\_data, unit\_id, time\_id, outcome, treatment, treatment\_time)}] 
Create a new panel event study task.
\item[\code{print()}] Print summary.

\end{description}

\end{Section}
%
\begin{Section}{Public fields}

\begin{description}

\item[\code{panel\_data}] The panel dataset in long format.

\item[\code{unit\_id}] Name of the unit identifier column.

\item[\code{time\_id}] Name of the time identifier column.

\item[\code{outcome}] Name of the outcome column.

\item[\code{treatment}] Name of the treatment indicator column (0/1).

\item[\code{treatment\_time}] Name of the treatment timing column
(period when unit first treated; NA for never-treated).

\item[\code{results}] Estimation results (populated after estimation).

\end{description}


\end{Section}
%
\begin{Section}{Methods}
%
\begin{SubSection}{Public methods}
\begin{itemize}

\item{} \Rhref{\#method-PanelEventStudyTask-initialize}{\code{PanelEventStudyTask\$new()}}
\item{} \Rhref{\#method-PanelEventStudyTask-print}{\code{PanelEventStudyTask\$print()}}
\item{} \Rhref{\#method-PanelEventStudyTask-clone}{\code{PanelEventStudyTask\$clone()}}

\end{itemize}

\end{SubSection}



\hypertarget{method-PanelEventStudyTask-initialize}{}
%
\begin{SubSection}{\code{PanelEventStudyTask\$new()}}
Create a new PanelEventStudyTask.
%
\begin{SubSubSection}{Usage}

\begin{alltt}PanelEventStudyTask$new(
  panel_data,
  unit_id = "unit_id",
  time_id = "time_id",
  outcome = "outcome",
  treatment = "treated",
  treatment_time = "treatment_time"
)\end{alltt}



\end{SubSubSection}

%
\begin{SubSubSection}{Arguments}

\begin{description}

\item[\code{panel\_data}] Long-format panel data frame.
\item[\code{unit\_id}] Name of the unit ID column.
\item[\code{time\_id}] Name of the time period column.
\item[\code{outcome}] Name of the outcome variable column.
\item[\code{treatment}] Name of the treatment indicator column (0/1).
\item[\code{treatment\_time}] Name of column indicating when each unit is
first treated (NA for never-treated controls).

\end{description}



\end{SubSubSection}

\end{SubSection}




\hypertarget{method-PanelEventStudyTask-print}{}
%
\begin{SubSection}{\code{PanelEventStudyTask\$print()}}
Print summary.
%
\begin{SubSubSection}{Usage}

\begin{alltt}PanelEventStudyTask$print(...)\end{alltt}



\end{SubSubSection}

\end{SubSection}




\hypertarget{method-PanelEventStudyTask-clone}{}
%
\begin{SubSection}{\code{PanelEventStudyTask\$clone()}}
The objects of this class are cloneable with this method.
%
\begin{SubSubSection}{Usage}

\begin{alltt}PanelEventStudyTask$clone(deep = FALSE)\end{alltt}



\end{SubSubSection}

%
\begin{SubSubSection}{Arguments}

\begin{description}

\item[\code{deep}] Whether to make a deep clone.

\end{description}



\end{SubSubSection}

\end{SubSection}


\end{Section}
\HeaderA{ParameterSet}{Event Study Parameter Set}{ParameterSet}
%
\begin{Description}
The parameter set defines the Event Study, e.g, the return
calculation, the event study model, the AR, CAR, AAR, and CAAR test
statistics that should be applied.
\end{Description}
%
\begin{Section}{Public fields}

\begin{description}

\item[\code{return\_calculation}] A R6 object for calculating returns.

\item[\code{return\_model}] A R6 object for fitting the desired model.

\item[\code{single\_event\_statistics}] single event test statistic R6 object.

\item[\code{multi\_event\_statistics}] multi event test statistic R6 object.

\item[\code{study\_type}] Type of event study: "return" (default),
"volume", or "volatility". Affects axis labels in plots.

\item[\code{degenerate\_handling}] Controls how degenerate estimation data
is handled. One of \code{"lenient"} (default), \code{"strict"},
or \code{NULL} (defer to package option or built-in default).
See \code{?degenerate-input-contract} for details.

\end{description}


\end{Section}
%
\begin{Section}{Methods}
%
\begin{SubSection}{Public methods}
\begin{itemize}

\item{} \Rhref{\#method-ParameterSet-initialize}{\code{ParameterSet\$new()}}
\item{} \Rhref{\#method-ParameterSet-print}{\code{ParameterSet\$print()}}
\item{} \Rhref{\#method-ParameterSet-clone}{\code{ParameterSet\$clone()}}

\end{itemize}

\end{SubSection}



\hypertarget{method-ParameterSet-initialize}{}
%
\begin{SubSection}{\code{ParameterSet\$new()}}
Initialize the parameters that defines the Event Study that should be applied.
%
\begin{SubSubSection}{Usage}

\begin{alltt}ParameterSet$new(
  return_calculation = SimpleReturn$new(),
  return_model = MarketModel$new(),
  single_event_statistics = SingleEventStatisticsSet$new(),
  multi_event_statistics = MultiEventStatisticsSet$new(),
  degenerate_handling = NULL
)\end{alltt}



\end{SubSubSection}

%
\begin{SubSubSection}{Arguments}

\begin{description}

\item[\code{return\_calculation}] An initialized return calculation class.
Defaults to SimpleReturn.
\item[\code{return\_model}] An initialized event study model.
Defaults to MarketModel.
\item[\code{single\_event\_statistics}] Definition of single event test statistics.
Defaults to SingleEventStatisticsSet (AR T and CAR T tests).
\item[\code{multi\_event\_statistics}] Definition of multiple event test statistics.
Defaults to MultiEventStatisticsSet (CSect T test).
\item[\code{degenerate\_handling}] One of \code{"lenient"}, \code{"strict"},
or \code{NULL}. See \code{?degenerate-input-contract}.

\end{description}



\end{SubSubSection}

\end{SubSection}




\hypertarget{method-ParameterSet-print}{}
%
\begin{SubSection}{\code{ParameterSet\$print()}}
Print a summary of the parameter set.
%
\begin{SubSubSection}{Usage}

\begin{alltt}ParameterSet$print(...)\end{alltt}



\end{SubSubSection}

\end{SubSection}




\hypertarget{method-ParameterSet-clone}{}
%
\begin{SubSection}{\code{ParameterSet\$clone()}}
The objects of this class are cloneable with this method.
%
\begin{SubSubSection}{Usage}

\begin{alltt}ParameterSet$clone(deep = FALSE)\end{alltt}



\end{SubSubSection}

%
\begin{SubSubSection}{Arguments}

\begin{description}

\item[\code{deep}] Whether to make a deep clone.

\end{description}



\end{SubSubSection}

\end{SubSection}


\end{Section}
\HeaderA{PatellZTest}{Patell or Standardized Residual Test (PatellZTest)}{PatellZTest}
%
\begin{Description}
The Patell or Standardized Residual Test is a statistical tool employed in event
studies to evaluate the null hypothesis that the average abnormal return at
the event date is zero. The test statistics is approximately distributed as
N(0, 1).

See also \url{https://eventstudy.de/statistics/aar_caar_statistics.html}
\end{Description}
%
\begin{Section}{Super class}
\code{\LinkA{TestStatisticBase}{TestStatisticBase}} -> \code{PatellZTest}
\end{Section}
%
\begin{Section}{Public fields}

\begin{description}

\item[\code{name}] Short code of the test statistic.

\end{description}


\end{Section}
%
\begin{Section}{Methods}
%
\begin{SubSection}{Public methods}
\begin{itemize}

\item{} \Rhref{\#method-PatellZTest-compute}{\code{PatellZTest\$compute()}}
\item{} \Rhref{\#method-PatellZTest-clone}{\code{PatellZTest\$clone()}}

\end{itemize}

\end{SubSection}




\hypertarget{method-PatellZTest-compute}{}
%
\begin{SubSection}{\code{PatellZTest\$compute()}}
Computes the Patell Z test statistics for multiple events.
%
\begin{SubSubSection}{Usage}

\begin{alltt}PatellZTest$compute(data_tbl, model)\end{alltt}



\end{SubSubSection}

%
\begin{SubSubSection}{Arguments}

\begin{description}

\item[\code{data\_tbl}] The data for a multiple event with
calculated abnormal returns.
\item[\code{model}] The fitted model.

\end{description}



\end{SubSubSection}

\end{SubSection}




\hypertarget{method-PatellZTest-clone}{}
%
\begin{SubSection}{\code{PatellZTest\$clone()}}
The objects of this class are cloneable with this method.
%
\begin{SubSubSection}{Usage}

\begin{alltt}PatellZTest$clone(deep = FALSE)\end{alltt}



\end{SubSubSection}

%
\begin{SubSubSection}{Arguments}

\begin{description}

\item[\code{deep}] Whether to make a deep clone.

\end{description}



\end{SubSubSection}

\end{SubSection}


\end{Section}
\HeaderA{plot\_car\_distribution}{Plot CAR Distribution}{plot.Rul.car.Rul.distribution}
%
\begin{Description}
Create a histogram of CARs across events, optionally colored by group.
\end{Description}
%
\begin{Usage}
\begin{verbatim}
plot_car_distribution(
  task,
  car_window = NULL,
  bins = 30,
  by_group = FALSE,
  title = NULL
)
\end{verbatim}
\end{Usage}
%
\begin{Arguments}
\begin{ldescription}
\item[\code{task}] A fitted EventStudyTask.

\item[\code{car\_window}] Optional two-element vector for CAR window.

\item[\code{bins}] Number of histogram bins. Default 30.

\item[\code{by\_group}] Logical. If TRUE, color histogram by group.

\item[\code{title}] Optional plot title.
\end{ldescription}
\end{Arguments}
%
\begin{Value}
A ggplot2 object.
\end{Value}
\HeaderA{plot\_diagnostics}{Plot Model Diagnostics}{plot.Rul.diagnostics}
%
\begin{Description}
Create diagnostic plots for fitted event study models, including
residual plots and Q-Q plots.
\end{Description}
%
\begin{Usage}
\begin{verbatim}
plot_diagnostics(task, event_id = NULL)
\end{verbatim}
\end{Usage}
%
\begin{Arguments}
\begin{ldescription}
\item[\code{task}] A fitted EventStudyTask.

\item[\code{event\_id}] The event identifier to plot diagnostics for.
\end{ldescription}
\end{Arguments}
%
\begin{Value}
A ggplot2 plot arranged with patchwork-style layout.
\end{Value}
\HeaderA{plot\_event\_study}{Plot Event Study Results}{plot.Rul.event.Rul.study}
%
\begin{Description}
Create an event study plot showing abnormal returns (AR), cumulative
abnormal returns (CAR), average abnormal returns (AAR), or cumulative
average abnormal returns (CAAR) with confidence bands.
\end{Description}
%
\begin{Usage}
\begin{verbatim}
plot_event_study(
  task,
  type = "car",
  event_id = NULL,
  group = NULL,
  stat_name = "CSectT",
  confidence_level = 0.95,
  title = NULL
)
\end{verbatim}
\end{Usage}
%
\begin{Arguments}
\begin{ldescription}
\item[\code{task}] A fitted EventStudyTask with statistics computed.

\item[\code{type}] Type of plot: "ar", "car", "aar", or "caar".

\item[\code{event\_id}] Event identifier for single-event plots ("ar", "car").

\item[\code{group}] Group name for multi-event plots ("aar", "caar").

\item[\code{stat\_name}] Name of the multi-event test statistic to use
for AAR/CAAR plots. Defaults to "CSectT".

\item[\code{confidence\_level}] Confidence level for the bands. Default is 0.95.

\item[\code{title}] Optional plot title.
\end{ldescription}
\end{Arguments}
%
\begin{Value}
A ggplot2 plot object.
\end{Value}
\HeaderA{plot\_panel\_event\_study}{Plot Panel Event Study Results}{plot.Rul.panel.Rul.event.Rul.study}
%
\begin{Description}
Create an event study plot from panel estimation results. Can also be
called via \code{plot\_event\_study()} when given a PanelEventStudyTask.
\end{Description}
%
\begin{Usage}
\begin{verbatim}
plot_panel_event_study(task, confidence_level = 0.95, title = NULL)
\end{verbatim}
\end{Usage}
%
\begin{Arguments}
\begin{ldescription}
\item[\code{task}] A PanelEventStudyTask with results.

\item[\code{confidence\_level}] Confidence level for error bars. Default 0.95.

\item[\code{title}] Optional plot title.
\end{ldescription}
\end{Arguments}
%
\begin{Value}
A ggplot2 object.
\end{Value}
\HeaderA{plot\_stocks}{Plot Stocks}{plot.Rul.stocks}
%
\begin{Description}
Visualize the adjusted close prices or abnormal returns of selected stocks
around a specified event date.
\end{Description}
%
\begin{Usage}
\begin{verbatim}
plot_stocks(
  task,
  target_variable = "firm_adjusted",
  add_event_date = FALSE,
  max_symbols = 6,
  do_sample = TRUE
)
\end{verbatim}
\end{Usage}
%
\begin{Arguments}
\begin{ldescription}
\item[\code{task}] An EventStudyTask object.

\item[\code{target\_variable}] A character string specifying the target variable
to plot. Default is "firm\_adjusted".

\item[\code{add\_event\_date}] Add vertical line for event date (TRUE/FALSE).

\item[\code{max\_symbols}] An integer specifying the maximum number of symbols
to display. Default is 6.

\item[\code{do\_sample}] A boolean specifying whether to randomly sample the
symbols if the number exceeds max\_symbols. Default is TRUE.
\end{ldescription}
\end{Arguments}
%
\begin{Value}
A plotly plot object.
\end{Value}
\HeaderA{plot\_synthetic\_control}{Plot Synthetic Control Results}{plot.Rul.synthetic.Rul.control}
%
\begin{Description}
Create plots for synthetic control analysis: trajectory comparison,
treatment effect gap, or placebo test results.
\end{Description}
%
\begin{Usage}
\begin{verbatim}
plot_synthetic_control(task, type = c("trajectory", "gap", "placebo"))
\end{verbatim}
\end{Usage}
%
\begin{Arguments}
\begin{ldescription}
\item[\code{task}] A \code{SyntheticControlTask} with results.

\item[\code{type}] Plot type: \code{"trajectory"} (default), \code{"gap"}, or
\code{"placebo"}.
\end{ldescription}
\end{Arguments}
%
\begin{Value}
A ggplot2 object.
\end{Value}
\HeaderA{prepare\_event\_study}{Prepare data for an Event Study}{prepare.Rul.event.Rul.study}
%
\begin{Description}
Perform return calculation for each stock and the corresponding reference
market defined in the task and the parameter set.
\end{Description}
%
\begin{Usage}
\begin{verbatim}
prepare_event_study(task, parameter_set)
\end{verbatim}
\end{Usage}
%
\begin{Arguments}
\begin{ldescription}
\item[\code{task}] An Event Study task.

\item[\code{parameter\_set}] A parameter set that defines the Event Study.
\end{ldescription}
\end{Arguments}
%
\begin{Value}
The task object with returns and windows appended.
\end{Value}
\HeaderA{prepare\_intraday\_event\_study}{Prepare Intraday Event Study}{prepare.Rul.intraday.Rul.event.Rul.study}
%
\begin{Description}
Compute returns and assign estimation/event windows for intraday data.
Windows are specified in number of observations (bars) rather than days.
\end{Description}
%
\begin{Usage}
\begin{verbatim}
prepare_intraday_event_study(task, parameter_set)
\end{verbatim}
\end{Usage}
%
\begin{Arguments}
\begin{ldescription}
\item[\code{task}] An IntradayEventStudyTask.

\item[\code{parameter\_set}] A ParameterSet.
\end{ldescription}
\end{Arguments}
%
\begin{Value}
The task with returns and windows appended.
\end{Value}
\HeaderA{pretrend\_test}{Pre-trend Test for Event Study}{pretrend.Rul.test}
%
\begin{Description}
Performs a joint F-test of pre-event abnormal returns to assess whether
there are statistically significant pre-event effects, which would
indicate potential model misspecification or confounding.
\end{Description}
%
\begin{Usage}
\begin{verbatim}
pretrend_test(task, group = NULL)
\end{verbatim}
\end{Usage}
%
\begin{Arguments}
\begin{ldescription}
\item[\code{task}] A fitted EventStudyTask with abnormal returns computed.

\item[\code{group}] Optional group to filter by.
\end{ldescription}
\end{Arguments}
%
\begin{Value}
A tibble with pre-trend test results for each group.
\end{Value}
\HeaderA{print.Advice}{Print method for Advice objects}{print.Advice}
%
\begin{Description}
Prints a structured summary of the grounded AI advice, including source,
task type, grounding guard status, interpretation, recommendations with
evidence, and caveats. Follows the package convention of
\code{print.es\_advice} (cat-based, invisible return).
\end{Description}
%
\begin{Usage}
\begin{verbatim}
## S3 method for class 'Advice'
print(x, ...)
\end{verbatim}
\end{Usage}
%
\begin{Arguments}
\begin{ldescription}
\item[\code{x}] An object of class \code{"Advice"}.

\item[\code{...}] Additional arguments (ignored).
\end{ldescription}
\end{Arguments}
%
\begin{Value}
Invisibly returns \code{x}.
\end{Value}
\HeaderA{print.es\_advice}{Print method for es\_advice objects}{print.es.Rul.advice}
%
\begin{Description}
Prints a structured summary of the offline advice, listing each matched rule
with its severity, citation key, and recommendation. Follows the package
convention of \code{print.es\_diagnostics} and \code{print.es\_simulation}
(cat-based, invisible return).
\end{Description}
%
\begin{Usage}
\begin{verbatim}
## S3 method for class 'es_advice'
print(x, ...)
\end{verbatim}
\end{Usage}
%
\begin{Arguments}
\begin{ldescription}
\item[\code{x}] An object of class \code{"es\_advice"}.

\item[\code{...}] Additional arguments (ignored).
\end{ldescription}
\end{Arguments}
%
\begin{Value}
Invisibly returns \code{x}.
\end{Value}
\HeaderA{print.es\_cross\_sectional}{Print Cross-Sectional Regression Results}{print.es.Rul.cross.Rul.sectional}
%
\begin{Description}
Print Cross-Sectional Regression Results
\end{Description}
%
\begin{Usage}
\begin{verbatim}
## S3 method for class 'es_cross_sectional'
print(x, ...)
\end{verbatim}
\end{Usage}
%
\begin{Arguments}
\begin{ldescription}
\item[\code{x}] An \code{es\_cross\_sectional} object.

\item[\code{...}] Additional arguments (unused).
\end{ldescription}
\end{Arguments}
\HeaderA{print.es\_diagnostics}{Print method for es\_diagnostics objects}{print.es.Rul.diagnostics}
%
\begin{Description}
Prints a concise summary of the diagnostics, following the package
convention of \code{print.es\_simulation} and \code{print.es\_cross\_sectional}.
\end{Description}
%
\begin{Usage}
\begin{verbatim}
## S3 method for class 'es_diagnostics'
print(x, ...)
\end{verbatim}
\end{Usage}
%
\begin{Arguments}
\begin{ldescription}
\item[\code{x}] An object of class \code{"es\_diagnostics"}.

\item[\code{...}] Additional arguments (ignored).
\end{ldescription}
\end{Arguments}
%
\begin{Value}
Invisibly returns \code{x}.
\end{Value}
\HeaderA{print.EventStudySummary}{Print method for EventStudySummary}{print.EventStudySummary}
%
\begin{Description}
Print method for EventStudySummary
\end{Description}
%
\begin{Usage}
\begin{verbatim}
## S3 method for class 'EventStudySummary'
print(x, ...)
\end{verbatim}
\end{Usage}
%
\begin{Arguments}
\begin{ldescription}
\item[\code{x}] An EventStudySummary object.

\item[\code{...}] Additional arguments (unused).
\end{ldescription}
\end{Arguments}
\HeaderA{provider}{Construct a Grounded AI Advisor provider}{provider}
%
\begin{Description}
Convenience factory that resolves the backend by 3-tier precedence
(explicit \code{type} argument -> \code{EVENTSTUDY\_ADVISOR\_PROVIDER} env
selector -> default \code{"custom"}) and constructs the matching provider
object. All three backends are wired: the in-process \code{"custom"} hook and
the HTTP \code{"openai"} (OpenAI-compatible) and \code{"anthropic"} (native
Anthropic Messages) providers.
\end{Description}
%
\begin{Usage}
\begin{verbatim}
provider(type = NULL, fn = NULL, model = NULL, base_url = NULL, ...)
\end{verbatim}
\end{Usage}
%
\begin{Arguments}
\begin{ldescription}
\item[\code{type}] Provider type, one of \code{"custom"}, \code{"openai"},
\code{"anthropic"}. When \code{NULL} (the default) it is resolved from the
\code{EVENTSTUDY\_ADVISOR\_PROVIDER} env var, falling back to \code{"custom"}.

\item[\code{fn}] For \code{type = "custom"}, the user function
\code{function(prompt, schema, ...)} returning a list or character. Required
for the custom branch.

\item[\code{model}] Optional character model identifier (resolved from
\code{EVENTSTUDY\_ADVISOR\_MODEL} when \code{NULL}).

\item[\code{base\_url}] Optional character base URL for HTTP providers (resolved from
\code{EVENTSTUDY\_ADVISOR\_BASE\_URL} when \code{NULL}).

\item[\code{...}] Additional arguments forwarded to the provider constructor.
\end{ldescription}
\end{Arguments}
%
\begin{Details}
The \code{"custom"} branch needs no network and neither \code{httr2} nor
\code{jsonlite}. Keys are never read here; the HTTP providers resolve them at
call time inside their own \code{complete()}.
\end{Details}
%
\begin{Value}
A \code{ProviderBase} subclass instance.
\end{Value}
%
\begin{Examples}
\begin{ExampleCode}
# Custom provider runs in-process, no network:
p <- provider("custom", fn = function(prompt, schema) "advice text")
p$complete("Summarise")$text

## Not run: 
# HTTP providers make live calls; never run in examples or on CRAN:
p <- provider("openai", model = "gpt-4o")
p$complete("Summarise these diagnostics")

a <- provider("anthropic", model = "claude-opus-4-8")
a$complete("Summarise these diagnostics")

## End(Not run)
\end{ExampleCode}
\end{Examples}
\HeaderA{ProviderBase}{ProviderBase}{ProviderBase}
%
\begin{Description}
Abstract base class for Grounded AI Advisor providers. Defines the
uniform `complete(prompt, schema, ...)` contract, mirroring the package's
`ModelBase` / `TestStatisticBase` strategy-pattern idiom. Concrete providers
(`CustomProvider`, and the HTTP `OpenAICompatProvider` / `AnthropicProvider`
delivered in later plans) inherit from this base and implement `complete()`.

`ProviderBase` itself is abstract: calling `complete()` on it errors. It
stores only non-secret configuration (`model`, `base\_url`); API keys are
NEVER read at construction — they are resolved at call time inside each
concrete provider's `complete()`.
\end{Description}
%
\begin{Section}{Public fields}

\begin{description}

\item[\code{model}] Character model identifier (or NULL for provider default).

\item[\code{base\_url}] Character base URL for HTTP providers (or NULL).

\end{description}


\end{Section}
%
\begin{Section}{Methods}
%
\begin{SubSection}{Public methods}
\begin{itemize}

\item{} \Rhref{\#method-ProviderBase-initialize}{\code{ProviderBase\$new()}}
\item{} \Rhref{\#method-ProviderBase-complete}{\code{ProviderBase\$complete()}}
\item{} \Rhref{\#method-ProviderBase-clone}{\code{ProviderBase\$clone()}}

\end{itemize}

\end{SubSection}



\hypertarget{method-ProviderBase-initialize}{}
%
\begin{SubSection}{\code{ProviderBase\$new()}}
Construct a provider. Stores non-secret config only; never
reads or stores API keys (keys are resolved at call time).
%
\begin{SubSubSection}{Usage}

\begin{alltt}ProviderBase$new(model = NULL, base_url = NULL, ...)\end{alltt}



\end{SubSubSection}

%
\begin{SubSubSection}{Arguments}

\begin{description}

\item[\code{model}] Optional character model identifier.
\item[\code{base\_url}] Optional character base URL (HTTP providers only).
\item[\code{...}] Ignored; accepted for subclass forward-compatibility.

\end{description}



\end{SubSubSection}

\end{SubSection}




\hypertarget{method-ProviderBase-complete}{}
%
\begin{SubSection}{\code{ProviderBase\$complete()}}
Complete a prompt. Abstract on the base class — concrete
providers override this.
%
\begin{SubSubSection}{Usage}

\begin{alltt}ProviderBase$complete(prompt, schema = NULL, ...)\end{alltt}



\end{SubSubSection}

%
\begin{SubSubSection}{Arguments}

\begin{description}

\item[\code{prompt}] Character prompt to send to the provider.
\item[\code{schema}] Optional structured-output schema (list) or NULL.
\item[\code{...}] Provider-specific arguments.

\end{description}



\end{SubSubSection}

%
\begin{SubSubSection}{Returns}
An `es\_provider\_response` S3 list with fields:
\code{source} (character provider label),
\code{is\_deterministic} (always \code{FALSE} for provider output),
\code{text} (character completion, or \code{NA\_character\_} on failure),
\code{error} (\code{NULL} on success, character reason on failure).
These field names match the Phase 5 \code{es\_advice} shape so the Phase 7
wrapper slots in trivially.

\end{SubSubSection}

\end{SubSection}




\hypertarget{method-ProviderBase-clone}{}
%
\begin{SubSection}{\code{ProviderBase\$clone()}}
The objects of this class are cloneable with this method.
%
\begin{SubSubSection}{Usage}

\begin{alltt}ProviderBase$clone(deep = FALSE)\end{alltt}



\end{SubSubSection}

%
\begin{SubSubSection}{Arguments}

\begin{description}

\item[\code{deep}] Whether to make a deep clone.

\end{description}



\end{SubSubSection}

\end{SubSection}


\end{Section}
\HeaderA{RankTest}{Rank Test (Corrado 1989)}{RankTest}
%
\begin{Description}
Non-parametric rank test for event studies. Ranks abnormal returns
across the combined estimation and event windows, which is robust
to non-normality of abnormal returns.
\end{Description}
%
\begin{Section}{Super class}
\code{\LinkA{TestStatisticBase}{TestStatisticBase}} -> \code{RankTest}
\end{Section}
%
\begin{Section}{Public fields}

\begin{description}

\item[\code{name}] Short code of the test statistic.

\end{description}


\end{Section}
%
\begin{Section}{Methods}
%
\begin{SubSection}{Public methods}
\begin{itemize}

\item{} \Rhref{\#method-RankTest-compute}{\code{RankTest\$compute()}}
\item{} \Rhref{\#method-RankTest-clone}{\code{RankTest\$clone()}}

\end{itemize}

\end{SubSection}




\hypertarget{method-RankTest-compute}{}
%
\begin{SubSection}{\code{RankTest\$compute()}}
Computes the Corrado rank test for multiple events.
%
\begin{SubSubSection}{Usage}

\begin{alltt}RankTest$compute(data_tbl, model)\end{alltt}



\end{SubSubSection}

%
\begin{SubSubSection}{Arguments}

\begin{description}

\item[\code{data\_tbl}] The data for a multiple event with
calculated abnormal returns.
\item[\code{model}] The fitted model (unused).

\end{description}



\end{SubSubSection}

\end{SubSection}




\hypertarget{method-RankTest-clone}{}
%
\begin{SubSection}{\code{RankTest\$clone()}}
The objects of this class are cloneable with this method.
%
\begin{SubSubSection}{Usage}

\begin{alltt}RankTest$clone(deep = FALSE)\end{alltt}



\end{SubSubSection}

%
\begin{SubSubSection}{Arguments}

\begin{description}

\item[\code{deep}] Whether to make a deep clone.

\end{description}



\end{SubSubSection}

\end{SubSection}


\end{Section}
\HeaderA{recommend\_stat}{Recommend Test Statistics via Offline KB Matching}{recommend.Rul.stat}
\methaliasA{recommend\_stat.default}{recommend\_stat}{recommend.Rul.stat.default}
\methaliasA{recommend\_stat.es\_diagnostics}{recommend\_stat}{recommend.Rul.stat.es.Rul.diagnostics}
\methaliasA{recommend\_stat.EventStudyTask}{recommend\_stat}{recommend.Rul.stat.EventStudyTask}
%
\begin{Description}
Evaluates the KB decision table (category \code{"stat\_choice"} rules) against
the diagnostic signals extracted from a fitted \code{EventStudyTask} or a
precomputed \code{es\_diagnostics} object, and returns a severity-ranked
\code{es\_advice} S3 object.
\end{Description}
%
\begin{Usage}
\begin{verbatim}
recommend_stat(x, provider = NULL, ...)

## Default S3 method:
recommend_stat(x, provider = NULL, ...)

## S3 method for class 'EventStudyTask'
recommend_stat(x, provider = NULL, ...)

## S3 method for class 'es_diagnostics'
recommend_stat(x, provider = NULL, ...)
\end{verbatim}
\end{Usage}
%
\begin{Arguments}
\begin{ldescription}
\item[\code{x}] A fitted \code{EventStudyTask} (after \code{fit\_model()}) or a
precomputed \code{es\_diagnostics} object returned by \code{es\_diagnostics()}.

\item[\code{provider}] Accepted but ignored in the offline path — present only so
the Phase 7 call shape is forward-compatible. Default \code{NULL}.

\item[\code{...}] Additional arguments (currently ignored).
\end{ldescription}
\end{Arguments}
%
\begin{Details}
No LLM provider, network connection, or API key is required. Both functions
are the always-available offline grounding layer (ADV-08). The returned
\code{es\_advice} object has the same shape as the Phase 7 Advice contract,
flagged \code{is\_deterministic = TRUE} and \code{source = "offline\_kb"}.
\end{Details}
%
\begin{Value}
An S3 object of class \code{"es\_advice"} — a named list with:
\begin{description}

\item[\code{source}] \code{"offline\_kb"} (character).
\item[\code{is\_deterministic}] \code{TRUE} — advice is rule-based, not LLM-generated.
\item[\code{rules\_matched}] Named list of matched rule records (severity-ranked:
\code{"error"} first, then \code{"warning"}, then \code{"info"}), each with
fields \code{id}, \code{recommendation}, \code{citation} (list of
\code{author}/\code{year}/\code{key}/\code{venue}), \code{severity},
\code{category}.
\item[\code{diagnostics\_ref}] The \code{es\_diagnostics} list that was evaluated
(possibly computed on-the-fly from the task).

\end{description}

\end{Value}
%
\begin{SeeAlso}
\code{\LinkA{flag\_robustness}{flag.Rul.robustness}}, \code{\LinkA{es\_diagnostics}{es.Rul.diagnostics}},
\code{\LinkA{es\_kb}{es.Rul.kb}}
\end{SeeAlso}
%
\begin{Examples}
\begin{ExampleCode}
## Not run: 
task <- run_event_study(my_task, ParameterSet$new())
advice <- recommend_stat(task)
print(advice)

## End(Not run)

\end{ExampleCode}
\end{Examples}
\HeaderA{ReturnCalculation}{R6 base class for return calculation}{ReturnCalculation}
%
\begin{Description}
R6 base class for return calculation
\end{Description}
%
\begin{Section}{Public fields}

\begin{description}

\item[\code{name}] Name of the return calculation.

\end{description}


\end{Section}
%
\begin{Section}{Methods}
%
\begin{SubSection}{Public methods}
\begin{itemize}

\item{} \Rhref{\#method-ReturnCalculation-calculate_return}{\code{ReturnCalculation\$calculate\_return()}}
\item{} \Rhref{\#method-ReturnCalculation-clone}{\code{ReturnCalculation\$clone()}}

\end{itemize}

\end{SubSection}



\hypertarget{method-ReturnCalculation-calculate_return}{}
%
\begin{SubSection}{\code{ReturnCalculation\$calculate\_return()}}
Calculates the return for a
single stock.
%
\begin{SubSubSection}{Usage}

\begin{alltt}ReturnCalculation$calculate_return(
  tbl,
  in_column = "adjusted",
  out_column = "adjusted_return"
)\end{alltt}



\end{SubSubSection}

%
\begin{SubSubSection}{Arguments}

\begin{description}

\item[\code{tbl}] The dataframe with the stock price.
\item[\code{in\_column}] The column name of the price
infromation.
\item[\code{out\_column}] The column name were the
return will be saved.

\end{description}



\end{SubSubSection}

\end{SubSection}




\hypertarget{method-ReturnCalculation-clone}{}
%
\begin{SubSection}{\code{ReturnCalculation\$clone()}}
The objects of this class are cloneable with this method.
%
\begin{SubSubSection}{Usage}

\begin{alltt}ReturnCalculation$clone(deep = FALSE)\end{alltt}



\end{SubSubSection}

%
\begin{SubSubSection}{Arguments}

\begin{description}

\item[\code{deep}] Whether to make a deep clone.

\end{description}



\end{SubSubSection}

\end{SubSection}


\end{Section}
\HeaderA{RollingWindowModel}{Rolling Window Model}{RollingWindowModel}
%
\begin{Description}
Event study model with time-varying parameters estimated via a rolling
OLS window over the estimation period. The last rolling window's
parameters are used for event-window prediction. This captures parameter
instability that is common in financial return data.
\end{Description}
%
\begin{Section}{Super class}
\code{\LinkA{ModelBase}{ModelBase}} -> \code{RollingWindowModel}
\end{Section}
%
\begin{Section}{Public fields}

\begin{description}

\item[\code{model\_name}] Name of the model.

\item[\code{window\_size}] Rolling window size. Default 60.

\item[\code{min\_obs}] Minimum observations required. Default 30.

\end{description}


\end{Section}
%
\begin{Section}{Methods}
%
\begin{SubSection}{Public methods}
\begin{itemize}

\item{} \Rhref{\#method-RollingWindowModel-initialize}{\code{RollingWindowModel\$new()}}
\item{} \Rhref{\#method-RollingWindowModel-fit}{\code{RollingWindowModel\$fit()}}
\item{} \Rhref{\#method-RollingWindowModel-abnormal_returns}{\code{RollingWindowModel\$abnormal\_returns()}}
\item{} \Rhref{\#method-RollingWindowModel-clone}{\code{RollingWindowModel\$clone()}}

\end{itemize}

\end{SubSection}



\hypertarget{method-RollingWindowModel-initialize}{}
%
\begin{SubSection}{\code{RollingWindowModel\$new()}}
Create a new RollingWindowModel.
%
\begin{SubSubSection}{Usage}

\begin{alltt}RollingWindowModel$new(window_size = 60L, min_obs = 30L)\end{alltt}



\end{SubSubSection}

%
\begin{SubSubSection}{Arguments}

\begin{description}

\item[\code{window\_size}] Size of the rolling window.
\item[\code{min\_obs}] Minimum observations for a valid window.

\end{description}



\end{SubSubSection}

\end{SubSection}




\hypertarget{method-RollingWindowModel-fit}{}
%
\begin{SubSection}{\code{RollingWindowModel\$fit()}}
Fit the rolling window model on the estimation window.
%
\begin{SubSubSection}{Usage}

\begin{alltt}RollingWindowModel$fit(data_tbl)\end{alltt}



\end{SubSubSection}

%
\begin{SubSubSection}{Arguments}

\begin{description}

\item[\code{data\_tbl}] Data frame or tibble with firm\_returns,
index\_returns, estimation\_window, event\_window columns.

\end{description}



\end{SubSubSection}

\end{SubSection}




\hypertarget{method-RollingWindowModel-abnormal_returns}{}
%
\begin{SubSection}{\code{RollingWindowModel\$abnormal\_returns()}}
Calculate abnormal returns using the last rolling window parameters.
%
\begin{SubSubSection}{Usage}

\begin{alltt}RollingWindowModel$abnormal_returns(data_tbl)\end{alltt}



\end{SubSubSection}

%
\begin{SubSubSection}{Arguments}

\begin{description}

\item[\code{data\_tbl}] Data frame or tibble.

\end{description}



\end{SubSubSection}

\end{SubSection}




\hypertarget{method-RollingWindowModel-clone}{}
%
\begin{SubSection}{\code{RollingWindowModel\$clone()}}
The objects of this class are cloneable with this method.
%
\begin{SubSubSection}{Usage}

\begin{alltt}RollingWindowModel$clone(deep = FALSE)\end{alltt}



\end{SubSubSection}

%
\begin{SubSubSection}{Arguments}

\begin{description}

\item[\code{deep}] Whether to make a deep clone.

\end{description}



\end{SubSubSection}

\end{SubSection}


\end{Section}
\HeaderA{run\_event\_study}{Run a Complete Event Study}{run.Rul.event.Rul.study}
%
\begin{Description}
Convenience wrapper that runs the full event study pipeline: prepare data,
fit models, and calculate test statistics in a single call.
\end{Description}
%
\begin{Usage}
\begin{verbatim}
run_event_study(task, parameter_set = ParameterSet$new())
\end{verbatim}
\end{Usage}
%
\begin{Arguments}
\begin{ldescription}
\item[\code{task}] An EventStudyTask object.

\item[\code{parameter\_set}] A ParameterSet object defining the event study.
Defaults to a new ParameterSet with default settings.
\end{ldescription}
\end{Arguments}
%
\begin{Value}
The task object with all results computed.
\end{Value}
\HeaderA{sc\_placebo\_test}{Placebo Test for Synthetic Control}{sc.Rul.placebo.Rul.test}
%
\begin{Description}
Performs a placebo (permutation) test by re-estimating the synthetic
control for each donor unit as a pseudo-treated unit. The p-value is
the rank of the treated unit's RMSPE ratio among all units.
\end{Description}
%
\begin{Usage}
\begin{verbatim}
sc_placebo_test(task, n_placebo = NULL)
\end{verbatim}
\end{Usage}
%
\begin{Arguments}
\begin{ldescription}
\item[\code{task}] A \code{SyntheticControlTask} with results.

\item[\code{n\_placebo}] Number of placebo units to use. Default NULL (all donors).
\end{ldescription}
\end{Arguments}
%
\begin{Value}
The task with \code{results\$placebo} populated, a list containing
\code{rmspe\_ratios}, \code{p\_value}, and \code{placebo\_gaps}.
\end{Value}
\HeaderA{SignTest}{Sign Test for Multiple Events}{SignTest}
%
\begin{Description}
Tests whether the proportion of positive abnormal returns differs
significantly from 0.5 under the null hypothesis. The test statistic
is approximately distributed as N(0, 1).
\end{Description}
%
\begin{Section}{Super class}
\code{\LinkA{TestStatisticBase}{TestStatisticBase}} -> \code{SignTest}
\end{Section}
%
\begin{Section}{Public fields}

\begin{description}

\item[\code{name}] Short code of the test statistic.

\end{description}


\end{Section}
%
\begin{Section}{Methods}
%
\begin{SubSection}{Public methods}
\begin{itemize}

\item{} \Rhref{\#method-SignTest-compute}{\code{SignTest\$compute()}}
\item{} \Rhref{\#method-SignTest-clone}{\code{SignTest\$clone()}}

\end{itemize}

\end{SubSection}




\hypertarget{method-SignTest-compute}{}
%
\begin{SubSection}{\code{SignTest\$compute()}}
Computes the sign test for multiple events.
%
\begin{SubSubSection}{Usage}

\begin{alltt}SignTest$compute(data_tbl, model)\end{alltt}



\end{SubSubSection}

%
\begin{SubSubSection}{Arguments}

\begin{description}

\item[\code{data\_tbl}] The data for a multiple event with
calculated abnormal returns.
\item[\code{model}] The fitted model (unused).

\end{description}



\end{SubSubSection}

\end{SubSection}




\hypertarget{method-SignTest-clone}{}
%
\begin{SubSection}{\code{SignTest\$clone()}}
The objects of this class are cloneable with this method.
%
\begin{SubSubSection}{Usage}

\begin{alltt}SignTest$clone(deep = FALSE)\end{alltt}



\end{SubSubSection}

%
\begin{SubSubSection}{Arguments}

\begin{description}

\item[\code{deep}] Whether to make a deep clone.

\end{description}



\end{SubSubSection}

\end{SubSection}


\end{Section}
\HeaderA{SimpleReturn}{R6 class for simple return calculation}{SimpleReturn}
%
\begin{Description}
R6 class for simple return calculation
\end{Description}
%
\begin{Section}{Super class}
\code{\LinkA{ReturnCalculation}{ReturnCalculation}} -> \code{SimpleReturn}
\end{Section}
%
\begin{Section}{Public fields}

\begin{description}

\item[\code{name}] Name of the return calculation.

\end{description}


\end{Section}
%
\begin{Section}{Methods}
%
\begin{SubSection}{Public methods}
\begin{itemize}

\item{} \Rhref{\#method-SimpleReturn-calculate_return}{\code{SimpleReturn\$calculate\_return()}}
\item{} \Rhref{\#method-SimpleReturn-clone}{\code{SimpleReturn\$clone()}}

\end{itemize}

\end{SubSection}



\hypertarget{method-SimpleReturn-calculate_return}{}
%
\begin{SubSection}{\code{SimpleReturn\$calculate\_return()}}
Calculates the simple return for a
single stock.
%
\begin{SubSubSection}{Usage}

\begin{alltt}SimpleReturn$calculate_return(
  tbl,
  in_column = "adjusted",
  out_column = "adjusted_return"
)\end{alltt}



\end{SubSubSection}

%
\begin{SubSubSection}{Arguments}

\begin{description}

\item[\code{tbl}] The dataframe with the stock price.
\item[\code{in\_column}] The column name of the price
infromation.
\item[\code{out\_column}] The column name were the
return will be saved.

\end{description}



\end{SubSubSection}

\end{SubSection}




\hypertarget{method-SimpleReturn-clone}{}
%
\begin{SubSection}{\code{SimpleReturn\$clone()}}
The objects of this class are cloneable with this method.
%
\begin{SubSubSection}{Usage}

\begin{alltt}SimpleReturn$clone(deep = FALSE)\end{alltt}



\end{SubSubSection}

%
\begin{SubSubSection}{Arguments}

\begin{description}

\item[\code{deep}] Whether to make a deep clone.

\end{description}



\end{SubSubSection}

\end{SubSection}


\end{Section}
\HeaderA{simulate\_event\_study}{Simulate Event Study for Power Analysis}{simulate.Rul.event.Rul.study}
%
\begin{Description}
Runs Monte Carlo simulations to estimate the statistical power (rejection
rate) of an event study design. Each simulation generates synthetic data
from a user-specified DGP, runs the full event study pipeline, and records
whether the null hypothesis is rejected.
\end{Description}
%
\begin{Usage}
\begin{verbatim}
simulate_event_study(
  n_events = 20,
  event_window = c(-5, 5),
  estimation_window_length = 120,
  abnormal_return = 0,
  return_model = MarketModel$new(),
  test_statistic = "CSectT",
  alpha = 0.05,
  n_simulations = 1000,
  dgp_params = list(),
  seed = NULL
)
\end{verbatim}
\end{Usage}
%
\begin{Arguments}
\begin{ldescription}
\item[\code{n\_events}] Number of events (firms) per simulation. Default 20.

\item[\code{event\_window}] Event window as \code{c(start, end)}. Default \code{c(-5, 5)}.

\item[\code{estimation\_window\_length}] Length of the estimation window. Default 120.

\item[\code{abnormal\_return}] Injected abnormal return on the event day (day 0).
Set to 0 for size analysis, positive for power analysis. Default 0.

\item[\code{return\_model}] An initialized return model object. Default \code{MarketModel\$new()}.

\item[\code{test\_statistic}] Name of the multi-event test statistic to evaluate.
Default \code{"CSectT"}.

\item[\code{alpha}] Significance level. Default 0.05.

\item[\code{n\_simulations}] Number of Monte Carlo replications. Default 1000.

\item[\code{dgp\_params}] List of DGP parameters: \code{alpha} (drift), \code{beta}
(market exposure), \code{sigma\_firm} (idiosyncratic volatility),
\code{sigma\_market} (market volatility).

\item[\code{seed}] Optional seed for reproducibility.
\end{ldescription}
\end{Arguments}
%
\begin{Value}
An S3 object of class \code{es\_simulation} with components:
\begin{description}

\item[power] Rejection rate at the event day
\item[rejection\_by\_day] Tibble of rejection rates for each event-window day
\item[test\_stats] Vector of test statistics at the event day from each sim
\item[params] List of simulation parameters

\end{description}

\end{Value}
\HeaderA{SingleEventStatisticsSet}{Single Event Statistic Set}{SingleEventStatisticsSet}
%
\begin{Description}
Contains a set of single event test statistics as, e.g., AR and CAR t test.
These are the default test statistics.
\end{Description}
%
\begin{Section}{Super class}
\code{\LinkA{StatisticsSetBase}{StatisticsSetBase}} -> \code{SingleEventStatisticsSet}
\end{Section}
%
\begin{Section}{Public fields}

\begin{description}

\item[\code{tests}] Container with the statistical tests.

\end{description}


\end{Section}
%
\begin{Section}{Methods}
%
\begin{SubSection}{Public methods}
\begin{itemize}

\item{} \Rhref{\#method-SingleEventStatisticsSet-clone}{\code{SingleEventStatisticsSet\$clone()}}

\end{itemize}

\end{SubSection}




\hypertarget{method-SingleEventStatisticsSet-clone}{}
%
\begin{SubSection}{\code{SingleEventStatisticsSet\$clone()}}
The objects of this class are cloneable with this method.
%
\begin{SubSubSection}{Usage}

\begin{alltt}SingleEventStatisticsSet$clone(deep = FALSE)\end{alltt}



\end{SubSubSection}

%
\begin{SubSubSection}{Arguments}

\begin{description}

\item[\code{deep}] Whether to make a deep clone.

\end{description}



\end{SubSubSection}

\end{SubSection}


\end{Section}
\HeaderA{StatisticsSetBase}{Adds a test statistic to this container.}{StatisticsSetBase}
%
\begin{Description}
Base class for test statistic container
\end{Description}
%
\begin{Section}{Public fields}

\begin{description}

\item[\code{tests}] Container with the statistical tests.
Test container initialization

\end{description}


\end{Section}
%
\begin{Section}{Methods}
%
\begin{SubSection}{Public methods}
\begin{itemize}

\item{} \Rhref{\#method-StatisticsSetBase-initialize}{\code{StatisticsSetBase\$new()}}
\item{} \Rhref{\#method-StatisticsSetBase-add_test}{\code{StatisticsSetBase\$add\_test()}}
\item{} \Rhref{\#method-StatisticsSetBase-clone}{\code{StatisticsSetBase\$clone()}}

\end{itemize}

\end{SubSection}



\hypertarget{method-StatisticsSetBase-initialize}{}
%
\begin{SubSection}{\code{StatisticsSetBase\$new()}}
Initializes a test statistic
container. Single event and multiple event
test statistics are collected in separated
containers.
%
\begin{SubSubSection}{Usage}

\begin{alltt}StatisticsSetBase$new(tests = NULL)\end{alltt}



\end{SubSubSection}

%
\begin{SubSubSection}{Arguments}

\begin{description}

\item[\code{tests}] List of test statistics.

\end{description}



\end{SubSubSection}

\end{SubSection}




\hypertarget{method-StatisticsSetBase-add_test}{}
%
\begin{SubSection}{\code{StatisticsSetBase\$add\_test()}}
%
\begin{SubSubSection}{Usage}

\begin{alltt}StatisticsSetBase$add_test(test)\end{alltt}



\end{SubSubSection}

%
\begin{SubSubSection}{Arguments}

\begin{description}

\item[\code{test}] A single event study tests.

\end{description}



\end{SubSubSection}

\end{SubSection}




\hypertarget{method-StatisticsSetBase-clone}{}
%
\begin{SubSection}{\code{StatisticsSetBase\$clone()}}
The objects of this class are cloneable with this method.
%
\begin{SubSubSection}{Usage}

\begin{alltt}StatisticsSetBase$clone(deep = FALSE)\end{alltt}



\end{SubSubSection}

%
\begin{SubSubSection}{Arguments}

\begin{description}

\item[\code{deep}] Whether to make a deep clone.

\end{description}



\end{SubSubSection}

\end{SubSection}


\end{Section}
\HeaderA{SyntheticControlTask}{Synthetic Control Task}{SyntheticControlTask}
%
\begin{Description}
Task container for synthetic control analysis. Holds treated unit data,
donor pool data, treatment time, and results after estimation.
\end{Description}
%
\begin{Section}{Public fields}

\begin{description}

\item[\code{treated\_data}] Tibble with \code{time} and \code{outcome} columns
for the treated unit.

\item[\code{donor\_data}] Tibble in long format with \code{unit}, \code{time},
and \code{outcome} columns for donor units.

\item[\code{treatment\_time}] The time period when treatment begins.

\item[\code{results}] Estimation results (populated after estimation).

\end{description}


\end{Section}
%
\begin{Section}{Methods}
%
\begin{SubSection}{Public methods}
\begin{itemize}

\item{} \Rhref{\#method-SyntheticControlTask-initialize}{\code{SyntheticControlTask\$new()}}
\item{} \Rhref{\#method-SyntheticControlTask-print}{\code{SyntheticControlTask\$print()}}
\item{} \Rhref{\#method-SyntheticControlTask-clone}{\code{SyntheticControlTask\$clone()}}

\end{itemize}

\end{SubSection}



\hypertarget{method-SyntheticControlTask-initialize}{}
%
\begin{SubSection}{\code{SyntheticControlTask\$new()}}
Create a new SyntheticControlTask.
%
\begin{SubSubSection}{Usage}

\begin{alltt}SyntheticControlTask$new(treated_data, donor_data, treatment_time)\end{alltt}



\end{SubSubSection}

%
\begin{SubSubSection}{Arguments}

\begin{description}

\item[\code{treated\_data}] Tibble with \code{time} and \code{outcome}.
\item[\code{donor\_data}] Long-format tibble with \code{unit}, \code{time},
\code{outcome}.
\item[\code{treatment\_time}] Period when treatment begins.

\end{description}



\end{SubSubSection}

\end{SubSection}




\hypertarget{method-SyntheticControlTask-print}{}
%
\begin{SubSection}{\code{SyntheticControlTask\$print()}}
Print summary.
%
\begin{SubSubSection}{Usage}

\begin{alltt}SyntheticControlTask$print(...)\end{alltt}



\end{SubSubSection}

\end{SubSection}




\hypertarget{method-SyntheticControlTask-clone}{}
%
\begin{SubSection}{\code{SyntheticControlTask\$clone()}}
The objects of this class are cloneable with this method.
%
\begin{SubSubSection}{Usage}

\begin{alltt}SyntheticControlTask$clone(deep = FALSE)\end{alltt}



\end{SubSubSection}

%
\begin{SubSubSection}{Arguments}

\begin{description}

\item[\code{deep}] Whether to make a deep clone.

\end{description}



\end{SubSubSection}

\end{SubSection}


\end{Section}
\HeaderA{TestStatisticBase}{Base class for Event Study test statistics}{TestStatisticBase}
%
\begin{Description}
Base class for Event Study test statistics
\end{Description}
%
\begin{Section}{Public fields}

\begin{description}

\item[\code{name}] Short code of the test statistic.

\item[\code{confidence\_level}] The chosen confidence
level.

\item[\code{confidence\_type}] Type of the test. Defaults
to 'two-sided'. Alternatives are 'less' or
greater'.

\end{description}


\end{Section}
%
\begin{Section}{Methods}
%
\begin{SubSection}{Public methods}
\begin{itemize}

\item{} \Rhref{\#method-TestStatisticBase-initialize}{\code{TestStatisticBase\$new()}}
\item{} \Rhref{\#method-TestStatisticBase-compute}{\code{TestStatisticBase\$compute()}}
\item{} \Rhref{\#method-TestStatisticBase-clone}{\code{TestStatisticBase\$clone()}}

\end{itemize}

\end{SubSection}



\hypertarget{method-TestStatisticBase-initialize}{}
%
\begin{SubSection}{\code{TestStatisticBase\$new()}}
Initializes the test statistic. This includes the
confidence level and the type of the test ('less',
greater' or 'two-sided')
%
\begin{SubSubSection}{Usage}

\begin{alltt}TestStatisticBase$new(confidence_level = 0.95, confidence_type = "two-sided")\end{alltt}



\end{SubSubSection}

%
\begin{SubSubSection}{Arguments}

\begin{description}

\item[\code{confidence\_level}] The confidence level for
the confidence band. Must be anumber between 0
and 1.
\item[\code{confidence\_type}] Side of the test statistic.

\end{description}



\end{SubSubSection}

\end{SubSection}




\hypertarget{method-TestStatisticBase-compute}{}
%
\begin{SubSection}{\code{TestStatisticBase\$compute()}}
Computes the test test statistics for a single event.
%
\begin{SubSubSection}{Usage}

\begin{alltt}TestStatisticBase$compute(data_tbl, model)\end{alltt}



\end{SubSubSection}

%
\begin{SubSubSection}{Arguments}

\begin{description}

\item[\code{data\_tbl}] The data for a single event with
calculated abnormal returns.
\item[\code{model}] The fitted model that includes the
necessary information for calculating the test
statistic.

\end{description}



\end{SubSubSection}

\end{SubSection}




\hypertarget{method-TestStatisticBase-clone}{}
%
\begin{SubSection}{\code{TestStatisticBase\$clone()}}
The objects of this class are cloneable with this method.
%
\begin{SubSubSection}{Usage}

\begin{alltt}TestStatisticBase$clone(deep = FALSE)\end{alltt}



\end{SubSubSection}

%
\begin{SubSubSection}{Arguments}

\begin{description}

\item[\code{deep}] Whether to make a deep clone.

\end{description}



\end{SubSubSection}

\end{SubSection}


\end{Section}
\HeaderA{tidy.EventStudyTask}{Tidy Event Study Results}{tidy.EventStudyTask}
%
\begin{Description}
Extract event study results in a tidy (long-format) tibble compatible
with broom conventions. This enables integration with standard tidyverse
workflows.
\end{Description}
%
\begin{Usage}
\begin{verbatim}
tidy.EventStudyTask(
  x,
  type = c("ar", "car", "aar", "model"),
  stat_name = "CSectT",
  ...
)
\end{verbatim}
\end{Usage}
%
\begin{Arguments}
\begin{ldescription}
\item[\code{x}] An EventStudyTask object.

\item[\code{type}] Type of results to extract: "ar" (abnormal returns),
"car" (cumulative abnormal returns), "aar" (average abnormal returns
with test statistics), or "model" (model fit statistics).

\item[\code{stat\_name}] For type "aar", the name of the multi-event test statistic.
Defaults to "CSectT".

\item[\code{...}] Additional arguments (unused).
\end{ldescription}
\end{Arguments}
%
\begin{Value}
A tibble with columns following broom conventions:
\begin{description}

\item[term] Identifier for the observation (e.g., relative\_index)
\item[estimate] The estimated value (AR, CAR, AAR, or coefficient)
\item[std.error] Standard error where available
\item[statistic] Test statistic value
\item[p.value] p-value where available

\end{description}

\end{Value}
\HeaderA{validate\_task}{Validate an Event Study experiment}{validate.Rul.task}
%
\begin{Description}
Performs validation checks on the event study task:
1. Does the event date exist in the stock data for each event?
2. Are there sufficient observations in the estimation window?
3. Are there any gaps in the time series?
4. Do estimation and event windows overlap?
\end{Description}
%
\begin{Usage}
\begin{verbatim}
validate_task(task, parameter_set = NULL, min_estimation_obs = 30)
\end{verbatim}
\end{Usage}
%
\begin{Arguments}
\begin{ldescription}
\item[\code{task}] The event study task.

\item[\code{parameter\_set}] The parameter set that defines the event study.

\item[\code{min\_estimation\_obs}] Minimum number of observations required in the
estimation window. Default is 30.
\end{ldescription}
\end{Arguments}
%
\begin{Value}
The task object (invisibly). Warnings are issued for each problem found.
\end{Value}
\HeaderA{VolatilityModel}{Volatility Event Study Model}{VolatilityModel}
%
\begin{Description}
Model for volatility-based event studies. Computes abnormal volatility
as the ratio of event-window squared returns to estimation-window variance.
The abnormal measure is written to the \code{abnormal\_returns} column
for compatibility with existing test statistics.
\end{Description}
%
\begin{Section}{Super class}
\code{\LinkA{ModelBase}{ModelBase}} -> \code{VolatilityModel}
\end{Section}
%
\begin{Section}{Public fields}

\begin{description}

\item[\code{model\_name}] Name of the model.

\end{description}


\end{Section}
%
\begin{Section}{Methods}
%
\begin{SubSection}{Public methods}
\begin{itemize}

\item{} \Rhref{\#method-VolatilityModel-fit}{\code{VolatilityModel\$fit()}}
\item{} \Rhref{\#method-VolatilityModel-abnormal_returns}{\code{VolatilityModel\$abnormal\_returns()}}
\item{} \Rhref{\#method-VolatilityModel-clone}{\code{VolatilityModel\$clone()}}

\end{itemize}

\end{SubSection}



\hypertarget{method-VolatilityModel-fit}{}
%
\begin{SubSection}{\code{VolatilityModel\$fit()}}
Fit the volatility model. Estimates expected variance from
estimation window.
%
\begin{SubSubSection}{Usage}

\begin{alltt}VolatilityModel$fit(data_tbl)\end{alltt}



\end{SubSubSection}

%
\begin{SubSubSection}{Arguments}

\begin{description}

\item[\code{data\_tbl}] Data frame or tibble.

\end{description}



\end{SubSubSection}

\end{SubSection}




\hypertarget{method-VolatilityModel-abnormal_returns}{}
%
\begin{SubSection}{\code{VolatilityModel\$abnormal\_returns()}}
Calculate abnormal volatility (squared returns / expected variance - 1).
%
\begin{SubSubSection}{Usage}

\begin{alltt}VolatilityModel$abnormal_returns(data_tbl)\end{alltt}



\end{SubSubSection}

%
\begin{SubSubSection}{Arguments}

\begin{description}

\item[\code{data\_tbl}] Data frame or tibble.

\end{description}



\end{SubSubSection}

\end{SubSection}




\hypertarget{method-VolatilityModel-clone}{}
%
\begin{SubSection}{\code{VolatilityModel\$clone()}}
The objects of this class are cloneable with this method.
%
\begin{SubSubSection}{Usage}

\begin{alltt}VolatilityModel$clone(deep = FALSE)\end{alltt}



\end{SubSubSection}

%
\begin{SubSubSection}{Arguments}

\begin{description}

\item[\code{deep}] Whether to make a deep clone.

\end{description}



\end{SubSubSection}

\end{SubSection}


\end{Section}
\HeaderA{VolumeModel}{Volume Event Study Model}{VolumeModel}
%
\begin{Description}
Model for volume-based event studies. Computes abnormal volume as the
difference between observed volume and expected volume from the estimation
window mean. The data must contain a \code{firm\_volume} column (and
optionally \code{index\_volume} for market-adjusted volume).

The existing test statistics infrastructure works on the
\code{abnormal\_returns} column, so this model writes abnormal volume
to that same column for compatibility.
\end{Description}
%
\begin{Section}{Super class}
\code{\LinkA{ModelBase}{ModelBase}} -> \code{VolumeModel}
\end{Section}
%
\begin{Section}{Public fields}

\begin{description}

\item[\code{model\_name}] Name of the model.

\item[\code{log\_transform}] Whether to log-transform volume. Default TRUE.

\end{description}


\end{Section}
%
\begin{Section}{Methods}
%
\begin{SubSection}{Public methods}
\begin{itemize}

\item{} \Rhref{\#method-VolumeModel-initialize}{\code{VolumeModel\$new()}}
\item{} \Rhref{\#method-VolumeModel-fit}{\code{VolumeModel\$fit()}}
\item{} \Rhref{\#method-VolumeModel-abnormal_returns}{\code{VolumeModel\$abnormal\_returns()}}
\item{} \Rhref{\#method-VolumeModel-clone}{\code{VolumeModel\$clone()}}

\end{itemize}

\end{SubSection}



\hypertarget{method-VolumeModel-initialize}{}
%
\begin{SubSection}{\code{VolumeModel\$new()}}
Create a new VolumeModel.
%
\begin{SubSubSection}{Usage}

\begin{alltt}VolumeModel$new(log_transform = TRUE)\end{alltt}



\end{SubSubSection}

%
\begin{SubSubSection}{Arguments}

\begin{description}

\item[\code{log\_transform}] Whether to log-transform volume before analysis.

\end{description}



\end{SubSubSection}

\end{SubSection}




\hypertarget{method-VolumeModel-fit}{}
%
\begin{SubSection}{\code{VolumeModel\$fit()}}
Fit the volume model. Computes expected volume from estimation window.
%
\begin{SubSubSection}{Usage}

\begin{alltt}VolumeModel$fit(data_tbl)\end{alltt}



\end{SubSubSection}

%
\begin{SubSubSection}{Arguments}

\begin{description}

\item[\code{data\_tbl}] Data frame or tibble with firm\_volume column.

\end{description}



\end{SubSubSection}

\end{SubSection}




\hypertarget{method-VolumeModel-abnormal_returns}{}
%
\begin{SubSection}{\code{VolumeModel\$abnormal\_returns()}}
Calculate abnormal volume.
%
\begin{SubSubSection}{Usage}

\begin{alltt}VolumeModel$abnormal_returns(data_tbl)\end{alltt}



\end{SubSubSection}

%
\begin{SubSubSection}{Arguments}

\begin{description}

\item[\code{data\_tbl}] Data frame or tibble.

\end{description}



\end{SubSubSection}

\end{SubSection}




\hypertarget{method-VolumeModel-clone}{}
%
\begin{SubSection}{\code{VolumeModel\$clone()}}
The objects of this class are cloneable with this method.
%
\begin{SubSubSection}{Usage}

\begin{alltt}VolumeModel$clone(deep = FALSE)\end{alltt}



\end{SubSubSection}

%
\begin{SubSubSection}{Arguments}

\begin{description}

\item[\code{deep}] Whether to make a deep clone.

\end{description}



\end{SubSubSection}

\end{SubSection}


\end{Section}
\printindex{}
\end{document}
